% ------------------------------------------------------------------------------
% Mathematics Project Report Template
% by Jon Shiach 2019
% Manchester Metropolitan University
% ------------------------------------------------------------------------------
\documentclass[11pt]{report} 

\usepackage{projectreport}

% ------------------------------------------------------------------------------
% Enter your details here
% ------------------------------------------------------------------------------
\newcommand{\name}{Bùi Đình Bảo}
\newcommand{\course}{Chương trình Viettel Digital Talent 2026}
% \newcommand{\course}{BSc. (Hons) Financial Mathematics}
% \newcommand{\course}{MMath. (Hons) Mathematics}
% \newcommand{\course}{BSc. (Hons) Secondary Mathematics Education with QTS}
\newcommand{\projecttitle}{Nghiên cứu Feature Flag và xây dựng hệ thống quản lý phát hành tính năng an toàn}
\newcommand{\submissiondate}{Tháng 7/2026}


% ------------------------------------------------------------------------------
% Document
% ------------------------------------------------------------------------------
\begin{document}

\maketitle

% ------------------------------------------------------------------------------
% Top matter
% ------------------------------------------------------------------------------

\chapter*{Lời cảm ơn}

Trước hết, em xin gửi lời cảm ơn chân thành đến Ban Tổ chức chương trình Viettel Digital Talent 2026 đã tạo điều kiện để em được tham gia học tập, rèn luyện và phát triển trong một môi trường chuyên nghiệp.

Em xin cảm ơn các anh mentor đã trực tiếp giảng dạy và chia sẻ kiến thức trong các buổi training của chương trình. Mỗi buổi đào tạo với một chủ đề và góc nhìn khác nhau đã mang lại cho em nhiều kiến thức, kinh nghiệm thực tế và những định hướng có giá trị trong quá trình học tập và hoàn thiện bản thân.

Em cũng xin cảm ơn mentor phụ trách Mini Project cùng Ban Tổ chức đã cung cấp định hướng, yêu cầu và những góp ý cần thiết trong quá trình thực hiện. Sự hỗ trợ của các anh đã giúp em xác định rõ mục tiêu, chủ động hơn trong quá trình nghiên cứu và hoàn thành sản phẩm.

Bên cạnh đó, em xin cảm ơn các bạn học viên trong chương trình đã cùng trao đổi, chia sẻ và hỗ trợ nhau trong suốt quá trình tham gia Viettel Digital Talent 2026.

Mặc dù đã dành nhiều thời gian và nỗ lực để hoàn thành báo cáo và sản phẩm, do kinh nghiệm còn hạn chế, nội dung thực hiện khó tránh khỏi những thiếu sót. Em rất mong nhận được những ý kiến đóng góp từ Ban Tổ chức và các anh/chị mentor để có thể tiếp tục cải thiện và hoàn thiện bản thân.

Em xin chân thành cảm ơn!

\hfill

\hfill TP. Hồ Chí Minh, tháng 7 năm 2026

\hfill Học viên thực hiện

\hfill Bùi Đình Bảo


\chapter*{Lời mở đầu}

Báo cáo này là sản phẩm Mini Project trong chương trình Viettel Digital Talent 2026, tập trung vào việc nghiên cứu và xây dựng một nền tảng Feature Flag có thể trình diễn được. Nội dung chính của báo cáo được tổ chức theo mạch: giới thiệu bài toán, cơ sở lý thuyết và thiết kế, kết quả triển khai, kiểm thử, đánh giá, hạn chế và hướng phát triển.

Các phần sau sẽ trình bày chi tiết khái niệm Feature Flag, lý do cần tách deploy khỏi release, kiến trúc control plane/data plane, Rule Evaluation Engine, Dashboard, SDK và Demo Application. Phần mở đầu này chỉ đóng vai trò định hướng cấu trúc báo cáo để tránh lặp lại nội dung chuyên môn ở các chương chính.

\chapter*{Tóm tắt nội dung và đóng góp}

Báo cáo nghiên cứu và xây dựng một mini Feature Flag Platform cho ứng dụng web, gồm backend, dashboard, demo app, SDK, PostgreSQL và các cơ chế hỗ trợ như audit log, cache, thống kê, RBAC demo, group kill switch và Docker Compose.

Các đóng góp, kết quả kiểm thử, hạn chế và hướng phát triển được trình bày chi tiết trong các chương chính; phần này chỉ tóm lược phạm vi sản phẩm để người đọc nắm cấu trúc báo cáo.

\tableofcontents
\listoffigures
\listoftables
\lstlistoflistings		% comment out if not needed

% ------------------------------------------------------------------------------
% Main matter
% Chapters are added to the document using the \include{chapterx} command
% ------------------------------------------------------------------------------
\newpage
\setcounter{page}{0}
\pagenumbering{arabic}
% ------------------------------------------------------------------------------
% Chapter 1
% Delete this content and replace with your own
% ------------------------------------------------------------------------------
\chapter{Giới thiệu} % enter the name of the chapter here

\section{Đặt vấn đề}

Trong quá trình phát triển phần mềm, việc đưa một tính năng mới vào hệ thống luôn đi kèm với rủi ro. Mặc dù tính năng đã được kiểm thử trước khi triển khai, các vấn đề về logic nghiệp vụ, hiệu năng, khả năng tương thích hoặc hành vi của người dùng vẫn có thể chỉ xuất hiện trên môi trường thực tế. Nếu tính năng được phát hành đồng thời cho toàn bộ người dùng, phạm vi ảnh hưởng của sự cố có thể lớn, trong khi việc khắc phục thường đòi hỏi rollback phiên bản hoặc thực hiện một lần triển khai mới.

Bên cạnh đó, nhu cầu phát hành phần mềm hiện nay không còn giới hạn ở hai trạng thái bật hoặc tắt cho toàn hệ thống. Đội phát triển có thể cần cung cấp tính năng cho một số người dùng cụ thể, một nhóm có vai trò nhất định hoặc một tỷ lệ người dùng được lựa chọn ổn định. Do đó, bài toán đặt ra là xây dựng một cơ chế cho phép kiểm soát việc phát hành tính năng độc lập với quá trình triển khai mã nguồn, đồng thời bảo đảm quyết định bật hoặc tắt tính năng được thực hiện nhất quán, có thể truy vết và dễ dàng quản lý.

Feature Flag là một giải pháp phù hợp cho bài toán này. Tuy nhiên, khi số lượng flag và quy tắc tăng lên, việc sử dụng các biến cấu hình rời rạc hoặc cài đặt logic trực tiếp trong từng ứng dụng có thể dẫn đến cấu hình không nhất quán, khó kiểm soát lịch sử thay đổi và khó mở rộng. Vì vậy, đề tài lựa chọn nghiên cứu và xây dựng một nền tảng quản lý Feature Flag tập trung cho ứng dụng web.

\section{Mục tiêu của báo cáo}

Báo cáo hướng đến hai nhóm mục tiêu chính: mục tiêu nghiên cứu và mục tiêu triển khai.

Về nghiên cứu, báo cáo làm rõ bản chất của Feature Flag, sự khác biệt giữa deploy và release, các loại flag phổ biến và vai trò của Feature Flag trong quản lý phát hành phần mềm. Báo cáo đồng thời phân tích các chiến lược đánh giá như bật toàn cục, nhắm mục tiêu theo người dùng, theo vai trò và rollout theo tỷ lệ phần trăm.

Về triển khai, đề tài xây dựng một hệ thống có khả năng quản lý project, feature flag và các rule tương ứng; tiếp nhận evaluation context từ ứng dụng tích hợp; đánh giá trạng thái flag theo một thứ tự ưu tiên xác định; trả về kết quả có lý do rõ ràng; và ghi nhận các thay đổi cấu hình thông qua audit log. Hệ thống cũng cung cấp dashboard quản trị và demo app để chứng minh khả năng ứng dụng trong thực tế. Các thành phần này phù hợp với yêu cầu về backend API, frontend dashboard, evaluation engine, audit log và ứng dụng minh họa của mini project.

Ngoài việc bảo đảm hệ thống có thể hoạt động, báo cáo còn tập trung giải thích các quyết định thiết kế và công nghệ, bao gồm cách tổ chức kiến trúc, mô hình dữ liệu, thứ tự đánh giá rule, cơ chế rollout ổn định, xử lý lỗi và khả năng mở rộng. Đây cũng là cơ sở để thể hiện khả năng giải quyết vấn đề, tư duy thiết kế và mức độ hiểu những nội dung đã triển khai.

\section{Câu hỏi nghiên cứu}

Báo cáo tập trung giải quyết các câu hỏi chính sau:

\begin{enumerate}
    \item Feature Flag giải quyết những hạn chế nào của quy trình triển khai và phát hành phần mềm truyền thống?
    \item Làm thế nào để tách quá trình deploy mã nguồn khỏi quyết định release tính năng cho người dùng?
    \item Một nền tảng Feature Flag cần những thành phần nào để hỗ trợ quản lý cấu hình, đánh giá rule, tích hợp với ứng dụng và truy vết thay đổi?
    \item Làm thế nào để xác định thứ tự ưu tiên giữa các trạng thái và rule nhằm bảo đảm kết quả đánh giá rõ ràng, nhất quán và có thể dự đoán?
    \item Làm thế nào để rollout tính năng theo tỷ lệ phần trăm mà vẫn giữ kết quả ổn định đối với cùng một người dùng?
    \item Hệ thống nên xử lý như thế nào khi feature flag không tồn tại, dữ liệu context bị thiếu hoặc xảy ra lỗi trong quá trình đánh giá?
    \item Kiến trúc và công nghệ nào phù hợp với quy mô của mini project, đồng thời vẫn tạo điều kiện để hệ thống được mở rộng trong tương lai?
\end{enumerate}

\section{Phạm vi triển khai}

Trong phạm vi báo cáo, đề tài triển khai một nền tảng Feature Flag phục vụ các ứng dụng web, gồm bốn thành phần chính:

\begin{itemize}
    \item Backend API chịu trách nhiệm quản lý dữ liệu, xử lý nghiệp vụ và cung cấp evaluation endpoint.
    \item Frontend dashboard hỗ trợ quản trị project, flag, rule, group, audit log và thống kê evaluation.
    \item Evaluation engine thực hiện việc đánh giá flag dựa trên trạng thái cấu hình và context của request.
    \item Demo app tích hợp evaluation API để minh họa việc hiển thị hoặc ẩn tính năng theo kết quả đánh giá.
\end{itemize}

Các chiến lược được hỗ trợ bao gồm bật hoặc tắt toàn cục, nhắm mục tiêu theo user ID, theo role và rollout theo phần trăm người dùng. Hệ thống lưu trữ project, môi trường, feature flag, cấu hình, rule, dữ liệu người dùng mẫu và audit log trong cơ sở dữ liệu. Ngoài ra, hệ thống thực hiện validation dữ liệu đầu vào, xử lý lỗi cơ bản, cung cấp dữ liệu mẫu và tài liệu hướng dẫn cài đặt.

Đề tài tập trung vào việc chứng minh đầy đủ luồng hoạt động của Feature Flag, từ khâu tạo và cấu hình flag đến khi ứng dụng gửi context và nhận kết quả đánh giá. Các vấn đề ở quy mô hệ thống phân tán lớn như đồng bộ cấu hình đa vùng, streaming update theo thời gian thực, độ sẵn sàng đa trung tâm dữ liệu, phân tích thử nghiệm thống kê chuyên sâu và quản lý hàng triệu request mỗi giây không thuộc phạm vi triển khai chính.

Các nội dung nâng cao như caching, SDK JavaScript, phân quyền, lịch sử cấu hình dựa trên audit log, thống kê lượt đánh giá, kill switch theo nhóm và Docker Compose được xem là phần mở rộng nhằm đánh giá khả năng phát triển tiếp của hệ thống, thay vì điều kiện bắt buộc để hoàn thành luồng chức năng cốt lõi.


% ------------------------------------------------------------------------------
% Chapter 2
% ------------------------------------------------------------------------------
\chapter{Nội dung} % enter the name of the chapter here
\label{cha:noi-dung} % enter the chapter label here (for cross referencing)

\section{Tổng quan về Feature Flag}

\subsection{Khái niệm Feature Flag}

\subsubsection{Định nghĩa Feature Flag}

Feature Flag, còn được gọi là Feature Toggle hoặc Feature Switch, là một kỹ thuật phát triển phần mềm cho phép hệ thống quyết định tại thời điểm chạy liệu một tính năng hoặc một nhánh xử lý có được kích hoạt hay không. Quyết định này được thực hiện dựa trên cấu hình bên ngoài mã nguồn và ngữ cảnh của đối tượng đang sử dụng hệ thống, thay vì được cố định trong quá trình biên dịch hoặc triển khai ứng dụng \cite{fowler_feature_toggles,unleash_feature_flags}.

Có thể xem Feature Flag như một công tắc logic nằm giữa mã nguồn đã được triển khai và người dùng cuối. Mã của tính năng mới có thể đã tồn tại trong môi trường chạy, nhưng tính năng chỉ được cung cấp cho người dùng khi flag tương ứng được đánh giá là bật. Nhờ đó, đội phát triển có thể thay đổi trạng thái phát hành của tính năng mà không cần chỉnh sửa mã nguồn hoặc triển khai lại toàn bộ ứng dụng.

Đặc điểm cốt lõi của Feature Flag là tách biệt hai hoạt động thường bị xem là đồng nhất: \textit{deploy} và \textit{release}. Deploy là quá trình đưa mã nguồn lên môi trường chạy, trong khi release là quá trình cho phép một nhóm người dùng thực sự truy cập và sử dụng tính năng. Sự tách biệt này giúp đội phát triển triển khai mã thường xuyên hơn nhưng vẫn kiểm soát được thời điểm, phạm vi và điều kiện phát hành tính năng \cite{microsoft_devops}.

\subsubsection{Các thành phần cơ bản của Feature Flag}

Một Feature Flag thường bao gồm các thành phần cơ bản sau:

\begin{itemize}
    \item \textbf{Khóa định danh của flag:} là giá trị duy nhất dùng để tham chiếu đến flag trong mã nguồn và trong các yêu cầu đánh giá, chẳng hạn như \texttt{new-checkout-page} hoặc \texttt{recommendation-engine}.

    \item \textbf{Trạng thái và cấu hình mặc định:} xác định flag đang được bật, tắt hoặc ngừng sử dụng. Giá trị mặc định được sử dụng khi không có quy tắc nào phù hợp hoặc khi hệ thống không thể đánh giá đầy đủ ngữ cảnh.

    \item \textbf{Quy tắc nhắm mục tiêu:} mô tả các điều kiện để xác định đối tượng được sử dụng tính năng. Quy tắc có thể dựa trên định danh người dùng, vai trò, tổ chức, thuộc tính ngữ cảnh hoặc các tiêu chí nghiệp vụ khác \cite{launchdarkly_targeting_rules}.

    \item \textbf{Chiến lược phân phối:} xác định cách tính năng được phát hành, chẳng hạn như bật toàn cục, bật cho một tập người dùng cụ thể hoặc rollout theo tỷ lệ phần trăm \cite{launchdarkly_percentage_rollouts}.

    \item \textbf{Ngữ cảnh đánh giá:} là tập hợp thông tin mô tả đối tượng đang yêu cầu sử dụng tính năng, ví dụ như \texttt{targetingKey}, mã người dùng, vai trò và các thuộc tính bổ sung. Ngữ cảnh này được truyền vào bộ máy đánh giá để đối chiếu với các quy tắc đã cấu hình \cite{openfeature_evaluation_context}.

    \item \textbf{Giá trị trả về:} là kết quả cuối cùng sau quá trình đánh giá, thường bao gồm trạng thái bật hoặc tắt, giá trị biến thể và lý do dẫn đến quyết định.
\end{itemize}

Trong một nền tảng quản lý Feature Flag, các thành phần trên thường được tổ chức theo project và environment. Cùng một flag có thể tồn tại trên nhiều môi trường như development, staging và production, nhưng có cấu hình và quy tắc khác nhau. Cách tổ chức này cho phép đội phát triển kiểm thử tính năng ở môi trường nội bộ trước khi mở rộng phạm vi phát hành trên môi trường production \cite{launchdarkly_feature_flags}.

\subsubsection{Cơ chế quyết định hành vi tại thời điểm chạy}

Feature Flag điều khiển hành vi của phần mềm thông qua quá trình đánh giá tại thời điểm chạy. Khi ứng dụng cần thực hiện một chức năng được kiểm soát bởi flag, ứng dụng gửi khóa của flag cùng với ngữ cảnh người dùng đến bộ máy đánh giá. Bộ máy này tải cấu hình hiện tại của flag, kiểm tra trạng thái tổng quát và lần lượt đối chiếu ngữ cảnh với các quy tắc đã được thiết lập.

Quá trình đánh giá có thể được biểu diễn khái quát như sau:

\begin{enumerate}
    \item Xác định flag theo project, environment và khóa định danh.
    \item Kiểm tra flag có tồn tại và có đang ở trạng thái cho phép phục vụ hay không.
    \item Kiểm tra các cơ chế vô hiệu hóa ưu tiên cao như trạng thái tắt hoặc kill switch.
    \item Nếu flag được bật toàn cục, trả về kết quả bật mà không cần đánh giá thêm.
    \item Nếu flag sử dụng chế độ nhắm mục tiêu, lần lượt đánh giá các quy tắc theo thứ tự ưu tiên.
    \item Nếu một quy tắc phù hợp, trả về giá trị được cấu hình cho quy tắc đó.
    \item Nếu không có quy tắc nào phù hợp, trả về giá trị mặc định của flag.
\end{enumerate}

Ví dụ, một tính năng có thể được cấu hình bật cho người dùng có vai trò \texttt{beta-customer}. Khi người dùng gửi yêu cầu, ứng dụng cung cấp vai trò này trong ngữ cảnh đánh giá. Nếu quy tắc vai trò được thỏa mãn, hệ thống trả về trạng thái bật và ứng dụng thực thi nhánh xử lý mới. Ngược lại, ứng dụng tiếp tục sử dụng hành vi cũ.

\begin{lstlisting}[
    caption={Điều khiển hành vi ứng dụng bằng Feature Flag},
    label={lst:feature-flag-runtime}
]
const result = await featureFlagClient.evaluate(
    "new-checkout-page",
    userContext
);

if (result.enabled) {
    return showNewCheckout();
}

return showCurrentCheckout();
\end{lstlisting}

Đối với rollout theo phần trăm, hệ thống thường sử dụng hàm băm ổn định trên khóa định danh của đối tượng và khóa của flag để gán đối tượng vào một nhóm cố định. Cơ chế này giúp cùng một người dùng nhận được kết quả nhất quán giữa nhiều lần truy cập, đồng thời cho phép đội phát triển tăng hoặc giảm tỷ lệ rollout mà không cần lưu trạng thái riêng cho từng người dùng \cite{launchdarkly_percentage_rollouts}.

Như vậy, Feature Flag không thay đổi mã nguồn tại thời điểm chạy mà thay đổi nhánh xử lý được lựa chọn trong mã nguồn đã triển khai. Hiệu quả của kỹ thuật này phụ thuộc vào việc thiết kế quy tắc rõ ràng, xác định thứ tự ưu tiên nhất quán và cung cấp giá trị mặc định an toàn khi quá trình đánh giá xảy ra lỗi.

\subsection{Deploy và Release}

\subsubsection{Khái niệm Deploy}

Deploy là quá trình đưa một phiên bản phần mềm từ môi trường phát triển hoặc kiểm thử lên một môi trường chạy cụ thể, chẳng hạn như staging hoặc production. Quá trình này có thể bao gồm biên dịch mã nguồn, đóng gói ứng dụng, cập nhật cơ sở dữ liệu, cấu hình hạ tầng, khởi động dịch vụ và kiểm tra trạng thái hoạt động sau triển khai.

Mục tiêu của deploy là làm cho mã nguồn mới có mặt và có khả năng thực thi trên môi trường đích. Tuy nhiên, việc mã nguồn đã được deploy không đồng nghĩa với việc tất cả chức năng chứa trong phiên bản đó phải được cung cấp ngay cho người dùng cuối. Một tính năng có thể đã tồn tại trong hệ thống production nhưng vẫn được ẩn, vô hiệu hóa hoặc chỉ cho phép một nhóm người dùng nhất định truy cập.

Trong quy trình phát triển hiện đại, deploy thường được tự động hóa thông qua các pipeline tích hợp và phân phối liên tục. Việc triển khai thường xuyên với các thay đổi nhỏ giúp giảm độ phức tạp của mỗi lần cập nhật, rút ngắn vòng phản hồi và hạn chế rủi ro so với việc tích lũy nhiều thay đổi trong một lần triển khai lớn \cite{microsoft_devops}.

\subsubsection{Khái niệm Release}

Release là quá trình đưa một tính năng hoặc một thay đổi trong phần mềm đến người sử dụng thực tế. Hoạt động này xác định thời điểm, phạm vi và đối tượng được phép truy cập chức năng, thay vì chỉ xác định mã nguồn đang được cài đặt ở môi trường nào.

Một release có thể được thực hiện cho toàn bộ người dùng hoặc giới hạn theo một nhóm cụ thể, chẳng hạn như người dùng nội bộ, khách hàng thử nghiệm, người dùng thuộc một vai trò nhất định hoặc một tỷ lệ phần trăm người dùng. Vì vậy, release mang tính chất nghiệp vụ và vận hành nhiều hơn deploy. Nó liên quan trực tiếp đến trải nghiệm người dùng, mức độ rủi ro được chấp nhận và chiến lược cung cấp tính năng của tổ chức.

Release cũng không nhất thiết yêu cầu một lần deploy mới. Khi mã nguồn của tính năng đã tồn tại trên production và được kiểm soát bằng Feature Flag, đội phát triển có thể release tính năng bằng cách thay đổi cấu hình của flag. Tương tự, tính năng có thể được tạm dừng hoặc thu hẹp phạm vi cung cấp mà không cần loại bỏ mã nguồn khỏi môi trường chạy.

\subsubsection{Sự khác biệt giữa Deploy và Release}

Deploy và release thường diễn ra gần nhau trong các quy trình phát triển truyền thống nên dễ bị xem là cùng một hoạt động. Tuy nhiên, hai khái niệm này giải quyết hai vấn đề khác nhau. Deploy trả lời câu hỏi mã nguồn đã được đưa lên môi trường chạy hay chưa, trong khi release trả lời câu hỏi người dùng đã được phép sử dụng tính năng hay chưa.

Sự khác biệt giữa deploy và release được thể hiện trong Bảng~\ref{tab:deploy-release-comparison}.

\begin{table}[H]
    \centering
    \caption{So sánh deploy và release}
    \label{tab:deploy-release-comparison}
    \begin{tabular}{p{0.22\textwidth} p{0.34\textwidth} p{0.34\textwidth}}
        \toprule
        \textbf{Tiêu chí} & \textbf{Deploy} & \textbf{Release} \\
        \midrule
        Mục tiêu
        & Đưa mã nguồn và thành phần phần mềm lên môi trường chạy
        & Cung cấp tính năng cho người dùng thực tế \\

        Đối tượng tác động
        & Hạ tầng, dịch vụ, cơ sở dữ liệu và phiên bản ứng dụng
        & Người dùng, nhóm người dùng hoặc phân khúc khách hàng \\

        Cách thực hiện
        & Pipeline CI/CD, script triển khai hoặc công cụ quản lý hạ tầng
        & Thay đổi cấu hình, bật flag hoặc điều chỉnh quy tắc nhắm mục tiêu \\

        Phạm vi
        & Thường áp dụng cho toàn bộ phiên bản ứng dụng trên môi trường đích
        & Có thể áp dụng cho toàn bộ hoặc một phần người dùng \\

        Khả năng đảo ngược
        & Có thể yêu cầu rollback phiên bản hoặc deploy lại
        & Có thể tắt tính năng bằng thay đổi cấu hình \\

        Thời điểm
        & Khi phiên bản phần mềm được đưa lên môi trường chạy
        & Khi tổ chức quyết định cho phép người dùng truy cập tính năng \\
        \bottomrule
    \end{tabular}
\end{table}

Nếu deploy và release bị ràng buộc chặt với nhau, mọi lần đưa mã mới lên production đều đồng thời làm thay đổi hành vi của toàn bộ hệ thống. Điều này làm tăng phạm vi ảnh hưởng khi xảy ra lỗi và khiến đội phát triển phải phụ thuộc vào các khung thời gian phát hành cố định. Ngược lại, khi hai hoạt động được tách biệt, việc triển khai kỹ thuật có thể diễn ra thường xuyên, còn quyết định phát hành được thực hiện độc lập dựa trên mức độ sẵn sàng và rủi ro nghiệp vụ.

\subsubsection{Tách Deploy khỏi Release bằng Feature Flag}

Feature Flag cung cấp một lớp điều khiển nằm giữa mã nguồn đã được deploy và hành vi được cung cấp cho người dùng. Mã của tính năng mới được triển khai cùng ứng dụng nhưng được đặt phía sau một flag. Khi flag ở trạng thái tắt, ứng dụng tiếp tục thực hiện hành vi hiện tại. Khi flag được bật, ứng dụng chuyển sang nhánh xử lý mới mà không cần deploy lại \cite{fowler_feature_toggles}.

Quy trình tách deploy khỏi release có thể được mô tả qua các bước sau:

\begin{enumerate}
    \item Đội phát triển triển khai mã nguồn chứa cả hành vi hiện tại và hành vi mới.
    \item Tính năng mới được bảo vệ bởi một Feature Flag và mặc định chưa được cung cấp rộng rãi.
    \item Tính năng được kiểm tra trên production với người dùng nội bộ hoặc một nhóm giới hạn.
    \item Phạm vi release được mở rộng theo user, role hoặc tỷ lệ phần trăm người dùng.
    \item Nếu phát hiện sự cố, flag được tắt để khôi phục hành vi cũ mà không cần rollback toàn bộ phiên bản.
    \item Khi tính năng đã ổn định và được phát hành hoàn toàn, flag cùng nhánh mã cũ có thể được loại bỏ.
\end{enumerate}

Ví dụ, mã nguồn của một giao diện thanh toán mới có thể được deploy lên production trong khi người dùng vẫn nhìn thấy giao diện hiện tại.



Sau khi hoàn tất kiểm thử, đội phát triển có thể bật flag cho nhân viên nội bộ, tiếp tục mở rộng cho khách hàng thử nghiệm và sau đó tăng dần tỷ lệ rollout. Cách tiếp cận này cho phép kiểm soát phạm vi ảnh hưởng của thay đổi, thu thập phản hồi sớm và giảm rủi ro khi phát hành \cite{launchdarkly_feature_flags,unleash_feature_flag_best_practices}.

Việc tách deploy khỏi release cũng hỗ trợ mô hình phát triển liên tục. Các thay đổi chưa hoàn thiện có thể được tích hợp vào nhánh mã chính nhưng chưa được hiển thị cho người dùng. Nhờ đó, đội phát triển giảm nhu cầu duy trì các nhánh tính năng tồn tại trong thời gian dài, hạn chế xung đột khi hợp nhất mã và có thể triển khai các thay đổi nhỏ thường xuyên hơn \cite{microsoft_devops}.

Tuy nhiên, Feature Flag không loại bỏ hoàn toàn nhu cầu rollback. Nếu lỗi nằm trong phần mã dùng chung, quá trình khởi động ứng dụng, thay đổi cơ sở dữ liệu hoặc hạ tầng, việc tắt flag có thể không đủ để khôi phục hệ thống. Do đó, Feature Flag nên được xem là cơ chế kiểm soát release và giảm phạm vi ảnh hưởng, không phải là phương án thay thế cho toàn bộ chiến lược rollback và khôi phục sự cố.

\subsection{Vai trò của Feature Flag trong vòng đời phần mềm}

\subsubsection{Hỗ trợ phát triển và tích hợp liên tục}

Feature Flag hỗ trợ quá trình phát triển và tích hợp liên tục bằng cách cho phép mã nguồn của tính năng chưa hoàn thiện được tích hợp sớm vào nhánh mã chính nhưng chưa được cung cấp cho người dùng. Thay vì duy trì các nhánh tính năng trong thời gian dài, đội phát triển có thể bảo vệ phần mã mới bằng một flag và giữ flag ở trạng thái tắt cho đến khi tính năng đạt mức độ sẵn sàng cần thiết.

Cách tiếp cận này giúp giảm xung đột khi hợp nhất mã, rút ngắn thời gian tồn tại của các nhánh phát triển và khuyến khích việc tích hợp các thay đổi nhỏ, thường xuyên. Các thành viên trong đội có thể cùng làm việc trên một phiên bản mã nguồn thống nhất, trong khi trạng thái phát hành của từng tính năng vẫn được kiểm soát độc lập \cite{microsoft_devops,fowler_feature_toggles}.

Feature Flag cũng hỗ trợ kiểm thử theo từng giai đoạn trong pipeline phát triển. Một tính năng có thể được bật trên môi trường development, tiếp tục được kiểm tra trên staging và vẫn giữ trạng thái tắt trên production. Cùng một phiên bản mã nguồn có thể được triển khai qua nhiều môi trường nhưng sử dụng cấu hình flag khác nhau, từ đó giảm sự khác biệt giữa các phiên bản phần mềm được kiểm thử và được vận hành thực tế.

Tuy nhiên, việc sử dụng Feature Flag trong phát triển liên tục đòi hỏi các nhánh xử lý cũ và mới đều phải được kiểm thử. Flag cần có mục đích, người chịu trách nhiệm và thời điểm loại bỏ rõ ràng để tránh tích lũy mã không còn cần thiết trong hệ thống \cite{unleash_feature_flag_best_practices,mahdavi_hezaveh_feature_toggles}.

\subsubsection{Giảm rủi ro khi phát hành}

Trong mô hình phát hành truyền thống, một tính năng mới thường được cung cấp đồng thời cho toàn bộ người dùng ngay sau khi deploy. Nếu tính năng chứa lỗi, phạm vi ảnh hưởng có thể lớn và đội vận hành phải rollback phiên bản hoặc thực hiện một lần triển khai khẩn cấp. Feature Flag giảm rủi ro này bằng cách cho phép kiểm soát phạm vi release độc lập với quá trình deploy.

Tính năng có thể được phát hành trước cho một nhóm nhỏ, chẳng hạn như nhân viên nội bộ, khách hàng thử nghiệm hoặc người dùng thuộc một vai trò cụ thể. Sau khi quan sát kết quả và xác nhận tính ổn định, đội phát triển có thể mở rộng phạm vi phát hành theo từng giai đoạn. Nếu xuất hiện lỗi, phạm vi ảnh hưởng được giới hạn trong nhóm người dùng đã được bật tính năng.

Rollout theo phần trăm là một chiến lược phổ biến để thực hiện phát hành từng bước. Hệ thống có thể bắt đầu với một tỷ lệ nhỏ người dùng, sau đó tăng dần từ \(5\%\), \(20\%\), \(50\%\) đến \(100\%\). Việc sử dụng cơ chế phân nhóm ổn định giúp một người dùng tiếp tục nhận cùng một kết quả qua nhiều lần truy cập, hạn chế trải nghiệm thay đổi ngẫu nhiên giữa hành vi cũ và mới \cite{launchdarkly_percentage_rollouts}.

Feature Flag không loại bỏ lỗi khỏi phần mềm, nhưng giúp giảm xác suất và mức độ ảnh hưởng của lỗi trong quá trình phát hành. Nhờ khả năng kiểm soát đối tượng và phạm vi tiếp cận, đội phát triển có thể đưa ra quyết định mở rộng hoặc thu hẹp release dựa trên dữ liệu quan sát được thay vì phát hành toàn bộ trong một lần.

\subsubsection{Hỗ trợ thử nghiệm trên môi trường thực tế}

Một số vấn đề chỉ xuất hiện trong điều kiện vận hành thực tế, nơi hệ thống phải xử lý dữ liệu, lưu lượng truy cập, cấu hình hạ tầng và hành vi người dùng đa dạng hơn môi trường kiểm thử. Feature Flag cho phép đội phát triển kiểm tra tính năng trên production nhưng giới hạn quyền truy cập cho một nhóm đối tượng được lựa chọn.

Ví dụ, một tính năng mới có thể chỉ được bật cho tài khoản nội bộ hoặc nhóm khách hàng beta. Nhóm này sử dụng cùng hạ tầng và dữ liệu vận hành như người dùng thông thường nhưng được phép truy cập hành vi mới. Qua đó, đội phát triển có thể phát hiện các vấn đề về hiệu năng, khả năng tương thích, trải nghiệm người dùng hoặc logic nghiệp vụ trước khi phát hành rộng rãi.

Feature Flag cũng là cơ sở để thực hiện các thử nghiệm so sánh. Những nhóm người dùng khác nhau có thể được phân phối đến các biến thể khác nhau của một tính năng, sau đó hệ thống thu thập các chỉ số để đánh giá hiệu quả. Khác với việc chỉ bật hoặc tắt tính năng, experiment flag thường cần quản lý biến thể, phân nhóm ổn định và dữ liệu đo lường để bảo đảm kết quả thử nghiệm có ý nghĩa \cite{fowler_feature_toggles}.

Việc thử nghiệm trên production cần được kiểm soát cẩn thận. Đội phát triển phải xác định rõ đối tượng tham gia, tiêu chí thành công, thời gian thử nghiệm và phương án dừng. Các dữ liệu nhạy cảm không nên được đưa trực tiếp vào cấu hình flag, đồng thời các thay đổi cấu hình cần được ghi nhận để có thể truy vết khi kết quả thử nghiệm bất thường \cite{unleash_feature_flag_best_practices}.

\subsubsection{Tăng khả năng phản ứng khi xảy ra sự cố}

Feature Flag giúp rút ngắn thời gian phản ứng khi một tính năng gây lỗi trên môi trường production. Thay vì phải sửa mã nguồn, xây dựng lại ứng dụng và thực hiện một lần deploy mới, đội vận hành có thể tắt flag để ngừng cung cấp tính năng và đưa người dùng trở lại hành vi ổn định trước đó.

Đối với các tính năng có nguy cơ ảnh hưởng lớn đến hệ thống, flag có thể được sử dụng như một kill switch. Kill switch là cơ chế ưu tiên cao cho phép vô hiệu hóa nhanh một tính năng hoặc một nhóm tính năng khi xuất hiện sự cố về hiệu năng, bảo mật, dữ liệu hoặc phụ thuộc bên ngoài. Cơ chế này đặc biệt hữu ích khi thời gian phục hồi dịch vụ là yếu tố quan trọng.

Ví dụ, nếu một dịch vụ gợi ý sản phẩm mới làm tăng thời gian phản hồi của hệ thống, đội vận hành có thể tắt flag tương ứng để ứng dụng bỏ qua dịch vụ này và sử dụng hành vi dự phòng.



Khả năng tắt nhanh giúp giảm thời gian gián đoạn và phạm vi ảnh hưởng, nhưng chỉ hiệu quả khi nhánh xử lý dự phòng vẫn hoạt động đúng. Vì vậy, hành vi khi flag tắt cần được thiết kế và kiểm thử như một phần chính thức của hệ thống, không nên được xem là nhánh mã tạm thời.

Mọi thao tác bật, tắt hoặc thay đổi quy tắc trong tình huống sự cố cũng cần được ghi lại trong audit log, bao gồm người thực hiện, thời điểm, đối tượng bị thay đổi và cấu hình trước và sau thao tác. Thông tin này hỗ trợ phân tích nguyên nhân, truy vết quyết định và đánh giá quá trình xử lý sau sự cố.

Feature Flag chủ yếu xử lý sự cố liên quan đến hành vi được đặt phía sau flag. Nếu lỗi nằm trong phần mã khởi động, cơ sở dữ liệu, hạ tầng hoặc thành phần dùng chung, việc tắt flag có thể không đủ để phục hồi hệ thống. Do đó, Feature Flag cần được kết hợp với giám sát, cảnh báo, rollback phiên bản và các cơ chế phục hồi khác trong chiến lược vận hành tổng thể.

\subsection{Các loại Feature Flag phổ biến}

Feature Flag có thể được phân loại theo mục đích sử dụng, thời gian tồn tại và đối tượng chịu trách nhiệm quản lý. Một cách phân loại phổ biến gồm release flag, experiment flag, ops flag và permission flag \cite{fowler_feature_toggles}. Việc xác định đúng loại flag giúp đội phát triển thiết kế quy tắc đánh giá, cơ chế giám sát và kế hoạch loại bỏ phù hợp.

\subsubsection{Release Flag}

Release flag được sử dụng để kiểm soát việc phát hành một tính năng mới sau khi mã nguồn của tính năng đã được deploy. Loại flag này cho phép đội phát triển tích hợp mã chưa hoàn thiện hoặc chưa sẵn sàng phát hành rộng rãi vào nhánh mã chính mà không làm thay đổi trải nghiệm của người dùng hiện tại.

Release flag thường có hai trạng thái cơ bản: tắt để tiếp tục sử dụng hành vi hiện tại và bật để sử dụng hành vi mới. Tuy nhiên, quá trình phát hành có thể được thực hiện theo nhiều giai đoạn, chẳng hạn như bật trước cho đội phát triển, người dùng nội bộ, khách hàng beta và sau đó mở rộng dần đến toàn bộ người dùng.

Ví dụ, một giao diện thanh toán mới có thể được triển khai lên production nhưng được đặt phía sau flag \texttt{new-checkout}. Khi flag chưa được bật, người dùng tiếp tục sử dụng giao diện cũ. Sau khi tính năng được kiểm thử đầy đủ, flag được bật cho toàn bộ người dùng và nhánh xử lý cũ có thể được loại bỏ.

Release flag thường là loại flag có thời gian tồn tại ngắn hoặc trung hạn. Sau khi tính năng đã được phát hành hoàn toàn và ổn định, flag cùng các nhánh mã không còn cần thiết nên được loại bỏ để giảm độ phức tạp của hệ thống \cite{unleash_feature_flag_best_practices}.

\subsubsection{Experiment Flag}

Experiment flag được sử dụng để phân phối các biến thể khác nhau của một tính năng cho những nhóm người dùng khác nhau nhằm đo lường và so sánh kết quả. Loại flag này thường được áp dụng trong thử nghiệm A/B, kiểm chứng giả thuyết sản phẩm hoặc đánh giá tác động của một thay đổi đối với hành vi người dùng.

Khác với release flag, kết quả của experiment flag không nhất thiết chỉ gồm hai trạng thái bật và tắt. Một flag có thể trả về nhiều biến thể, chẳng hạn như \texttt{control}, \texttt{variant-a} và \texttt{variant-b}. Mỗi biến thể đại diện cho một hành vi hoặc giao diện khác nhau.

Việc phân nhóm người dùng cần bảo đảm tính ổn định. Cùng một người dùng nên tiếp tục nhận cùng một biến thể trong suốt thời gian thử nghiệm, trừ khi cấu hình được thay đổi có chủ đích. Điều này thường được thực hiện bằng cách áp dụng hàm băm ổn định trên khóa định danh của người dùng và khóa của flag.

Kết quả của experiment flag cần được kết hợp với dữ liệu đo lường, chẳng hạn như tỷ lệ chuyển đổi, thời gian sử dụng, số lần tương tác hoặc tỷ lệ lỗi. Chỉ riêng kết quả đánh giá flag không đủ để xác định biến thể nào hiệu quả hơn. Vì vậy, loại flag này thường gắn với hệ thống phân tích và quy trình đánh giá thống kê \cite{fowler_feature_toggles}.

Experiment flag thường có thời gian tồn tại giới hạn theo thời gian thử nghiệm. Sau khi thu thập đủ dữ liệu và xác định được biến thể phù hợp, thử nghiệm cần được kết thúc, biến thể được lựa chọn trở thành hành vi chính và các nhánh mã còn lại nên được loại bỏ.

\subsubsection{Ops và Kill-Switch Flag}

Ops flag được sử dụng để kiểm soát hành vi vận hành của hệ thống, đặc biệt đối với các chức năng có ảnh hưởng đến hiệu năng, độ ổn định hoặc các dịch vụ phụ thuộc. Loại flag này cho phép đội vận hành thay đổi cách hệ thống hoạt động mà không cần sửa mã nguồn và deploy lại ứng dụng.

Ví dụ, một ops flag có thể được sử dụng để bật hoặc tắt tác vụ chạy nền, cơ chế gửi thông báo, dịch vụ gợi ý sản phẩm hoặc một tích hợp với hệ thống bên ngoài động mà không cần sửa mã nguồn và deploy lại ứng dụng.

Ví dụ, một ops flag có thể được sử dụng để bật hoặc tắt tác. Khi thành phần phụ thuộc gặp lỗi hoặc gây quá tải, đội vận hành có thể tạm thời vô hiệu hóa chức năng liên quan để duy trì hoạt động của các phần quan trọng hơn.

Kill-switch flag là một dạng ops flag có mức ưu tiên cao, được thiết kế để tắt nhanh một tính năng hoặc một nhóm tính năng khi xảy ra sự cố. Khi kill switch được kích hoạt, quyết định này thường được thực hiện trước các quy tắc nhắm mục tiêu khác. Do đó, ngay cả khi người dùng thỏa mãn user allowlist, role targeting hoặc percentage rollout, tính năng vẫn bị tắt.

Một kill switch hiệu quả cần có các đặc điểm sau:

\begin{itemize}
    \item Có thể được kích hoạt nhanh thông qua dashboard hoặc API quản trị.
    \item Có thứ tự ưu tiên rõ ràng trong quá trình đánh giá.
    \item Có hành vi dự phòng an toàn khi tính năng bị vô hiệu hóa.
    \item Có audit log ghi nhận người thực hiện, thời điểm và thay đổi cấu hình.
    \item Được kiểm thử định kỳ để bảo đảm cơ chế tắt khẩn cấp vẫn hoạt động.
\end{itemize}

Ops flag và kill-switch flag có thể tồn tại lâu hơn release flag nếu chúng đóng vai trò như một cơ chế kiểm soát vận hành lâu dài. Tuy nhiên, các flag không còn gắn với rủi ro thực tế hoặc không còn được sử dụng vẫn cần được rà soát và loại bỏ.

\subsubsection{Permission Flag}

Permission flag được sử dụng để xác định người dùng hoặc nhóm người dùng nào có quyền truy cập một tính năng. Quyết định đánh giá thường dựa trên vai trò, gói dịch vụ, loại tài khoản, tổ chức hoặc các thuộc tính nghiệp vụ khác.

Ví dụ, một tính năng quản trị nâng cao có thể chỉ được bật cho người dùng có vai trò \texttt{admin}, trong khi một chức năng thử nghiệm có thể chỉ dành cho nhóm \texttt{beta-customer}. Tương tự, một tính năng trả phí có thể được bật cho khách hàng thuộc gói \texttt{premium} và tắt đối với người dùng gói miễn phí.

Permission flag có điểm tương đồng với cơ chế phân quyền, nhưng không nên được xem là phương án thay thế hoàn toàn cho authorization. Feature Flag chủ yếu quyết định tính năng có được cung cấp hay không, trong khi hệ thống phân quyền phải bảo vệ tài nguyên và kiểm tra quyền truy cập tại phía máy chủ. Việc chỉ ẩn giao diện bằng flag mà không kiểm tra quyền tại backend có thể tạo ra rủi ro bảo mật.

Permission flag thường có thời gian tồn tại dài vì nó có thể phản ánh chính sách sản phẩm hoặc mô hình kinh doanh ổn định. Do đó, loại flag này cần được đặt tên rõ ràng, quản lý quyền thay đổi chặt chẽ và kiểm thử như một phần của logic nghiệp vụ.

\subsubsection{Vòng đời và thời gian tồn tại của từng loại Flag}

Mỗi loại Feature Flag có mục đích và thời gian tồn tại khác nhau. Việc áp dụng cùng một quy trình quản lý cho mọi loại flag có thể dẫn đến tình trạng các flag tạm thời tồn tại quá lâu hoặc các flag dài hạn bị loại bỏ không phù hợp.

Vòng đời tổng quát của một Feature Flag thường gồm các giai đoạn sau:

\begin{enumerate}
    \item \textbf{Tạo flag:} xác định mục đích, người chịu trách nhiệm, giá trị mặc định và phạm vi sử dụng.
    \item \textbf{Tích hợp:} thêm điểm kiểm tra flag vào mã nguồn và cấu hình các quy tắc đánh giá.
    \item \textbf{Triển khai và vận hành:} sử dụng flag để kiểm soát phát hành, thử nghiệm, quyền truy cập hoặc hành vi vận hành.
    \item \textbf{Ổn định:} theo dõi kết quả, xử lý sự cố và điều chỉnh phạm vi áp dụng.
    \item \textbf{Kết thúc:} lưu trữ, vô hiệu hóa hoặc loại bỏ flag cùng các nhánh mã không còn cần thiết.
\end{enumerate}

Bảng~\ref{tab:feature-flag-lifecycle} trình bày đặc điểm vòng đời của các loại flag phổ biến.

\begin{table}[H]
    \centering
    \caption{Vòng đời của các loại Feature Flag phổ biến}
    \label{tab:feature-flag-lifecycle}
    \begin{tabular}{p{0.20\textwidth} p{0.25\textwidth} p{0.20\textwidth} p{0.25\textwidth}}
        \toprule
        \textbf{Loại flag} &
        \textbf{Mục đích chính} &
        \textbf{Thời gian tồn tại} &
        \textbf{Điều kiện kết thúc} \\
        \midrule

        Release flag &
        Kiểm soát quá trình phát hành tính năng mới &
        Ngắn hoặc trung hạn &
        Tính năng đã phát hành hoàn toàn và hành vi cũ không còn cần thiết \\

        Experiment flag &
        So sánh các biến thể và kiểm chứng giả thuyết &
        Ngắn hạn &
        Thử nghiệm kết thúc và đã lựa chọn được biến thể chính thức \\

        Ops/Kill-switch flag &
        Kiểm soát rủi ro vận hành và phản ứng sự cố &
        Trung hoặc dài hạn &
        Rủi ro không còn tồn tại hoặc cơ chế kiểm soát đã được thay thế \\

        Permission flag &
        Kiểm soát quyền truy cập theo vai trò hoặc chính sách sản phẩm &
        Dài hạn &
        Chính sách quyền hoặc mô hình sản phẩm thay đổi \\
        \bottomrule
    \end{tabular}
\end{table}

Release flag và experiment flag thường được xem là các flag tạm thời. Nếu không được loại bỏ sau khi hoàn thành mục đích, chúng tạo ra nhiều nhánh xử lý dư thừa, làm tăng số trường hợp cần kiểm thử và gây khó khăn khi bảo trì. Tình trạng này được gọi là \textit{feature flag debt} \cite{mahdavi_hezaveh_feature_toggles}.

Ngược lại, ops flag và permission flag có thể được duy trì lâu dài nếu chúng tiếp tục mang lại giá trị vận hành hoặc nghiệp vụ. Tuy nhiên, flag dài hạn vẫn cần được rà soát định kỳ để xác nhận rằng cấu hình, người sở hữu và mục đích sử dụng còn phù hợp.

Để quản lý vòng đời hiệu quả, mỗi flag nên có metadata mô tả loại flag, người chịu trách nhiệm, ngày tạo, thời điểm dự kiến rà soát và trạng thái vòng đời. Hệ thống cũng nên hỗ trợ lưu trữ flag thay vì xóa trực tiếp, ghi nhận lịch sử thay đổi và cảnh báo các flag đã tồn tại vượt quá thời gian dự kiến \cite{unleash_feature_flag_best_practices}.

\subsection{Các chiến lược phát hành tính năng}

Feature Flag cho phép đội phát triển kiểm soát phạm vi phát hành tính năng theo nhiều chiến lược khác nhau. Thay vì chỉ bật hoặc tắt tính năng cho toàn bộ hệ thống, flag có thể được đánh giá dựa trên danh tính, vai trò hoặc nhóm phân phối của người dùng. Việc lựa chọn chiến lược phù hợp phụ thuộc vào mục tiêu phát hành, mức độ rủi ro và loại thông tin có trong ngữ cảnh đánh giá.

\subsubsection{Bật hoặc tắt toàn cục}

Bật hoặc tắt toàn cục là chiến lược đơn giản nhất, trong đó một tính năng được cung cấp cho toàn bộ người dùng hoặc bị vô hiệu hóa hoàn toàn. Kết quả đánh giá không phụ thuộc vào danh tính, vai trò hay thuộc tính của người dùng.

Khi flag được bật toàn cục, mọi yêu cầu đánh giá hợp lệ đều nhận kết quả bật. Chiến lược này phù hợp khi tính năng đã được kiểm thử đầy đủ và sẵn sàng phát hành rộng rãi. Ngược lại, trạng thái tắt toàn cục được sử dụng khi tính năng chưa sẵn sàng, đã kết thúc hoặc cần được vô hiệu hóa nhanh do xảy ra sự cố.

Bật hoặc tắt toàn cục thường được sử dụng trong các trường hợp sau:

\begin{itemize}
    \item Phát hành hoàn toàn một tính năng đã ổn định.
    \item Ẩn một tính năng chưa hoàn thiện nhưng mã nguồn đã được deploy.
    \item Tạm dừng một tính năng gây lỗi hoặc ảnh hưởng đến hiệu năng.
    \item Kiểm soát các chức năng vận hành như tác vụ nền hoặc tích hợp bên ngoài.
\end{itemize}

Do không cần xử lý nhiều điều kiện, chiến lược này có chi phí đánh giá thấp và hành vi dễ dự đoán. Tuy nhiên, nó không giới hạn được phạm vi ảnh hưởng khi tính năng được bật, vì toàn bộ người dùng sẽ nhận cùng một kết quả.

\subsubsection{Targeting theo người dùng}

Targeting theo người dùng cho phép bật tính năng cho một tập hợp đối tượng cụ thể dựa trên khóa định danh của họ. Khóa này có thể là mã người dùng, mã khách hàng hoặc một \texttt{targetingKey} ổn định do ứng dụng tích hợp cung cấp trong ngữ cảnh đánh giá \cite{openfeature_evaluation_context}.

Chiến lược này thường được triển khai dưới dạng danh sách cho phép. Hệ thống so sánh khóa định danh trong ngữ cảnh với danh sách đã được cấu hình cho flag. Nếu tìm thấy giá trị phù hợp, tính năng được bật; nếu không, hệ thống tiếp tục đánh giá quy tắc tiếp theo hoặc trả về giá trị mặc định.

Targeting theo người dùng phù hợp trong các trường hợp như:

\begin{itemize}
    \item Bật tính năng cho thành viên đội phát triển hoặc kiểm thử.
    \item Cho phép một số khách hàng tham gia chương trình beta.
    \item Kiểm tra lỗi với một tài khoản cụ thể trên production.
    \item Cung cấp tính năng sớm cho một nhóm khách hàng chiến lược.
\end{itemize}

Ví dụ, một flag có thể được cấu hình bật cho các khóa \texttt{user-001}, \texttt{user-024} và \texttt{user-105}. Khi một trong các đối tượng này gửi yêu cầu, hệ thống trả về trạng thái bật bất kể phần lớn người dùng còn lại đang nhận trạng thái tắt.

Khóa dùng để targeting cần ổn định và duy nhất trong phạm vi phù hợp. Không nên phụ thuộc vào các giá trị có thể thay đổi thường xuyên, chẳng hạn như tên hiển thị hoặc địa chỉ email chưa được chuẩn hóa. Đồng thời, dữ liệu định danh nhạy cảm không nên được đưa trực tiếp vào cấu hình nếu không thực sự cần thiết \cite{unleash_feature_flag_best_practices}.

\subsubsection{Targeting theo vai trò}

Targeting theo vai trò quyết định trạng thái tính năng dựa trên vai trò nghiệp vụ của người dùng, chẳng hạn như \texttt{admin}, \texttt{beta-customer} hoặc \texttt{regular-customer}. Ứng dụng tích hợp cung cấp danh sách vai trò trong ngữ cảnh đánh giá, sau đó hệ thống kiểm tra xem ít nhất một vai trò có phù hợp với quy tắc đã cấu hình hay không.

Chiến lược này hữu ích khi quyền sử dụng tính năng gắn với một nhóm người dùng thay vì từng cá nhân. So với targeting theo người dùng, việc cấu hình theo vai trò dễ quản lý hơn khi số lượng người dùng lớn và thành viên của nhóm thường xuyên thay đổi.

Ví dụ, một tính năng quản trị mới có thể được bật cho vai trò \texttt{admin}, trong khi một giao diện thử nghiệm được bật cho vai trò \texttt{beta-customer}. Khi vai trò của người dùng thay đổi trong hệ thống nguồn, kết quả đánh giá flag cũng thay đổi tương ứng mà không cần cập nhật từng định danh người dùng trong cấu hình.

Một quy tắc targeting theo vai trò có thể được mô tả như sau.



Cần phân biệt targeting theo vai trò với cơ chế phân quyền bảo mật. Vai trò trong Feature Flag được sử dụng để quyết định cách phát hành tính năng, nhưng backend vẫn phải kiểm tra quyền truy cập đối với các tài nguyên và thao tác nhạy cảm. Việc flag trả về trạng thái bật không tự động chứng minh rằng người dùng có quyền thực hiện mọi hành động liên quan.

\subsubsection{Rollout theo phần trăm người dùng}

Rollout theo phần trăm cho phép phát hành tính năng cho một tỷ lệ xác định trong tổng số người dùng đủ điều kiện. Chiến lược này giúp đội phát triển mở rộng phạm vi phát hành từng bước thay vì bật tính năng đồng thời cho toàn bộ hệ thống.

Một quy trình rollout có thể bắt đầu với \(5\%\) người dùng, sau đó tăng lên \(20\%\), \(50\%\) và cuối cùng là \(100\%\). Sau mỗi giai đoạn, đội phát triển theo dõi tỷ lệ lỗi, thời gian phản hồi, phản hồi của người dùng và các chỉ số nghiệp vụ trước khi quyết định tiếp tục mở rộng.

Để bảo đảm kết quả ổn định, hệ thống không nên lựa chọn ngẫu nhiên lại người dùng trong mỗi lần đánh giá. Thay vào đó, một hàm băm ổn định được áp dụng trên khóa của flag và khóa định danh của đối tượng. Giá trị băm sau đó được ánh xạ vào một khoảng bucket cố định. Tính năng được bật nếu bucket của người dùng nằm trong tỷ lệ rollout đã cấu hình \cite{launchdarkly_percentage_rollouts}.

Cơ chế này có thể được biểu diễn khái quát như sau:

\[
b = H(\textit{flagKey}, \textit{targetingKey}) \bmod 10000
\]

Trong đó \(H\) là hàm băm ổn định và \(b\) là bucket của người dùng trong khoảng từ \(0\) đến \(9999\). Với tỷ lệ rollout \(p\%\), tính năng được bật khi:

\[
b < p \times 100
\]

Ví dụ, với rollout \(20\%\), các bucket từ \(0\) đến \(1999\) nhận trạng thái bật. Do cùng một đầu vào tạo ra cùng một giá trị băm, một người dùng tiếp tục nhận cùng một kết quả giữa nhiều lần truy cập, miễn là khóa định danh, khóa flag và thuật toán phân bucket không thay đổi.

Rollout theo phần trăm giúp giới hạn phạm vi ảnh hưởng nhưng không bảo đảm rằng tỷ lệ quan sát được trong một tập người dùng nhỏ luôn chính xác tuyệt đối. Tỷ lệ cấu hình thể hiện tỷ lệ bucket mục tiêu; khi số lượng người dùng đủ lớn, tỷ lệ thực tế thường tiến gần đến giá trị này.

\subsubsection{Kết hợp nhiều chiến lược Targeting}

Trong thực tế, một flag thường cần kết hợp nhiều chiến lược để đáp ứng quy trình phát hành theo từng giai đoạn. Ví dụ, tính năng có thể được bật cho một số tài khoản kiểm thử cụ thể, tiếp tục bật cho toàn bộ khách hàng beta và sau đó rollout cho \(20\%\) người dùng còn lại.

Khi nhiều quy tắc cùng tồn tại, hệ thống cần xác định thứ tự đánh giá rõ ràng. Một thứ tự phổ biến có thể gồm:

\begin{enumerate}
    \item Kiểm tra trạng thái vòng đời và trạng thái bật hoặc tắt của flag.
    \item Kiểm tra kill switch hoặc các cơ chế vô hiệu hóa ưu tiên cao.
    \item Trả về trạng thái bật nếu flag đang ở chế độ bật toàn cục.
    \item Kiểm tra danh sách người dùng được cho phép.
    \item Kiểm tra các quy tắc targeting theo vai trò.
    \item Kiểm tra rollout theo phần trăm.
    \item Trả về giá trị mặc định nếu không có quy tắc nào phù hợp.
\end{enumerate}

Với cách tổ chức này, quy tắc cụ thể thường được đánh giá trước quy tắc có phạm vi rộng hơn. Người dùng nằm trong allowlist có thể được bật tính năng ngay cả khi bucket phần trăm của họ không thuộc nhóm rollout. Tương tự, người dùng có vai trò beta có thể được truy cập tính năng trước khi tính năng được phát hành cho người dùng thông thường.

Ví dụ, một chiến lược phát hành kết hợp có thể được cấu hình như sau:

\begin{itemize}
    \item Bật cho các tài khoản nội bộ được xác định bằng \texttt{targetingKey}.
    \item Bật cho người dùng có vai trò \texttt{beta-customer}.
    \item Bật cho \(10\%\) người dùng còn lại bằng stable hashing.
    \item Tắt mặc định đối với các đối tượng không thỏa mãn bất kỳ điều kiện nào.
\end{itemize}

Các quy tắc cùng loại cũng cần có mức độ ưu tiên xác định. Bộ máy đánh giá thường sắp xếp quy tắc theo trường \texttt{priority} và dừng tại quy tắc đầu tiên phù hợp \cite{launchdarkly_targeting_rules}. Vì vậy, thay đổi thứ tự quy tắc có thể làm thay đổi kết quả ngay cả khi nội dung từng quy tắc không đổi.

Khi kết hợp nhiều chiến lược, hệ thống nên trả về không chỉ trạng thái bật hoặc tắt mà còn cả lý do đánh giá, chẳng hạn như \texttt{USER\_ALLOWLIST}, \texttt{ROLE\_MATCH}, \texttt{PERCENTAGE\_ROLLOUT} hoặc \texttt{DEFAULT\_OFF}. Thông tin này giúp giải thích quyết định, hỗ trợ kiểm thử và rút ngắn thời gian phân tích khi hành vi thực tế không đúng như mong đợi \cite{openfeature_flag_evaluation}.

\subsection{Lợi ích và rủi ro của Feature Flag}

Feature Flag mang lại khả năng kiểm soát linh hoạt đối với quá trình phát hành và vận hành phần mềm. Tuy nhiên, việc bổ sung các điểm rẽ nhánh động vào hệ thống cũng làm tăng độ phức tạp của mã nguồn, cấu hình và quy trình quản trị. Vì vậy, Feature Flag chỉ phát huy hiệu quả khi được sử dụng có mục đích, có cơ chế giám sát và có kế hoạch quản lý vòng đời rõ ràng \cite{unleash_feature_flag_best_practices,mahdavi_hezaveh_feature_toggles}.

\subsubsection{Tăng tốc độ phát hành}

Feature Flag giúp tăng tốc độ phát hành bằng cách tách quá trình deploy mã nguồn khỏi quyết định release tính năng. Đội phát triển có thể tích hợp và triển khai mã của tính năng mới lên môi trường production trong khi vẫn giữ tính năng ở trạng thái tắt đối với người dùng cuối. Khi tính năng đã sẵn sàng, việc phát hành được thực hiện bằng cách thay đổi cấu hình flag thay vì xây dựng và triển khai một phiên bản mới \cite{fowler_feature_toggles,microsoft_devops}.

Cơ chế này cho phép các thay đổi nhỏ được tích hợp thường xuyên vào nhánh mã chính, giảm thời gian tồn tại của các nhánh phát triển dài hạn và hạn chế xung đột khi hợp nhất mã. Nhiều tính năng có thể được triển khai trong cùng một phiên bản ứng dụng nhưng được phát hành ở các thời điểm khác nhau, tùy theo mức độ hoàn thiện và yêu cầu nghiệp vụ.

Feature Flag cũng giảm sự phụ thuộc giữa đội phát triển và đội sản phẩm. Đội kỹ thuật có thể hoàn tất deploy theo lịch trình kỹ thuật, trong khi đội sản phẩm lựa chọn thời điểm release phù hợp với kế hoạch kinh doanh, hoạt động truyền thông hoặc mức độ sẵn sàng của người dùng.

Tuy nhiên, tốc độ phát hành chỉ được cải thiện khi quy trình cấu hình flag đủ rõ ràng và đáng tin cậy. Nếu việc bật flag yêu cầu nhiều thao tác thủ công, thiếu kiểm thử hoặc không có cơ chế phê duyệt phù hợp, Feature Flag có thể chuyển rủi ro từ deploy sang quá trình quản lý cấu hình.

\subsubsection{Kiểm soát phạm vi ảnh hưởng}

Feature Flag cho phép giới hạn phạm vi người dùng chịu ảnh hưởng của một thay đổi. Thay vì phát hành đồng thời cho toàn bộ hệ thống, tính năng có thể được bật trước cho tài khoản nội bộ, người dùng beta, một vai trò cụ thể hoặc một tỷ lệ nhỏ người dùng.

Cách tiếp cận này giúp đội phát triển quan sát hành vi thực tế của tính năng trước khi mở rộng phạm vi phát hành. Nếu xuất hiện lỗi, chỉ một nhóm người dùng giới hạn bị ảnh hưởng. Nếu hệ thống hoạt động ổn định, phạm vi rollout có thể được tăng dần đến toàn bộ người dùng.

Các chiến lược targeting còn cho phép tổ chức kiểm soát phạm vi theo nhu cầu nghiệp vụ. Ví dụ, một tính năng có thể chỉ được cung cấp cho khách hàng thuộc một gói dịch vụ, một tổ chức cụ thể hoặc một nhóm có vai trò phù hợp. Nhờ đó, cùng một phiên bản ứng dụng có thể phục vụ nhiều nhóm người dùng với các hành vi khác nhau.

Việc giới hạn phạm vi ảnh hưởng không đồng nghĩa với việc loại bỏ hoàn toàn rủi ro. Các thay đổi trong mã dùng chung, cơ sở dữ liệu hoặc hạ tầng vẫn có thể ảnh hưởng đến toàn bộ hệ thống, ngay cả khi tính năng được đặt phía sau flag. Do đó, phạm vi bảo vệ của từng flag cần được xác định rõ trong quá trình thiết kế.

\subsubsection{Rollback và khôi phục dịch vụ}

Khi một tính năng gây lỗi trên production, Feature Flag cho phép vô hiệu hóa nhanh hành vi mới và đưa người dùng trở lại nhánh xử lý ổn định trước đó. Thao tác này thường chỉ yêu cầu thay đổi cấu hình và có thể được thực hiện nhanh hơn so với rollback phiên bản hoặc triển khai một bản sửa lỗi khẩn cấp.

Đối với các thành phần có rủi ro vận hành cao, flag có thể đóng vai trò như một kill switch. Khi kill switch được kích hoạt, tính năng bị tắt bất kể các quy tắc targeting khác có phù hợp hay không. Cơ chế này giúp đội vận hành giảm thời gian phản ứng và giới hạn tác động của sự cố.

Ví dụ, nếu một dịch vụ bên ngoài làm tăng thời gian phản hồi, ứng dụng có thể bỏ qua dịch vụ đó khi flag bị tắt và chuyển sang hành vi dự phòng.



Khả năng tắt flag không hoàn toàn tương đương với rollback. Rollback đưa toàn bộ ứng dụng hoặc một thành phần về phiên bản trước, trong khi tắt flag chỉ thay đổi nhánh hành vi được lựa chọn trong phiên bản hiện tại. Nếu lỗi nằm trong quá trình khởi động, migration cơ sở dữ liệu, thư viện dùng chung hoặc đoạn mã được thực thi trước khi đánh giá flag, việc tắt flag có thể không khôi phục được dịch vụ.

Vì vậy, Feature Flag nên được kết hợp với cơ chế rollback phiên bản, giám sát, cảnh báo và kế hoạch khôi phục sự cố. Nhánh xử lý khi flag tắt cũng cần được kiểm thử thường xuyên để bảo đảm vẫn hoạt động khi cần sử dụng.

\subsubsection{Độ phức tạp của các nhánh xử lý}

Mỗi Feature Flag bổ sung ít nhất một điểm rẽ nhánh vào mã nguồn. Khi nhiều flag cùng tác động đến một luồng xử lý, số lượng tổ hợp trạng thái có thể tăng nhanh. Với \(n\) flag nhị phân độc lập, về lý thuyết có thể tồn tại tối đa:
$
2^n
$
tổ hợp bật và tắt khác nhau.

Không phải mọi tổ hợp đều xuất hiện trong thực tế, nhưng đội phát triển vẫn phải xác định các trạng thái hợp lệ và kiểm thử các trường hợp quan trọng. Nếu hai hoặc nhiều flag phụ thuộc lẫn nhau, hành vi của hệ thống có thể trở nên khó dự đoán và khó truy vết.

Độ phức tạp còn xuất hiện trong cấu hình. Một flag có thể có nhiều quy tắc user targeting, role targeting và percentage rollout với thứ tự ưu tiên khác nhau. Kết quả đánh giá không chỉ phụ thuộc vào trạng thái của flag mà còn phụ thuộc vào ngữ cảnh người dùng, môi trường, thứ tự quy tắc và giá trị mặc định.

Để hạn chế độ phức tạp, điểm kiểm tra flag nên được đặt tại các ranh giới chức năng rõ ràng. Đội phát triển nên tránh kiểm tra cùng một flag tại quá nhiều vị trí hoặc lồng nhiều flag trong cùng một khối điều kiện. Bộ máy đánh giá cũng cần trả về lý do quyết định để hỗ trợ kiểm thử và phân tích lỗi.

\subsubsection{Feature Flag Debt}

Feature flag debt là khoản nợ kỹ thuật hình thành khi các flag tạm thời không được loại bỏ sau khi đã hoàn thành mục đích. Flag có thể đã được bật cho toàn bộ người dùng nhưng nhánh xử lý cũ vẫn tồn tại, hoặc thử nghiệm đã kết thúc nhưng các biến thể không còn sử dụng vẫn được giữ lại trong mã nguồn và cấu hình \cite{mahdavi_hezaveh_feature_toggles}.

Các flag lỗi thời làm tăng số lượng nhánh cần đọc, kiểm thử và bảo trì. Thành viên mới có thể không biết flag nào còn hiệu lực, flag nào có thể thay đổi và flag nào chỉ còn tồn tại do chưa được dọn dẹp. Theo thời gian, việc xóa flag trở nên khó khăn hơn vì không còn rõ những thành phần nào phụ thuộc vào nó.

Feature flag debt có thể dẫn đến các vấn đề sau:

\begin{itemize}
    \item Mã nguồn chứa nhiều nhánh xử lý không còn cần thiết.
    \item Số lượng trường hợp kiểm thử và cấu hình tăng lên.
    \item Khó xác định trạng thái thực tế của tính năng.
    \item Tăng nguy cơ bật lại nhầm một hành vi đã lỗi thời.
    \item Tăng chi phí bảo trì và phân tích sự cố.
\end{itemize}

Để kiểm soát feature flag debt, mỗi flag nên có loại, người chịu trách nhiệm, ngày tạo và thời điểm dự kiến rà soát. Release flag và experiment flag cần có tiêu chí kết thúc rõ ràng. Khi tính năng đã ổn định, đội phát triển nên loại bỏ nhánh cũ, điểm kiểm tra trong mã nguồn và cấu hình flag không còn sử dụng \cite{unleash_feature_flag_best_practices}.

\subsubsection{Rủi ro từ cấu hình sai}

Feature Flag chuyển một phần hành vi của hệ thống từ mã nguồn sang cấu hình có thể thay đổi tại thời điểm chạy. Điều này tạo ra sự linh hoạt nhưng cũng làm tăng rủi ro từ thao tác cấu hình sai.

Một lỗi cấu hình có thể bao gồm bật nhầm flag trên production, đặt sai tỷ lệ rollout, nhập sai định danh người dùng, thay đổi thứ tự ưu tiên của quy tắc hoặc kích hoạt nhầm kill switch. Không giống lỗi mã nguồn phải trải qua quá trình build và deploy, lỗi cấu hình có thể ảnh hưởng đến hệ thống ngay sau khi thay đổi được áp dụng.

Rủi ro cũng có thể xuất hiện do khác biệt cấu hình giữa các environment. Một flag hoạt động đúng trên staging nhưng có thể cho kết quả khác trên production nếu trạng thái, quy tắc hoặc dữ liệu ngữ cảnh không đồng nhất.

Để giảm rủi ro từ cấu hình sai, nền tảng quản lý Feature Flag nên hỗ trợ:

\begin{itemize}
    \item Validation dữ liệu và giới hạn giá trị hợp lệ.
    \item Hiển thị rõ environment đang được chỉnh sửa.
    \item Cảnh báo trước các thay đổi có phạm vi ảnh hưởng lớn.
    \item Phân quyền đối với thao tác tạo và thay đổi flag.
    \item Audit log lưu cấu hình trước và sau thay đổi.
    \item Versioning hoặc lịch sử để khôi phục cấu hình trước đó.
    \item Quy trình phê duyệt đối với các thay đổi quan trọng.
\end{itemize}

Hệ thống cũng cần xác định giá trị mặc định an toàn khi không thể tải cấu hình hoặc khi quá trình đánh giá xảy ra lỗi. Việc lựa chọn mặc định bật hay tắt phụ thuộc vào bản chất của tính năng, nhưng quyết định này phải được xác định có chủ đích thay vì phụ thuộc vào hành vi ngẫu nhiên.

\subsubsection{Rủi ro bảo mật và quyền truy cập}

Nền tảng Feature Flag có khả năng thay đổi hành vi của ứng dụng production nên các API và dashboard quản trị trở thành mục tiêu nhạy cảm. Nếu quyền truy cập không được kiểm soát, một người dùng không hợp lệ có thể bật tính năng chưa hoàn thiện, tắt chức năng quan trọng hoặc thay đổi phạm vi phát hành.

Các thao tác quản trị cần được bảo vệ bằng cơ chế xác thực và phân quyền. Vai trò quản trị viên có thể được phép tạo, lưu trữ và kích hoạt kill switch, trong khi developer chỉ được chỉnh sửa một số cấu hình và viewer chỉ có quyền xem. Nguyên tắc quyền tối thiểu giúp giới hạn hậu quả nếu một tài khoản bị lộ hoặc sử dụng sai mục đích.

Endpoint evaluation cũng cần được bảo vệ phù hợp với mô hình triển khai. Dữ liệu đầu vào phải được validation để ngăn chặn các giá trị không hợp lệ hoặc payload quá lớn. Giao tiếp giữa client và hệ thống Feature Flag cần sử dụng kênh mã hóa, đặc biệt khi evaluation context chứa thông tin liên quan đến người dùng \cite{owasp_rest_security}.

Evaluation context và cấu hình targeting không nên chứa dữ liệu cá nhân nhạy cảm nếu không cần thiết. Các thuộc tính như mật khẩu, token truy cập hoặc thông tin định danh trực tiếp không nên được lưu trong rule. Khi cần nhận diện người dùng, hệ thống nên sử dụng một targeting key ổn định với phạm vi và mục đích rõ ràng \cite{unleash_feature_flag_best_practices}.

Feature Flag cũng không nên được sử dụng như cơ chế authorization duy nhất. Việc flag trả về trạng thái tắt có thể ẩn một nút hoặc một màn hình, nhưng backend vẫn phải kiểm tra quyền trước khi thực hiện thao tác nhạy cảm. Nếu chỉ dựa vào flag ở phía client, người dùng có thể bỏ qua giao diện và gọi trực tiếp API được bảo vệ không đầy đủ.

Cuối cùng, mọi thay đổi quan trọng cần được ghi nhận trong audit log, bao gồm chủ thể thực hiện, thời điểm, hành động, đối tượng và kết quả thay đổi. Audit log hỗ trợ truy vết sự cố và trách nhiệm vận hành, nhưng không nên lưu token, mật khẩu hoặc dữ liệu nhạy cảm không cần thiết \cite{owasp_logging}.

\section{Khảo sát các nền tảng Feature Flag hiện có}

Phần này khảo sát hai nền tảng tiêu biểu là LaunchDarkly và Unleash nhằm rút ra các nguyên tắc thiết kế cần kế thừa cho mini project. Nội dung chỉ tập trung vào điểm khác biệt và bài học thiết kế, tránh lặp lại các khái niệm Feature Flag đã trình bày ở phần tổng quan.

\subsection{Nền tảng LaunchDarkly}

LaunchDarkly là nền tảng Feature Flag thương mại dạng dịch vụ được quản lý. Nền tảng tổ chức tài nguyên theo project, environment và flag; cung cấp SDK đa nền tảng; hỗ trợ targeting rule có thứ tự, percentage rollout, nhiều variation, audit log và quản trị quyền truy cập \cite{launchdarkly_feature_flags,launchdarkly_targeting_rules}.

Điểm đáng chú ý đối với mini project là mô hình ordered rule. Người quản trị có thể đặt rule cụ thể trước rule tổng quát, ví dụ user allowlist trước percentage rollout. Cách biểu diễn này giúp kết quả evaluation dễ giải thích vì engine dừng tại rule đầu tiên khớp.

Hạn chế chính khi tham chiếu LaunchDarkly cho mini project là phạm vi tính năng rất rộng. Các khả năng như segment phức tạp, nhiều variation, workflow phê duyệt và phân tích thử nghiệm chuyên sâu được xem là hướng mở rộng, không phải yêu cầu triển khai ban đầu.

\subsection{Nền tảng Unleash}

Unleash là nền tảng Feature Flag mã nguồn mở, có thể self-hosted hoặc dùng dịch vụ được quản lý. Nền tảng sử dụng activation strategy để quyết định flag có được bật hay không trong từng environment; SDK thường tải cấu hình về và đánh giá tại ứng dụng \cite{unleash_feature_flags,unleash_feature_flag_best_practices}.

Điểm đáng chú ý đối với mini project là khả năng tự triển khai và tư duy tách control plane khỏi runtime evaluation. Unleash cũng nhấn mạnh lifecycle của flag, metric evaluation và việc tránh để hệ thống Feature Flag trở thành phụ thuộc đồng bộ bắt buộc của mọi request nghiệp vụ.

Hạn chế của mô hình self-hosted là tổ chức phải tự chịu trách nhiệm về database, backup, monitoring, nâng cấp và mở rộng. Vì mini project cần dễ chạy và dễ trình diễn, hệ thống chỉ mô phỏng phần lõi bằng một backend duy nhất, PostgreSQL và cache đơn giản.

\subsection{So sánh và bài học rút ra}

Bảng~\ref{tab:tong-hop-so-sanh-launchdarkly-unleash} tóm tắt các điểm so sánh có ảnh hưởng trực tiếp đến thiết kế mini project.

\begin{table}[H]
    \centering
    \caption{Tổng hợp so sánh LaunchDarkly và Unleash}
    \label{tab:tong-hop-so-sanh-launchdarkly-unleash}
    \begin{tabular}{p{0.22\textwidth} p{0.30\textwidth} p{0.30\textwidth}}
        \toprule
        \textbf{Tiêu chí} & \textbf{LaunchDarkly} & \textbf{Unleash} \\
        \midrule
        Mô hình triển khai & Dịch vụ được quản lý & Managed hoặc self-hosted \\
        Đơn vị logic & Ordered targeting rule & Activation strategy và constraint \\
        Evaluation & Chủ yếu qua SDK & Chủ yếu qua SDK hoặc proxy/edge \\
        Rollout & Percentage rollout ổn định & Gradual rollout ổn định \\
        Audit/RBAC & Tích hợp phong phú & Phụ thuộc phiên bản/mô hình triển khai \\
        Bài học cho mini project & Precedence rõ, reason dễ giải thích & Tách plane, self-hosted workflow, lifecycle \\
        \bottomrule
    \end{tabular}
\end{table}

Từ khảo sát, mini project kế thừa các nguyên tắc sau:

\begin{itemize}
    \item tổ chức tài nguyên theo project, environment và flag;
    \item tách control plane quản trị khỏi data plane evaluation;
    \item dùng evaluation context thay vì hard-code điều kiện trong ứng dụng;
    \item áp dụng stable rollout bằng khóa không chứa PII;
    \item trả reason code để giải thích kết quả evaluation;
    \item ghi audit log cho mọi mutation quan trọng;
    \item cung cấp SDK và fallback an toàn cho ứng dụng tích hợp.
\end{itemize}

\subsection{Định hướng của hệ thống đề xuất}

Hệ thống đề xuất không cố gắng sao chép đầy đủ LaunchDarkly hoặc Unleash. Phạm vi được chọn là một lát cắt đủ nhỏ để triển khai trong mini project nhưng vẫn thể hiện được bản chất của nền tảng Feature Flag: quản lý cấu hình, đánh giá runtime, thay đổi hành vi ứng dụng, auditability và trình diễn release an toàn.

Các chức năng được kế thừa trực tiếp gồm project/environment, flag config theo environment, ordered rule, user/role targeting, percentage rollout, kill switch, audit log, statistics, SDK và Docker Compose. Các chức năng được đơn giản hóa gồm multi-variant testing, segment phức tạp, approval workflow, realtime streaming, multi-region deployment và phân quyền enterprise.

Điểm khác biệt của MiniProject là ưu tiên tính minh bạch và khả năng giải thích hơn độ rộng tính năng. Vì vậy, Rule Evaluation Engine dùng precedence cố định, response có reason code rõ ràng, audit log append-only và Dashboard tập trung vào các kịch bản demo quan trọng nhất.

\section{Phân tích bài toán và yêu cầu hệ thống}

Mục này xác định phạm vi bài toán trước khi đi vào thiết kế. Các mô tả dưới đây dừng ở mức yêu cầu; chi tiết kiến trúc, dữ liệu, API và kiểm thử được trình bày ở các mục riêng để tránh trùng lặp.

\subsection{Các đối tượng sử dụng hệ thống}

Hệ thống phục vụ bốn nhóm tác nhân: người quản trị, nhà phát triển, người xem và ứng dụng tích hợp Feature Flag.

\subsubsection{Quản trị viên}

Quản trị viên chịu trách nhiệm tạo project, cấu hình flag, quản lý rule/group, kích hoạt kill switch và theo dõi lịch sử thay đổi.

\subsubsection{Nhà phát triển}

Nhà phát triển sử dụng nền tảng để chuẩn bị rollout, kiểm thử cấu hình và tích hợp Evaluation API hoặc SDK vào ứng dụng.

\subsubsection{Người xem}

Người xem cần quan sát trạng thái flag, audit log và thống kê nhưng không được thực hiện mutation cấu hình.

\subsubsection{Ứng dụng tích hợp Feature Flag}

Ứng dụng tích hợp gửi evaluation request kèm context và sử dụng kết quả trả về để bật hoặc ẩn hành vi tương ứng.

\subsection{Các ca sử dụng chính}

Các ca sử dụng được gom theo chuỗi thao tác từ quản trị đến runtime evaluation.

\subsubsection{Quản lý project và environment}

Người dùng cần tạo project và sử dụng environment làm phạm vi cấu hình/evaluation.

\subsubsection{Quản lý Feature Flag}

Hệ thống cần hỗ trợ tạo, cập nhật, archive/restore và xem trạng thái của Feature Flag.

\subsubsection{Cấu hình chiến lược rollout}

Người dùng cần cấu hình serving mode, user allowlist, role targeting, percentage rollout và kill switch.

\subsubsection{Đánh giá Feature Flag}

Ứng dụng khách cần nhận quyết định bật/tắt có reason code từ Evaluation API.

\subsubsection{Theo dõi thay đổi và thống kê}

Hệ thống cần ghi lại mutation cấu hình và tổng hợp số lượt evaluation để hỗ trợ truy vết và trình diễn.

\subsection{Yêu cầu chức năng mức MVP}

MVP tập trung vào một luồng end-to-end hoàn chỉnh: quản trị flag, lưu dữ liệu, evaluate từ ứng dụng demo và kiểm chứng kết quả.

\subsubsection{Yêu cầu đối với Backend API}

Backend phải cung cấp Management API, Evaluation API, validation, lỗi chuẩn hóa, persistence, audit log và seed data.

\subsubsection{Yêu cầu đối với Frontend Dashboard}

Dashboard phải cho phép quản lý project/flag/rule, quan sát trạng thái, xem audit log và thống kê ở mức đủ cho demo.

\subsubsection{Yêu cầu đối với Demo Application}

Demo Application phải cho thấy thay đổi cấu hình trên Dashboard ảnh hưởng đến hành vi ứng dụng mà không cần deploy lại.

\subsubsection{Yêu cầu đối với dữ liệu và tài liệu}

Sản phẩm cần có schema/migration, dữ liệu mẫu, tài liệu thiết kế ngắn và hướng dẫn chạy để phục vụ chấm điểm.

\subsection{Yêu cầu phi chức năng}

Các yêu cầu phi chức năng bảo đảm hệ thống không chỉ chạy được mà còn có hành vi đáng tin cậy khi cấu hình phức tạp hơn.

\subsubsection{Tính đúng đắn và xác định của kết quả}

Cùng một cấu hình và context phải cho cùng một kết quả; percentage rollout phải dùng stable hashing.

\subsubsection{Hiệu năng evaluation}

Evaluation API cần có độ trễ thấp và có thể tận dụng snapshot cache mà không làm sai quyết định.

\subsubsection{Tính nhất quán của cấu hình}

Mutation cấu hình phải làm mới dữ liệu runtime liên quan và tránh để engine dùng snapshot lỗi thời quá lâu.

\subsubsection{Khả năng bảo trì và mở rộng}

Logic evaluation cần tách khỏi controller để có thể bổ sung rule hoặc SDK mới trong tương lai.

\subsubsection{Khả năng kiểm thử}

Các nhánh precedence, fallback, validation và persistence phải có test tự động hoặc kịch bản kiểm chứng rõ ràng.

\subsubsection{Bảo mật và khả năng truy vết}

Management API cần RBAC demo, audit log append-only và không tin tưởng dữ liệu context do client tự khai báo.

\subsection{Các yêu cầu khuyến khích}

Ngoài MVP, đề tài bổ sung các năng lực giúp sản phẩm gần hơn với nền tảng Feature Flag thực tế.

\subsubsection{Caching evaluation snapshot}

Cache snapshot cấu hình để giảm truy vấn lặp của Evaluation API.

\subsubsection{JavaScript SDK}

SDK giúp ứng dụng web gọi Evaluation API bằng một giao diện ổn định và fail-closed.

\subsubsection{Unit test cho Rule Evaluation Engine}

Engine cần test cho precedence, stable rollout và các trường hợp lỗi/thiếu context.

\subsubsection{Versioning và lịch sử cấu hình}

Lịch sử cấu hình được dựng từ audit log thay vì bảng version riêng.

\subsubsection{Phân quyền người dùng quản trị}

RBAC demo phân biệt admin, developer và viewer cho các thao tác quản trị.

\subsubsection{Thống kê số lần evaluation}

Hệ thống tổng hợp số lượt evaluation theo flag, kết quả và thời gian.

\subsubsection{Group kill switch}

Group kill switch cho phép vô hiệu hóa nhanh nhiều flag liên quan.

\subsubsection{Triển khai bằng Docker Compose}

Docker Compose cung cấp quy trình khởi chạy thống nhất cho database, backend, dashboard và demo.

\section{Kiến trúc tổng thể của hệ thống}

Phần này trình bày kiến trúc ở mức hệ thống. Các chi tiết về mô hình dữ liệu, Rule Evaluation Engine, API, cache và bảo mật được tách sang các mục chuyên sâu phía sau để tránh lặp lại nội dung.

\subsection{Các thành phần chính}

Hệ thống được tổ chức thành sáu thành phần: Admin Dashboard, Backend API, Rule Evaluation Engine, PostgreSQL/Cache, JavaScript SDK và Demo Application. Dashboard phục vụ quản trị cấu hình; Backend API là ranh giới truy cập duy nhất; Rule Evaluation Engine quyết định trạng thái flag tại runtime; database lưu cấu hình và audit log; cache tăng tốc truy xuất snapshot; SDK và Demo Application minh họa cách ứng dụng khách tích hợp Feature Flag.

\begin{figure}[H]
    \centering
    \includegraphics[width=0.95\textwidth]{\detokenize{Images/mermaid/Kiến trúc tổng thể của hệ thống.png}}
    \caption{Kiến trúc tổng thể của hệ thống Feature Flag Platform}
    \label{fig:system-overall-architecture}
\end{figure}


\subsubsection{Admin Dashboard}

Dashboard cung cấp giao diện quản trị cho project, Feature Flag, rule, flag group, audit log, history, RBAC demo và thống kê evaluation. Dashboard không tự quyết định kết quả evaluation mà gửi yêu cầu đến Backend API, nhờ đó logic nghiệp vụ và quyền truy cập vẫn được kiểm soát tập trung ở phía máy chủ.

\subsubsection{Backend API}

Backend API gom toàn bộ management endpoint và Evaluation API trong một service NestJS. Các endpoint quản trị xử lý validation, phân quyền, persistence và audit log; Evaluation API phục vụ ứng dụng tích hợp với response nhất quán gồm trạng thái bật/tắt và reason code.

\subsubsection{Rule Evaluation Engine}

Rule Evaluation Engine là lớp domain thuần, độc lập với controller. Engine nhận evaluation request, sử dụng snapshot cấu hình đã được chuẩn hóa và trả về quyết định xác định. Thứ tự ưu tiên chi tiết được trình bày trong mục thiết kế engine, vì vậy mục này chỉ xác định vị trí của engine trong kiến trúc tổng thể.

\subsubsection{Database và cache}

PostgreSQL là nguồn dữ liệu chính cho project, environment, flag, rule, group, audit log và thống kê aggregate. Cache chỉ được xem là tầng tăng tốc cho evaluation snapshot; dữ liệu trong cache có thể bị loại bỏ mà không làm mất tính đúng đắn vì database vẫn là nguồn sự thật.

\subsubsection{JavaScript SDK}

JavaScript SDK đóng gói cách gọi Evaluation API, chuẩn hóa timeout, lỗi mạng và fallback fail-closed cho ứng dụng web. SDK giúp Demo Application không phụ thuộc trực tiếp vào chi tiết endpoint và là cơ sở để mở rộng SDK cho các ngôn ngữ khác.

\subsubsection{Demo Application}

Demo Application mô phỏng một hệ thống thương mại điện tử nhỏ, sử dụng SDK để quyết định hiển thị tính năng theo kết quả evaluation. Vai trò của ứng dụng này là chứng minh giá trị thực tế của nền tảng thay vì chứa logic quản trị riêng.

\subsection{Control plane và data plane}

Kiến trúc được phân chia theo hai luồng trách nhiệm. \textit{Control plane} quản lý cấu hình và lịch sử thay đổi; \textit{data plane} xử lý evaluation ở runtime. Cách phân chia này phù hợp với đặc điểm của Feature Flag: thao tác quản trị ít hơn nhưng yêu cầu evaluation có thể xuất hiện liên tục từ ứng dụng tích hợp \cite{launchdarkly_feature_flags,unleash_feature_flag_best_practices}.

\begin{figure}[H]
    \centering
    \includegraphics[width=0.88\textwidth]{\detokenize{Images/mermaid/Control Plane - Data Plane.png}}
    \caption{Phân tách control plane và data plane trong hệ thống}
    \label{fig:control-plane-data-plane}
\end{figure}


\subsubsection{Vai trò của control plane}

Control plane bao gồm Dashboard, management endpoints, persistence và audit logging. Mọi mutation quan trọng đều đi qua validation, kiểm tra quyền, ghi database và tạo audit entry append-only.

\subsubsection{Vai trò của data plane}

Data plane bao gồm Evaluation API, Rule Evaluation Engine và snapshot cache. Luồng này được tối ưu cho quyết định nhanh, xác định và fail-closed khi thiếu cấu hình hoặc phát sinh lỗi.

\subsubsection{Luồng dữ liệu giữa hai thành phần}

Người quản trị thay đổi cấu hình qua Dashboard; Backend API lưu thay đổi, ghi audit log và invalidation snapshot liên quan. Ứng dụng khách gửi evaluation context qua SDK hoặc gọi trực tiếp Evaluation API; engine đọc snapshot hiện tại và trả kết quả cho ứng dụng.

\subsubsection{Lợi ích của việc tách biệt hai plane}

Sự tách biệt giúp giảm coupling giữa giao diện quản trị và runtime evaluation, cho phép tối ưu hai luồng theo tải khác nhau, đồng thời làm rõ ranh giới bảo mật: dữ liệu quản trị và lịch sử thay đổi không cần phơi bày cho ứng dụng khách.

\subsection{Luồng quản trị Feature Flag}

Luồng quản trị tập trung vào việc biến ý định phát hành thành cấu hình có kiểm soát. Các bước dưới đây chỉ mô tả trình tự kiến trúc; cấu trúc bảng dữ liệu và endpoint cụ thể được trình bày ở các mục sau.

\subsubsection{Tạo và chỉnh sửa Feature Flag}

Người quản trị khai báo flag key, tên và mô tả. Backend kiểm tra ràng buộc duy nhất trong project, tạo hoặc cập nhật bản ghi liên quan và trả kết quả cho Dashboard.

\subsubsection{Cấu hình trạng thái và serving mode}

Sau khi có flag, người quản trị cấu hình lifecycle/status và serving mode theo từng environment. Thiết kế này tách trạng thái quản lý của flag khỏi cách flag được phục vụ ở runtime.

\subsubsection{Cấu hình và sắp xếp rule}

Khi flag chạy ở chế độ targeted, người quản trị cấu hình các rule hỗ trợ bởi engine. Dashboard chỉ biểu diễn và gửi cấu hình; việc diễn giải rule thuộc trách nhiệm của Rule Evaluation Engine.

\subsubsection{Lưu audit log và invalidation cache}

Sau mỗi mutation, Backend API ghi audit log và làm mới snapshot liên quan. Nhờ vậy, hệ thống vừa có khả năng truy vết vừa hạn chế việc data plane sử dụng cấu hình cũ.

\subsection{Luồng đánh giá Feature Flag}

Luồng evaluation bắt đầu từ ứng dụng khách và kết thúc bằng một quyết định có reason code. Evaluation API không thay đổi cấu hình, không ghi audit log quản trị và không yêu cầu Dashboard tham gia trong runtime.

\begin{figure}[H]
    \centering
    \includegraphics[width=0.95\textwidth]{\detokenize{Images/mermaid/Luồng đánh giá Feature Flag.png}}
    \caption{Luồng đánh giá Feature Flag từ ứng dụng khách đến Rule Evaluation Engine}
    \label{fig:evaluation-flow}
\end{figure}


\subsection{Luồng tích hợp ứng dụng khách}

Ứng dụng khách có thể gọi trực tiếp Evaluation API hoặc sử dụng JavaScript SDK. Trong mini project, Demo Application dùng SDK để giảm mã lặp ở phía client và thể hiện cách tích hợp có thể tái sử dụng.

\subsubsection{Khởi tạo evaluation context}

Ứng dụng tạo context từ các định danh không nhạy cảm như targeting key, user ID mẫu, role và organization ID mẫu. Context chỉ là đầu vào cho engine, không phải bằng chứng phân quyền.

\subsubsection{Gọi trực tiếp Evaluation API}

Cách gọi trực tiếp phù hợp cho kiểm thử hoặc hệ thống muốn tự quản lý HTTP client. Client gửi project key, environment key, flag key và context, sau đó tự xử lý response.

\subsubsection{Gọi thông qua JavaScript SDK}

SDK bao bọc cùng contract nhưng bổ sung cấu hình base URL, timeout và fallback. Nội dung triển khai SDK cụ thể được trình bày trong chương kết quả.

\subsubsection{Áp dụng kết quả trong ứng dụng}

Ứng dụng chỉ sử dụng trường kết quả để bật hoặc ẩn hành vi giao diện tương ứng. Việc này giữ cho logic sản phẩm tách khỏi cách cấu hình flag được quản trị.

\subsubsection{Xử lý khi Evaluation API không khả dụng}

Khi không nhận được kết quả hợp lệ, client áp dụng mặc định an toàn là tắt tính năng. Quy tắc này tránh việc lỗi hạ tầng vô tình mở tính năng chưa sẵn sàng.

\subsection{Lựa chọn công nghệ}

Các công nghệ được chọn theo tiêu chí phù hợp phạm vi mini project, dễ kiểm thử và đủ gần với một hệ thống thực tế.

\subsubsection{Lựa chọn NestJS cho Backend API}

NestJS hỗ trợ cấu trúc module rõ ràng, dependency injection, DTO validation và test bằng Jest, phù hợp với backend REST cần tách controller, service và domain logic.

\subsubsection{Lựa chọn React cho giao diện web}

React phù hợp để xây dựng Dashboard và Demo Application dạng SPA, tái sử dụng component và phản hồi nhanh với trạng thái cấu hình thay đổi.

\subsubsection{Lựa chọn PostgreSQL và Prisma}

PostgreSQL cung cấp giao dịch và kiểu JSONB cho rule parameters/audit snapshot; Prisma giúp mô hình hóa schema, migration và truy cập dữ liệu có kiểu.

\subsubsection{Lựa chọn cache in-memory hoặc Redis}

In-memory cache đủ cho demo một instance; Redis là hướng mở rộng khi cần chia sẻ snapshot giữa nhiều backend instance. Thiết kế giữ cache là tùy chọn để không phụ thuộc vào hạ tầng ngoài.

\subsubsection{Lựa chọn Docker và Docker Compose}

Docker Compose gom backend, frontend, demo và database vào một quy trình khởi chạy thống nhất, thuận tiện cho chấm điểm và trình diễn.

\subsubsection{Các phương án thay thế và lý do không lựa chọn}

Các lựa chọn như microservice tách riêng evaluation service, message broker hoặc frontend framework khác chưa cần thiết trong phạm vi mini project. Hệ thống ưu tiên kiến trúc đơn giản, có ranh giới rõ và có thể mở rộng sau.

\section{Thiết kế dữ liệu}

Mô hình dữ liệu được thiết kế để lưu cấu hình Feature Flag, phục vụ Evaluation API và ghi nhận lịch sử thay đổi. Nội dung phần này chỉ mô tả cấu trúc lưu trữ; thứ tự đánh giá nằm ở mục Rule Evaluation Engine và contract HTTP nằm ở mục API.

\subsection{Mô hình dữ liệu tổng thể}

Các thực thể chính của hệ thống được tóm tắt trong Bảng~\ref{tab:quan-he-thuc-the}. PostgreSQL được lựa chọn vì hỗ trợ giao dịch, khóa ngoại và kiểu \texttt{JSONB} cho rule parameters hoặc audit snapshot \cite{postgresql_json_types}.

\begin{table}[H]
    \centering
    \caption{Các thực thể và quan hệ chính trong mô hình dữ liệu}
    \label{tab:quan-he-thuc-the}
    \begin{tabular}{p{0.26\textwidth} p{0.64\textwidth}}
        \toprule
        \textbf{Thực thể} & \textbf{Vai trò chính} \\
        \midrule
        \texttt{Project} & Phạm vi quản lý cao nhất của một ứng dụng/sản phẩm. \\
        \texttt{Environment} & Môi trường triển khai như staging hoặc production. \\
        \texttt{FeatureFlag} & Định nghĩa ổn định của flag trong project. \\
        \texttt{FlagEnvironmentConfig} & Cấu hình runtime của một flag trong một environment. \\
        \texttt{FlagRule} & Rule targeting thuộc một cấu hình runtime. \\
        \texttt{FlagGroup} & Nhóm các flag liên quan trong cùng project. \\
        \texttt{FlagGroupConfig} & Trạng thái group kill switch theo environment. \\
        \texttt{SampleUserContext} & Dữ liệu người dùng mẫu phục vụ demo/kiểm thử. \\
        \texttt{AuditLogEntry} & Bản ghi append-only cho mutation quản trị. \\
        \bottomrule
    \end{tabular}
\end{table}

\subsubsection{Các thực thể chính}

Các thực thể được tổ chức theo chuỗi: project chứa environment và flag; flag có cấu hình theo environment; cấu hình có rule; group và audit log bổ sung khả năng vận hành/truy vết.

\subsubsection{Quan hệ giữa các thực thể}

Quan hệ quan trọng nhất là \texttt{FeatureFlag} -- \texttt{FlagEnvironmentConfig} -- \texttt{Environment}, cho phép một flag có cấu hình khác nhau theo từng môi trường. Group và audit log đều nằm trong phạm vi project để tránh lẫn dữ liệu giữa các ứng dụng.

\subsubsection{Phân tách định nghĩa flag và cấu hình runtime}

\texttt{FeatureFlag} lưu danh tính, tên, mô tả và lifecycle; \texttt{FlagEnvironmentConfig} lưu enabled/disabled, serving mode và kill switch theo environment. Việc tách này tránh nhân bản flag khi production và staging cần cấu hình khác nhau.

\subsubsection{Sơ đồ thực thể liên kết}

Sơ đồ rút gọn trong Hình~\ref{fig:feature-flag-er-diagram} nhấn mạnh đường đi dữ liệu phục vụ evaluation.

\begin{figure}[H]
    \centering
    \begin{tikzpicture}[
        entity/.style={rectangle, draw, rounded corners, minimum width=3.1cm, minimum height=0.9cm, text width=3.0cm, align=center},
        relation/.style={-{Latex[length=2mm]}, thick},
        node distance=1.4cm and 1.7cm
    ]
        \node[entity] (project) {\textbf{Project}};
        \node[entity, below left=of project] (environment) {\textbf{Environment}};
        \node[entity, below right=of project] (flag) {\textbf{FeatureFlag}};
        \node[entity, below=2.4cm of project] (config) {\textbf{FlagEnvironmentConfig}};
        \node[entity, below=of config] (rule) {\textbf{FlagRule}};
        \node[entity, left=of config] (group) {\textbf{FlagGroup / Config}};
        \node[entity, right=of config] (audit) {\textbf{AuditLogEntry}};
        \draw[relation] (project) -- node[left] {\small 1:N} (environment);
        \draw[relation] (project) -- node[right] {\small 1:N} (flag);
        \draw[relation] (environment) -- node[left] {\small 1:N} (config);
        \draw[relation] (flag) -- node[right] {\small 1:N} (config);
        \draw[relation] (config) -- node[right] {\small 1:N} (rule);
        \draw[relation] (project) -- node[above] {\small 1:N} (group);
        \draw[relation] (project) -- node[above] {\small 1:N} (audit);
    \end{tikzpicture}
    \caption{Sơ đồ thực thể liên kết rút gọn của hệ thống Feature Flag}
    \label{fig:feature-flag-er-diagram}
\end{figure}

\subsection{Project và environment}

Project xác định phạm vi nghiệp vụ; environment xác định phạm vi triển khai. Khóa nghiệp vụ của hai thực thể này được dùng trong API, cache key và seed data.

\subsubsection{Thực thể Project}

\texttt{Project} gồm \texttt{id}, \texttt{key}, \texttt{name}, \texttt{description}, \texttt{createdAt} và \texttt{updatedAt}. \texttt{key} là duy nhất toàn hệ thống và ổn định trong vòng đời project.

\subsubsection{Thực thể Environment}

\texttt{Environment} gồm \texttt{id}, \texttt{projectId}, \texttt{key}, \texttt{name}, \texttt{description}, \texttt{isDefault} và timestamp. \texttt{key} là duy nhất trong phạm vi project.

\subsubsection{Environment mặc định}

Mỗi project có một environment mặc định để tạo cấu hình ban đầu cho flag mới và giúp demo chạy ngay sau seed.

\subsubsection{Ràng buộc environment trong project}

Environment luôn thuộc một project cụ thể; các truy vấn evaluation dùng cặp \texttt{projectKey}/\texttt{environmentKey} để tránh nhầm cấu hình giữa project.

\subsection{Feature Flag và cấu hình theo environment}

Thiết kế tách metadata của flag khỏi trạng thái runtime nhằm hỗ trợ nhiều environment.

\subsubsection{Thực thể FeatureFlag}

\texttt{FeatureFlag} gồm \texttt{id}, \texttt{projectId}, \texttt{key}, \texttt{name}, \texttt{description}, \texttt{status} và timestamp.

\subsubsection{Trạng thái vòng đời của flag}

Lifecycle status gồm active và archived. Archived giữ lịch sử nhưng không còn phục vụ runtime.

\subsubsection{Thực thể FlagEnvironmentConfig}

\texttt{FlagEnvironmentConfig} nối flag với environment và lưu \texttt{enabled}, \texttt{servingMode}, \texttt{killSwitch} cùng timestamp.

\subsubsection{Trạng thái cấu hình và serving mode}

Config disabled dừng flag trong environment; serving mode quyết định bật toàn cục hay xét rule khi config enabled.

\subsubsection{Cấu hình flag theo từng environment}

Cùng một flag có thể bật ở staging, rollout một phần ở production và disabled ở environment khác mà không cần tạo nhiều flag.

\subsection{Rule và tham số rule}

Rule được lưu riêng để có thể thay đổi chiến lược rollout mà không sửa định nghĩa flag.

\subsubsection{Thực thể FlagRule}

\texttt{FlagRule} gồm \texttt{id}, \texttt{flagConfigId}, \texttt{type}, \texttt{priority}, \texttt{enabled}, \texttt{parameters} và timestamp.

\subsubsection{Kiểu rule}

Mini project hỗ trợ ba kiểu: \texttt{USER\_ALLOWLIST}, \texttt{ROLE\_TARGETING} và \texttt{PERCENTAGE\_ROLLOUT}.

\subsubsection{Thứ tự ưu tiên của rule}

Database lưu \texttt{priority} cho rule cùng loại; type precedence của engine được định nghĩa riêng trong mục evaluation.

\subsubsection{Lưu trữ tham số bằng JSON}

\texttt{parameters} dùng \texttt{JSONB} để lưu danh sách user, danh sách role hoặc tỷ lệ rollout mà không cần tạo bảng phụ cho từng kiểu rule.

\subsubsection{Validation tham số theo loại rule}

Backend kiểm tra shape của \texttt{parameters}: user/role phải là danh sách chuỗi, percentage phải nằm trong miền hợp lệ.

\subsection{User context mẫu}

\texttt{SampleUserContext} phục vụ demo và kiểm thử, không phải nguồn danh tính bắt buộc của Evaluation API.

\subsubsection{Targeting key}

\texttt{targetingKey} là khóa ổn định cho percentage rollout và là duy nhất trong project demo.

\subsubsection{Các trường nghiệp vụ và attributes}

Dữ liệu mẫu có thể chứa \texttt{userId}, \texttt{role} và \texttt{attributes} như organization ID.

\subsubsection{Phạm vi sử dụng dữ liệu mẫu}

Dữ liệu này giúp chọn nhanh context trên Dashboard/Demo Application và không đại diện cho user store sản xuất.

\subsection{Flag group và group kill switch}

Flag group hỗ trợ quản lý các flag liên quan và cho phép ngắt nhanh cả nhóm trong một environment.

\subsubsection{Thực thể FlagGroup và FlagGroupConfig}

\texttt{FlagGroup} lưu định danh nhóm; \texttt{FlagGroupConfig} lưu kill switch của nhóm theo environment.

\subsubsection{Quan hệ giữa group và flag}

Một group có nhiều flag; trong phạm vi mini project, một flag thuộc tối đa một group.

\subsubsection{Trạng thái group kill switch theo environment}

Kill switch của group được lưu theo environment để tránh ảnh hưởng đồng thời staging và production.

\subsubsection{Audit log và cache invalidation cho group}

Mutation group tạo audit entry và invalidation snapshot của các flag thành viên bị ảnh hưởng.

\subsection{Audit log và lịch sử cấu hình}

Audit log là nguồn truy vết chính cho mọi mutation quản trị và cũng là cơ sở cho audit-backed history \cite{owasp_logging}.

\subsubsection{Cấu trúc AuditLogEntry}

\texttt{AuditLogEntry} gồm \texttt{id}, \texttt{projectId}, \texttt{actor}, \texttt{requestId}, \texttt{targetType}, \texttt{targetId}, \texttt{targetKey}, \texttt{action}, \texttt{before}, \texttt{after} và \texttt{createdAt}.

\subsubsection{Thông tin actor và request}

Actor và request ID được xác định ở backend, không lấy trực tiếp từ body request, để tránh giả mạo lịch sử.

\subsubsection{Snapshot trước và sau thay đổi}

Snapshot \texttt{before}/\texttt{after} dùng \texttt{JSONB}, đủ để so sánh thay đổi nhưng không lưu bí mật hoặc dữ liệu cá nhân không cần thiết.

\subsubsection{Tính chất append-only}

Ứng dụng chỉ thêm audit entry mới, không cung cấp luồng sửa hoặc xóa bản ghi đã tạo.

\subsubsection{Lịch sử cấu hình dựa trên audit log}

History API dựng lại diễn biến cấu hình từ các audit entry liên quan thay vì duy trì bảng version riêng.

\subsection{Ràng buộc toàn vẹn dữ liệu}

Ràng buộc dữ liệu được đặt ở cả Prisma/database và tầng validation để giữ cấu hình hợp lệ trước khi đi vào engine.

\subsubsection{Ràng buộc khóa và tính duy nhất}

\begin{table}[H]
    \centering
    \caption{Các ràng buộc duy nhất tiêu biểu}
    \label{tab:unique-constraints}
    \begin{tabular}{p{0.30\textwidth} p{0.55\textwidth}}
        \toprule
        \textbf{Thực thể} & \textbf{Ràng buộc} \\
        \midrule
        \texttt{Project} & \texttt{key} duy nhất toàn hệ thống. \\
        \texttt{Environment} & \texttt{projectId, key} duy nhất trong project. \\
        \texttt{FeatureFlag} & \texttt{projectId, key} duy nhất trong project. \\
        \texttt{FlagEnvironmentConfig} & \texttt{featureFlagId, environmentId} duy nhất. \\
        \texttt{FlagGroupConfig} & \texttt{groupId, environmentId} duy nhất. \\
        \bottomrule
    \end{tabular}
\end{table}

\subsubsection{Ràng buộc quan hệ tham chiếu}

Khóa ngoại bảo đảm flag, config, rule, group và audit entry không trỏ đến project hoặc environment không tồn tại.

\subsubsection{Ràng buộc giá trị phần trăm rollout}

Percentage rollout phải nằm trong miền hợp lệ, thường từ 0 đến 100.

\subsubsection{Ràng buộc không sửa hoặc xóa audit log}

Audit log không có API update/delete; triển khai production có thể bổ sung trigger database để cưỡng chế thêm.

\subsubsection{Chiến lược archive thay cho xóa vật lý}

Flag không còn dùng được archive thay vì xóa vật lý để giữ lịch sử và tránh phá vỡ audit/history.

\section{Thiết kế Rule Evaluation Engine}

Rule Evaluation Engine là thành phần quyết định một Feature Flag bật hay tắt cho từng request. Mục này là nơi duy nhất trình bày chi tiết thứ tự đánh giá, cách dùng context, stable hashing, reason code và fallback; các mục dữ liệu/API chỉ tham chiếu đến contract này ở mức tổng quan.

\subsection{Evaluation request và context}

Evaluation request gồm khóa project, environment, flag và một evaluation context do ứng dụng khách cung cấp. Context chỉ là đầu vào cho rollout/targeting, không thay thế cơ chế xác thực hay phân quyền nghiệp vụ.

\begin{lstlisting}[
    caption={Cấu trúc evaluation request},
    label={lst:evaluation-request}
]
{
  "projectKey": "ecommerce-demo",
  "environmentKey": "production",
  "flagKey": "new-checkout",
  "context": {
    "targetingKey": "user-123",
    "userId": "user-123",
    "role": "BETA_CUSTOMER",
    "attributes": {
      "organizationId": "org-01"
    }
  }
}
\end{lstlisting}

\subsubsection{Project key, environment key và flag key}

Ba khóa này xác định chính xác cấu hình cần evaluate. Nếu project hoặc flag không tồn tại, engine trả kết quả tắt an toàn với reason \texttt{NOT\_FOUND}.

\subsubsection{Targeting key}

\texttt{targetingKey} là định danh ổn định, không chứa PII, dùng cho percentage rollout. Khi thiếu trường này, rule percentage không đủ điều kiện để match.

\subsubsection{User ID, role và organization ID}

\texttt{userId} phục vụ allowlist; \texttt{role} phục vụ role targeting; \texttt{attributes.organizationId} được giữ như dữ liệu mở rộng cho demo và hướng phát triển sau này.

\subsubsection{Thuộc tính mở rộng}

\texttt{attributes} cho phép bổ sung dữ liệu mà không đổi contract chính. Engine hiện chỉ dùng một phần nhỏ, các rule theo thuộc tính nâng cao được để ở hướng phát triển.

\subsubsection{Validation evaluation request}

Backend kiểm tra kiểu dữ liệu và định dạng bắt buộc trước khi gọi engine. Request sai định dạng bị từ chối ở biên API, còn request hợp lệ nhưng thiếu dữ liệu cho một rule sẽ được engine xử lý theo hướng bỏ qua rule đó.

\subsection{Evaluation snapshot}

Snapshot là ảnh chụp cấu hình cần thiết tại thời điểm evaluate, gồm trạng thái flag, config theo environment, group config và danh sách rule đã chuẩn hóa.

\subsubsection{Khái niệm evaluation snapshot}

Snapshot giúp engine không phụ thuộc trực tiếp vào ORM hoặc controller. Cùng một snapshot và cùng một context phải luôn tạo ra cùng một kết quả.

\subsubsection{Dữ liệu cấu hình cần truy xuất}

Snapshot chứa lifecycle flag, trạng thái enabled/disabled, serving mode, kill switch, group kill switch và rule parameters. Dữ liệu audit/history không nằm trong snapshot runtime.

\subsubsection{Sắp xếp rule trước khi evaluation}

Repository chuẩn bị rule theo loại và priority để engine chỉ tập trung vào logic quyết định. Cách sắp xếp này làm kết quả ổn định giữa các lần gọi.

\subsubsection{Tính nhất quán của snapshot}

Snapshot được lấy từ database hoặc cache đã invalidation sau mutation. Nếu cache miss, backend nạp lại snapshot từ nguồn dữ liệu chính.

\subsection{Thứ tự ưu tiên khi đánh giá}

Thứ tự ưu tiên là hợp đồng quan trọng nhất của engine. Các điều kiện an toàn được xét trước rule rollout để tránh việc một flag bị tắt nhưng vẫn vô tình được bật bởi rule phía sau.

\begin{lstlisting}[
    caption={Thứ tự ưu tiên của Rule Evaluation Engine},
    label={lst:evaluation-precedence}
]
ARCHIVED_FLAG
FLAG_DISABLED
GROUP_KILL_SWITCH
KILL_SWITCH
GLOBAL_ON
USER_ALLOWLIST
ROLE_TARGETING
PERCENTAGE_ROLLOUT
DEFAULT_OFF
\end{lstlisting}

\begin{figure}[H]
    \centering
    \includegraphics[width=0.90\textwidth]{\detokenize{Images/mermaid/Thứ tự ưu tiên khi đánh giá.png}}
    \caption{Flowchart thứ tự ưu tiên của Rule Evaluation Engine}
    \label{fig:evaluation-precedence-flowchart}
\end{figure}

\subsubsection{Flag đã được archive}


Flag archived không tham gia phục vụ runtime, vì vậy engine trả tắt trước mọi cấu hình khác.

\subsubsection{Flag bị disabled}

Nếu config theo environment đang disabled, engine trả tắt ngay cả khi còn rule đang enabled.

\subsubsection{Group kill switch}

Khi group kill switch bật cho environment hiện tại, toàn bộ flag thuộc group bị tắt sau khi đã qua kiểm tra archived/disabled.

\subsubsection{Flag kill switch}

Flag-level kill switch là lớp dừng khẩn cấp cho một flag riêng lẻ.

\subsubsection{Global on}

Nếu config enabled và serving mode là global on, engine bật flag cho mọi context hợp lệ mà không cần xét rule targeting.

\subsubsection{User allowlist}

Allowlist được xét trước các rule mở rộng phạm vi vì đây là điều kiện cụ thể nhất.

\subsubsection{Role targeting}

Role targeting được xét sau allowlist để hỗ trợ phát hành theo nhóm người dùng.

\subsubsection{Percentage rollout}

Percentage rollout là lớp cuối trong nhóm rule vì nó mở rộng phát hành theo tỷ lệ.

\subsubsection{Default off}

Nếu không có điều kiện nào match, engine trả tắt an toàn.

\subsection{Đánh giá trạng thái toàn cục}

Trạng thái lifecycle, enabled/disabled, serving mode và kill switch quyết định các nhánh đầu của engine. Bảng~\ref{tab:status-serving-mode-interaction} tóm tắt các trường hợp chính.

\begin{table}[H]
    \centering
    \caption{Tương tác giữa trạng thái và serving mode}
    \label{tab:status-serving-mode-interaction}
    \begin{tabular}{p{0.22\textwidth} p{0.22\textwidth} p{0.22\textwidth} p{0.24\textwidth}}
        \toprule
        \textbf{Lifecycle/config} & \textbf{Kill switch} & \textbf{Serving mode} & \textbf{Kết quả ưu tiên} \\
        \midrule
        Archived & Bất kỳ & Bất kỳ & Tắt, \texttt{FLAG\_ARCHIVED} \\
        Disabled & Bất kỳ & Bất kỳ & Tắt, \texttt{FLAG\_DISABLED} \\
        Enabled & Group on & Bất kỳ & Tắt, \texttt{GROUP\_KILL\_SWITCH} \\
        Enabled & Flag on & Bất kỳ & Tắt, \texttt{FLAG\_KILL\_SWITCH} \\
        Enabled & Off & Global on & Bật, \texttt{GLOBAL\_ON} \\
        Enabled & Off & Targeted & Xét rule theo precedence \\
        \bottomrule
    \end{tabular}
\end{table}

\subsection{Targeting theo người dùng}

User targeting dựa trên \texttt{context.userId}. Rule chỉ match khi user ID trong context thuộc danh sách đã cấu hình.

\subsubsection{User allowlist}

Allowlist phù hợp cho kiểm thử nội bộ hoặc mở tính năng cho một nhóm nhỏ đã biết trước.

\subsubsection{So khớp user ID}

Việc so khớp dùng giá trị chính xác, không suy diễn từ targeting key để tránh nhầm giữa định danh rollout và định danh nghiệp vụ.

\subsubsection{Tách userId và targetingKey}

\texttt{userId} dùng cho allowlist; \texttt{targetingKey} dùng cho hashing. Tách hai trường giúp hệ thống dùng một khóa rollout không chứa PII.

\subsubsection{Xử lý khi thiếu định danh người dùng}

Nếu thiếu \texttt{userId}, allowlist không match nhưng engine vẫn có thể xét role hoặc percentage nếu context còn đủ dữ liệu cho các rule đó.

\subsection{Targeting theo vai trò}

Role targeting dựa trên vai trò do context cung cấp và phù hợp với demo phân nhóm người dùng.

\subsubsection{Cấu trúc role trong evaluation context}

Role có thể là một chuỗi hoặc danh sách chuỗi tùy contract triển khai; backend chuẩn hóa trước khi so khớp.

\subsubsection{So khớp một hoặc nhiều role}

Rule match khi ít nhất một role của context nằm trong danh sách role được cấu hình.

\subsubsection{Quan hệ giữa role targeting và user targeting}

Allowlist có precedence cao hơn role targeting, vì một người dùng được chỉ định trực tiếp nên được quyết định trước nhóm vai trò rộng hơn.

\subsubsection{Xử lý khi context không chứa role}

Nếu thiếu role, các rule role targeting bị bỏ qua và engine tiếp tục xét lớp rule tiếp theo.

\subsection{Rollout theo phần trăm người dùng}

Percentage rollout dùng stable hashing để phân bổ một targeting key vào bucket cố định từ 0 đến 99.

\subsubsection{Mục tiêu của percentage rollout}

Mục tiêu là phát hành dần tính năng cho một tỷ lệ người dùng ổn định, tránh thay đổi kết quả ngẫu nhiên giữa các request.

\subsubsection{Stable hashing}

Hệ thống băm tổ hợp project/environment/flag/targeting key để cùng một người dùng giữ kết quả ổn định trong phạm vi flag và environment.

\begin{lstlisting}[
    caption={Tạo giá trị băm ổn định bằng SHA-256},
    label={lst:stable-hashing}
]
const hashInput = `${projectKey}:${environmentKey}:${flagKey}:${targetingKey}`;
const digest = sha256(hashInput);
\end{lstlisting}

\subsubsection{Thành phần đầu vào của hàm băm}

Đầu vào băm bao gồm các khóa cấu hình và targeting key. Không sử dụng thời gian hiện tại hoặc số ngẫu nhiên để đảm bảo tính xác định.

\subsubsection{Chuyển đổi kết quả băm thành bucket}

Một phần digest được chuyển thành số nguyên rồi lấy modulo 100 để tạo bucket phần trăm.

\begin{lstlisting}[
    caption={Chuyển đổi SHA-256 digest thành bucket phần trăm},
    label={lst:hash-to-percentage-bucket}
]
const bucket = parseInt(digest.slice(0, 8), 16) % 100;
const enabled = bucket < percentage;
\end{lstlisting}

\subsubsection{Persistent assignment}

Nếu input không đổi, bucket không đổi; vì vậy người dùng không bị nhảy qua lại giữa enabled và disabled.

\subsubsection{Ảnh hưởng khi thay đổi targeting key}

Thay đổi targeting key làm thay đổi bucket, do đó ứng dụng cần chọn khóa ổn định ngay từ đầu.

\subsection{Thực thi rule theo độ ưu tiên}

Rule chỉ được thực thi sau khi các điều kiện hệ thống đã cho phép flag tham gia evaluation.

\subsubsection{Sắp xếp rule}

Trong từng loại rule, các rule enabled được xét theo priority tăng dần.

\subsubsection{Type precedence và first match}

Engine áp dụng type precedence cố định: allowlist, role targeting, percentage rollout. Khi một rule match, engine trả kết quả ngay.

\subsubsection{Rule không hợp lệ hoặc không đủ context}

Rule sai tham số hoặc thiếu dữ liệu context bị bỏ qua thay vì làm bật tính năng ngoài ý muốn.

\subsection{Variant, EvaluationResult và reason code}

Kết quả evaluation phải vừa đủ cho ứng dụng dùng, Dashboard/debug hiểu được, nhưng không phơi bày toàn bộ cấu hình quản trị.

\subsubsection{Variant và response shape}

\begin{lstlisting}[
    caption={Cấu trúc EvaluationResult},
    label={lst:evaluation-result-structure}
]
{
  "enabled": true,
  "variant": "on",
  "reason": "USER_ALLOWLIST",
  "projectKey": "ecommerce-demo",
  "flagKey": "new-checkout"
}
\end{lstlisting}

\subsubsection{Tập reason code}

\begin{lstlisting}[
    caption={Tập reason code của Evaluation API},
    label={lst:evaluation-reason-codes}
]
NOT_FOUND
FLAG_ARCHIVED
FLAG_DISABLED
GROUP_KILL_SWITCH
KILL_SWITCH
GLOBAL_ON
USER_ALLOWLIST
ROLE_MATCH
PERCENTAGE_ROLLOUT
DEFAULT_OFF
INVALID_CONTEXT
ERROR
\end{lstlisting}

\begin{table}[H]
    \centering
    \caption{Reason code cho trạng thái hệ thống}
    \label{tab:system-reason-codes}
    \begin{tabular}{p{0.36\textwidth} p{0.54\textwidth}}
        \toprule
        \textbf{Reason} & \textbf{Ý nghĩa} \\
        \midrule
        \texttt{NOT\_FOUND} & Không tìm thấy project hoặc flag. \\
        \texttt{FLAG\_ARCHIVED} & Flag không còn phục vụ runtime. \\
        \texttt{FLAG\_DISABLED} & Config theo environment đang disabled. \\
        \texttt{GROUP\_KILL\_SWITCH} & Group kill switch đang bật. \\
        \texttt{FLAG\_KILL\_SWITCH} & Kill switch riêng của flag đang bật. \\
        \texttt{ERROR} & Lỗi không mong muốn, trả fail-closed. \\
        \bottomrule
    \end{tabular}
\end{table}

\begin{table}[H]
    \centering
    \caption{Quan hệ giữa mã lý do và matched rule}
    \label{tab:reason-matched-rule}
    \begin{tabular}{p{0.36\textwidth} p{0.54\textwidth}}
        \toprule
        \textbf{Reason} & \textbf{Matched rule} \\
        \midrule
        \texttt{USER\_TARGET\_MATCH} & Rule allowlist match \texttt{userId}. \\
        \texttt{ROLE\_TARGET\_MATCH} & Rule role targeting match role của context. \\
        \texttt{PERCENTAGE\_ROLLOUT\_MATCH} & Bucket stable hashing nằm trong tỷ lệ rollout. \\
        \texttt{DEFAULT\_OFF} & Không rule nào match. \\
        \bottomrule
    \end{tabular}
\end{table}

\subsubsection{Metadata phục vụ truy vết}

Response có thể kèm metadata như matched rule ID hoặc cache status khi cần debug, nhưng metadata này không bắt buộc đối với ứng dụng khách.

\subsection{Default off và error fallback}

Nguyên tắc chung là fail-closed: khi thiếu thông tin hoặc gặp lỗi, hệ thống ưu tiên tắt tính năng mới.

\subsubsection{Nguyên tắc fail-safe và fail-closed}

Fail-closed giúp lỗi hạ tầng không trở thành hành vi mở tính năng ngoài kế hoạch.

\begin{lstlisting}[
    caption={Kết quả fallback mặc định},
    label={lst:default-fallback-result}
]
{
  "enabled": false,
  "variant": "off",
  "reason": "ERROR"
}
\end{lstlisting}

\subsubsection{Flag không tồn tại}

Không tìm thấy project hoặc flag trả \texttt{enabled=false} và \texttt{reason=NOT\_FOUND}.

\subsubsection{Context không đầy đủ}

Context thiếu dữ liệu chỉ làm rule tương ứng không match; nếu không còn rule nào match thì kết quả cuối là \texttt{DEFAULT\_OFF} hoặc \texttt{INVALID\_CONTEXT} tùy loại lỗi.

\subsubsection{Lỗi database hoặc cache}

Khi không tải được snapshot, backend không suy đoán kết quả mà trả fallback tắt.

\subsubsection{Lỗi nội bộ của Evaluation Engine}

\begin{lstlisting}[
    caption={Xử lý lỗi nội bộ của Evaluation Engine},
    label={lst:evaluation-engine-error-handling}
]
try {
  return evaluateWithSnapshot(snapshot, context);
} catch (error) {
  return disabled("ERROR");
}
\end{lstlisting}

\subsection{Khả năng mở rộng Rule Evaluation Engine}

Engine được thiết kế theo hướng thêm loại rule mới mà không thay đổi contract response. Các hướng mở rộng phù hợp gồm targeting theo thuộc tính, organization, thời gian, vị trí và phân phối nhiều variant cho A/B testing.

\section{Thiết kế API}

API của hệ thống được thiết kế theo REST, dùng JSON, đặt dưới tiền tố \texttt{/v1} và dùng quy ước \texttt{camelCase}. Backend NestJS sinh Swagger/OpenAPI từ controller và DTO; khi chạy, Swagger UI có tại \texttt{/docs} và đặc tả JSON có tại \texttt{/docs/json}. Thiết kế API tách \textit{control plane} cho quản trị cấu hình, \textit{data plane} cho runtime evaluation và \textit{monitoring} cho audit/statistics; các endpoint dùng khóa nghiệp vụ ổn định như \texttt{projectKey}, \texttt{flagKey}, \texttt{groupKey} và request ID \texttt{X-Request-Id} để hỗ trợ truy vết.

\subsection{Control plane API}

Control plane quản lý project, flag, rule, group và sample user. Nhóm endpoint này yêu cầu bearer token demo \texttt{demoBearer}; backend resolve actor theo token, kiểm tra quyền theo vai trò \texttt{ADMIN}/\texttt{DEVELOPER}/\texttt{VIEWER}, chuẩn hóa pagination bằng \texttt{items}/\texttt{page}, và ghi audit log cho các mutation quan trọng. Environment chưa có CRUD riêng; dữ liệu environment được dùng qua seed/model và được tham chiếu bằng \texttt{environmentKey} trong evaluation, group config và statistics.

\subsubsection{Project API}

\begin{table}[H]
    \centering
    \scriptsize
    \caption{Đặc tả Project API}
    \label{tab:api-projects}
    \begin{tabular}{@{}p{0.12\textwidth} p{0.34\textwidth} p{0.19\textwidth} p{0.26\textwidth}@{}}
        \toprule
        \textbf{Method} & \textbf{Endpoint} & \textbf{Quyền} & \textbf{Mục đích và dữ liệu chính} \\
        \midrule
        GET & \path|/v1/projects| & \texttt{CONTROL\_PLANE\_READ} & Liệt kê project, trả \texttt{items} và \texttt{page}. Query kế thừa pagination. \\
        POST & \path|/v1/projects| & \texttt{PROJECT\_MANAGE} & Tạo project với \texttt{key}, \texttt{name}, \texttt{description}. Trả \texttt{201 Created}. \\
        GET & \path|/v1/projects/{projectKey}| & \texttt{CONTROL\_PLANE\_READ} & Xem chi tiết project theo khóa nghiệp vụ. \\
        PATCH & \path|/v1/projects/{projectKey}| & \texttt{PROJECT\_MANAGE} & Cập nhật \texttt{name} hoặc \texttt{description}; \texttt{key} là bất biến. \\
        DELETE & \path|/v1/projects/{projectKey}| & \texttt{PROJECT\_MANAGE} & Xóa project ở mức quản trị, trả \texttt{204 No Content}. \\
        \bottomrule
    \end{tabular}
\end{table}

\subsubsection{Feature Flag API}

\begin{table}[H]
    \centering
    \scriptsize
    \caption{Đặc tả Feature Flag API}
    \label{tab:api-feature-flags}
    \begin{tabular}{@{}p{0.12\textwidth} p{0.38\textwidth} p{0.17\textwidth} p{0.24\textwidth}@{}}
        \toprule
        \textbf{Method} & \textbf{Endpoint} & \textbf{Quyền} & \textbf{Mục đích và dữ liệu chính} \\
        \midrule
        GET & \path|/v1/projects/{projectKey}/flags| & \texttt{CONTROL\_PLANE\_READ} & Liệt kê flag chưa soft-delete. Query hỗ trợ \texttt{search}, \texttt{status}, \texttt{lifecycleStatus} và pagination. \\
        GET & \path|/v1/projects/{projectKey}/flags/deleted| & \texttt{CONTROL\_PLANE\_READ} & Liệt kê flag đã soft-delete để có thể khôi phục khi cần. \\
        POST & \path|/v1/projects/{projectKey}/flags| & \texttt{FLAG\_MANAGE} & Tạo flag với \texttt{key}, \texttt{name}, \texttt{description}; backend tạo config mặc định cho environment mặc định. \\
        GET & \path|/v1/projects/{projectKey}/flags/{flagKey}| & \texttt{CONTROL\_PLANE\_READ} & Xem chi tiết một flag, gồm group summary nếu flag đang được gán group. \\
        PATCH & \path|/v1/projects/{projectKey}/flags/{flagKey}| & \texttt{FLAG\_MANAGE} & Cập nhật metadata và config runtime: \texttt{name}, \texttt{description}, \texttt{status}, \texttt{servingMode}, \texttt{killSwitch}. \\
        POST & \path|/v1/projects/{projectKey}/flags/{flagKey}/archive| & \texttt{FLAG\_LIFECYCLE\_MANAGE} & Archive flag; data plane sẽ trả tắt với reason \texttt{FLAG\_ARCHIVED}. \\
        POST & \path|/v1/projects/{projectKey}/flags/{flagKey}/restore| & \texttt{FLAG\_LIFECYCLE\_MANAGE} & Khôi phục flag đã archive về vòng đời active. \\
        DELETE & \path|/v1/projects/{projectKey}/flags/{flagKey}| & \texttt{FLAG\_LIFECYCLE\_MANAGE} & Soft-delete flag, trả \texttt{204 No Content}. \\
        POST & \path|/v1/projects/{projectKey}/flags/{flagKey}/restore-deleted| & \texttt{FLAG\_LIFECYCLE\_MANAGE} & Khôi phục flag đã soft-delete. \\
        \bottomrule
    \end{tabular}
\end{table}

\subsubsection{Rule API}

\begin{table}[H]
    \centering
    \scriptsize
    \caption{Đặc tả Rule API}
    \label{tab:api-rules}
    \begin{tabular}{@{}p{0.12\textwidth} p{0.38\textwidth} p{0.18\textwidth} p{0.23\textwidth}@{}}
        \toprule
        \textbf{Method} & \textbf{Endpoint} & \textbf{Quyền} & \textbf{Mục đích và dữ liệu chính} \\
        \midrule
        GET & \path|/v1/projects/{projectKey}/flags/{flagKey}/rules| & \texttt{CONTROL\_PLANE\_READ} & Liệt kê rule của flag. Query hỗ trợ \texttt{type} và pagination. \\
        PUT & \path|/v1/projects/{projectKey}/flags/{flagKey}/rules| & \texttt{RULE\_MANAGE} & Thay thế toàn bộ rule set bằng mảng \texttt{rules}. Mỗi rule gồm \texttt{type}, \texttt{priority}, \texttt{enabled}, \texttt{parameters}. \\
        \bottomrule
    \end{tabular}
\end{table}

\subsubsection{Flag Group API}

\begin{table}[H]
    \centering
    \scriptsize
    \caption{Đặc tả Flag Group API}
    \label{tab:api-flag-groups}
    \begin{tabular}{@{}p{0.12\textwidth} p{0.38\textwidth} p{0.18\textwidth} p{0.23\textwidth}@{}}
        \toprule
        \textbf{Method} & \textbf{Endpoint} & \textbf{Quyền} & \textbf{Mục đích và dữ liệu chính} \\
        \midrule
        GET & \path|/v1/projects/{projectKey}/groups| & \texttt{CONTROL\_PLANE\_READ} & Liệt kê group. Query hỗ trợ \texttt{search}, \texttt{environmentKey} và pagination. \\
        POST & \path|/v1/projects/{projectKey}/groups| & \texttt{GROUP\_MANAGE} & Tạo group với \texttt{key} và \texttt{name}. \texttt{key} là bất biến trong project. \\
        PATCH & \path|/v1/projects/{projectKey}/groups/{groupKey}| & \texttt{GROUP\_MANAGE} & Đổi \texttt{name} của group. \\
        DELETE & \path|/v1/projects/{projectKey}/groups/{groupKey}| & \texttt{GROUP\_MANAGE} & Xóa group, trả \texttt{204 No Content}. \\
        PUT & \path|/v1/projects/{projectKey}/groups/{groupKey}/config| & \texttt{GROUP\_KILL\_SWITCH} & Cập nhật \texttt{killSwitch} cho một \texttt{environmentKey}. \\
        PUT & \path|/v1/projects/{projectKey}/flags/{flagKey}/group| & \texttt{GROUP\_ASSIGN} & Gán flag vào group bằng body \texttt{groupKey}; trả lại flag response mới. \\
        DELETE & \path|/v1/projects/{projectKey}/flags/{flagKey}/group| & \texttt{GROUP\_ASSIGN} & Gỡ flag khỏi group; trả lại flag response mới. \\
        \bottomrule
    \end{tabular}
\end{table}

\subsubsection{Sample User API}

\begin{table}[H]
    \centering
    \scriptsize
    \caption{Đặc tả Sample User API}
    \label{tab:api-sample-users}
    \begin{tabular}{@{}p{0.12\textwidth} p{0.36\textwidth} p{0.20\textwidth} p{0.23\textwidth}@{}}
        \toprule
        \textbf{Method} & \textbf{Endpoint} & \textbf{Quyền} & \textbf{Mục đích và dữ liệu chính} \\
        \midrule
        GET & \path|/v1/projects/{projectKey}/sample-users| & \texttt{CONTROL\_PLANE\_READ} & Liệt kê sample user. Query hỗ trợ \texttt{search}, \texttt{role} và pagination. \\
        POST & \path|/v1/projects/{projectKey}/sample-users| & \texttt{SAMPLE\_USER\_MANAGE} & Tạo sample user với \texttt{displayName}, \texttt{targetingKey}, \texttt{userId}, \texttt{roles}, \texttt{attributes}. \\
        DELETE & \path|/v1/projects/{projectKey}/sample-users/{targetingKey}| & \texttt{SAMPLE\_USER\_MANAGE} & Xóa sample user theo \texttt{targetingKey}, trả \texttt{204 No Content}. \\
        \bottomrule
    \end{tabular}
\end{table}

\subsection{Data plane API}

Data plane chỉ có một endpoint chính: \texttt{POST /v1/evaluate}. Endpoint này được đánh dấu public trong controller, không dùng bearer token quản trị, và được tối ưu cho tính xác định, fail-closed và an toàn khi thiếu cấu hình. Đây là API mà Demo Application gọi thông qua \texttt{@ffp/js-sdk} để quyết định có hiển thị một tính năng hay không.

\begin{table}[H]
    \centering
    \scriptsize
    \caption{Đặc tả Evaluation API}
    \label{tab:api-evaluation}
    \begin{tabular}{@{}p{0.12\textwidth} p{0.24\textwidth} p{0.20\textwidth} p{0.35\textwidth}@{}}
        \toprule
        \textbf{Method} & \textbf{Endpoint} & \textbf{Quyền} & \textbf{Mục đích và dữ liệu chính} \\
        \midrule
        POST & \path|/v1/evaluate| & Public & Trả quyết định runtime cho một \texttt{projectKey}, \texttt{flagKey}, \texttt{environmentKey} tùy chọn và \texttt{context}. Response luôn có dạng evaluation result với \texttt{enabled}, \texttt{variant}, \texttt{reason}. \\
        \bottomrule
    \end{tabular}
\end{table}

Request body của Evaluation API gồm khóa cấu hình và context runtime:

\begin{lstlisting}[
    caption={Ví dụ request Evaluation API},
    label={lst:evaluation-api-request}
]
{
  "projectKey": "demo-project",
  "environmentKey": "production",
  "flagKey": "new-checkout",
  "context": {
    "targetingKey": "demo-user-beta",
    "userId": "demo-user-beta",
    "roles": ["beta-tester"],
    "attributes": {
      "plan": "pro",
      "country": "VN"
    }
  }
}
\end{lstlisting}

Response body trả lại quyết định đã được giải thích bằng reason code:

\begin{lstlisting}[
    caption={Ví dụ response Evaluation API},
    label={lst:evaluation-api-response}
]
{
  "projectKey": "demo-project",
  "flagKey": "new-checkout",
  "enabled": true,
  "variant": "on",
  "reason": "ROLE_MATCH",
  "matchedRuleId": "rule_123"
}
\end{lstlisting}

Nếu project, environment, flag hoặc config theo environment không tồn tại, API không trả lỗi \texttt{404} kiểu control plane mà trả kết quả an toàn: \texttt{enabled = false}, \texttt{variant = "off"}, \texttt{reason = "NOT\_FOUND"}. Nếu request hợp lệ nhưng thiếu context cần thiết cho percentage rollout, engine trả \texttt{INVALID\_CONTEXT}. Nếu có lỗi bất ngờ trong đường evaluate, service bắt lỗi và trả \texttt{ERROR} theo nguyên tắc fail-closed. Các reason code hiện có gồm \texttt{GLOBAL\_ON}, \texttt{FLAG\_DISABLED}, \texttt{FLAG\_ARCHIVED}, \texttt{GROUP\_KILL\_SWITCH}, \texttt{KILL\_SWITCH}, \texttt{USER\_ALLOWLIST}, \texttt{ROLE\_MATCH}, \texttt{PERCENTAGE\_ROLLOUT}, \texttt{DEFAULT\_OFF}, \texttt{NOT\_FOUND}, \texttt{INVALID\_CONTEXT} và \texttt{ERROR}.

Thứ tự ưu tiên của data plane là: archived flag, disabled config, group kill switch, flag kill switch, global on, ordered enabled rules và default off. Vì vậy một rule rollout không thể vô tình bật flag khi flag đã bị archive, disabled hoặc bị kill switch.

\subsection{Monitoring API}

Monitoring API gồm audit/history và evaluation statistics. Audit log là nguồn append-only để truy vết actor, action, target, request ID và snapshot trước/sau; flag history là projection từ audit log, không tạo nguồn sự thật thứ hai. Statistics chỉ lưu dữ liệu aggregate theo flag, environment, reason, enabled/disabled và bucket giờ UTC, không lưu raw evaluation context hay định danh người dùng cuối. Ngoài ra backend có health probe công khai \path|GET /v1/health|.

\subsubsection{Audit Log API}

\begin{table}[H]
    \centering
    \scriptsize
    \caption{Đặc tả Audit Log và Flag History API}
    \label{tab:api-audit-history}
    \begin{tabular}{@{}p{0.12\textwidth} p{0.38\textwidth} p{0.19\textwidth} p{0.22\textwidth}@{}}
        \toprule
        \textbf{Method} & \textbf{Endpoint} & \textbf{Quyền} & \textbf{Mục đích và dữ liệu chính} \\
        \midrule
        GET & \path|/v1/projects/{projectKey}/audit-logs| & \texttt{CONTROL\_PLANE\_READ} & Liệt kê audit log của project. Query hỗ trợ \texttt{targetType}, \texttt{targetKey}, \texttt{actor}, \texttt{action}, \texttt{from}, \texttt{to} và pagination. \\
        GET & \path|/v1/projects/{projectKey}/flags/{flagKey}/history| & \texttt{CONTROL\_PLANE\_READ} & Lịch sử cấu hình của một flag, được projection từ audit log liên quan đến flag, config và rule. Query hỗ trợ pagination và \texttt{sort=createdAt}. \\
        \bottomrule
    \end{tabular}
\end{table}

\subsubsection{Evaluation Statistics API}

\begin{table}[H]
    \centering
    \scriptsize
    \caption{Đặc tả Evaluation Statistics API}
    \label{tab:api-statistics}
    \begin{tabular}{@{}p{0.12\textwidth} p{0.36\textwidth} p{0.19\textwidth} p{0.24\textwidth}@{}}
        \toprule
        \textbf{Method} & \textbf{Endpoint} & \textbf{Quyền} & \textbf{Mục đích và dữ liệu chính} \\
        \midrule
        GET & \path|/v1/projects/{projectKey}/stats/flags| & \texttt{CONTROL\_PLANE\_READ} & Thống kê tổng hợp theo các flag trong project. Query hỗ trợ \texttt{environmentKey}, \texttt{from}, \texttt{to}, pagination và sort theo \texttt{totalEvaluations}, \texttt{enabledCount}, \texttt{disabledCount}, \texttt{flagKey}. \\
        GET & \path|/v1/projects/{projectKey}/flags/{flagKey}/stats| & \texttt{CONTROL\_PLANE\_READ} & Thống kê chi tiết một flag: tổng lượt evaluate, số enabled/disabled, phần trăm enabled, breakdown theo reason và bucket theo giờ. Query hỗ trợ \texttt{environmentKey}, \texttt{from}, \texttt{to}. \\
        \bottomrule
    \end{tabular}
\end{table}

\subsection{Validation và mô hình lỗi}

Validation được thực hiện ở biên API bằng DTO, \texttt{class-validator}, \texttt{class-transformer} và global \texttt{ValidationPipe}. Pipe bật \texttt{whitelist}, \texttt{forbidNonWhitelisted} và \texttt{transform}, vì vậy field ngoài DTO bị từ chối, query string như \texttt{limit} được chuyển sang số, và lỗi validation được chuẩn hóa trước khi vào service. Các khóa nghiệp vụ như \texttt{projectKey}, \texttt{flagKey}, \texttt{groupKey} và \texttt{environmentKey} dùng regex:

\begin{lstlisting}[
    caption={Định dạng key nghiệp vụ},
    label={lst:key-regex}
]
^[a-z0-9][a-z0-9-]{1,62}[a-z0-9]$
\end{lstlisting}

Quy tắc này tương ứng độ dài 3--64 ký tự, chỉ gồm chữ thường, chữ số và dấu gạch ngang, không bắt đầu hoặc kết thúc bằng dấu gạch ngang. Các DTO còn giới hạn độ dài \texttt{name} tối đa 120 ký tự, \texttt{description} tối đa 1000 ký tự, pagination \texttt{limit} từ 1 đến 100, \texttt{offset} không âm, thời gian \texttt{from}/\texttt{to} theo ISO-8601 và enum phải thuộc giá trị hợp lệ.

Control-plane error response có dạng thống nhất:

\begin{lstlisting}[
    caption={Mẫu error response của control plane},
    label={lst:api-error-response}
]
{
  "code": "VALIDATION_ERROR",
  "message": "Request validation failed.",
  "details": [
    {
      "field": "projectKey",
      "message": "projectKey must be a valid 3-64 character key."
    }
  ],
  "requestId": "req_123"
}
\end{lstlisting}

Các mã lỗi chuẩn gồm \texttt{VALIDATION\_ERROR}, \texttt{UNAUTHORIZED}, \texttt{FORBIDDEN}, \texttt{NOT\_FOUND}, \texttt{CONFLICT} và \texttt{INTERNAL\_ERROR}. \texttt{UNAUTHORIZED} dùng khi thiếu hoặc sai bearer token demo; \texttt{FORBIDDEN} dùng khi role đã xác thực nhưng không có permission; \texttt{CONFLICT} dùng cho xung đột như khóa đã tồn tại; \texttt{INTERNAL\_ERROR} che giấu chi tiết lỗi nội bộ nhưng vẫn log kèm request ID. Riêng Evaluation API có thêm mô hình lỗi domain bằng reason code trong response 200 để giữ fail-closed, ví dụ \texttt{NOT\_FOUND}, \texttt{INVALID\_CONTEXT} hoặc \texttt{ERROR} đều trả \texttt{enabled = false} và \texttt{variant = off} thay vì để client tự suy diễn.

\section{Hiệu năng, nhất quán và bảo mật}

Mục này tổng hợp các quyết định phi chức năng quan trọng của hệ thống dựa trên hiện trạng triển khai trong source code. Trọng tâm là cách data plane giữ độ trễ thấp nhưng vẫn fail-closed, cách control plane bảo vệ các thao tác thay đổi cấu hình, và cách hệ thống duy trì khả năng truy vết khi cấu hình feature flag thay đổi. Các chi tiết giao diện và kết quả chạy thử được trình bày ở chương kết quả.

\subsection{Hiệu năng data plane: snapshot cache, TTL và thống kê}

Evaluation API là đường chạy thường xuyên nhất vì demo application và JavaScript SDK gọi endpoint này để quyết định bật/tắt tính năng. Để giảm số lần đọc PostgreSQL lặp lại, backend cache \textit{evaluation snapshot} thay vì cache quyết định cuối cùng theo từng người dùng. Snapshot chứa phần cấu hình cần thiết để engine đánh giá một flag trong một environment: trạng thái lifecycle của flag, trạng thái cấu hình theo environment, serving mode, kill switch, group kill switch nếu có, và danh sách rule đã nạp từ database. Khi cache miss, \texttt{EvaluationRepository.findSnapshot} đọc lại cấu hình từ PostgreSQL; khi cache hit, engine vẫn tự tính kết quả dựa trên evaluation context của request.

Cách cache snapshot này phù hợp với yêu cầu an toàn của mini project. Hệ thống không ưu tiên cache final decision theo user context, vì quyết định cuối cùng phụ thuộc vào \texttt{targetingKey}, \texttt{userId}, \texttt{roles} và các thuộc tính đi kèm. Nếu cache trực tiếp kết quả cuối, hệ thống dễ phải đưa định danh người dùng vào cache key hoặc lưu dữ liệu context nhạy cảm. Do đó thiết kế hiện tại đánh đổi một phần hiệu năng để giữ mô hình đơn giản hơn, ít rủi ro dữ liệu cá nhân hơn, và không làm thay đổi tính xác định của engine.

Cache provider được đóng gói sau interface \texttt{EvaluationSnapshotCache}. Source code hiện hỗ trợ ba lựa chọn: \texttt{memory}, \texttt{redis} và \texttt{none}. Mặc định hệ thống dùng in-memory cache với TTL 30 giây thông qua \texttt{DEFAULT\_EVALUATION\_CACHE\_TTL\_MS}; giá trị này có thể cấu hình bằng biến môi trường \texttt{EVALUATION\_CACHE\_TTL\_MS}. In-memory cache phù hợp cho demo, kiểm thử cục bộ và một backend instance. Redis được triển khai như provider mở rộng cho trường hợp nhiều backend instance cần chia sẻ snapshot hoặc cần invalidation phân tán; nếu không muốn cache, provider \texttt{none} cho phép tắt cache mà không đổi API đánh giá.

Ngoài cache, backend còn ghi thống kê evaluation theo hướng không chặn response. Sau khi engine trả kết quả, \texttt{EvaluationMetricsService} ghi tăng bộ đếm theo project, environment, flag, reason, enabled result và bucket thời gian theo giờ UTC. Nếu ghi thống kê thất bại, hệ thống chỉ log cảnh báo và vẫn giữ nguyên response evaluation. Điều này giúp statistics dashboard có dữ liệu quan sát vận hành nhưng không biến thống kê thành điểm lỗi của data plane.

\subsection{Cache key, invalidation và tính nhất quán cấu hình}

Cache key được thiết kế theo dữ liệu cấu hình, không theo dữ liệu người dùng. Hàm tạo cache key ghép bốn thành phần: prefix \texttt{evaluation-snapshot}, \texttt{projectKey}, \texttt{environmentKey} và \texttt{flagKey}. Khi request không truyền environment, backend dùng scope mặc định \texttt{\_\_default\_\_} để phân biệt với các environment cụ thể. Nhờ đó snapshot của staging, production hoặc default environment không bị dùng nhầm lẫn cho nhau. Evaluation context không nằm trong cache key vì hệ thống chỉ cache snapshot cấu hình, không cache quyết định cuối cùng.

Invalidation được kích hoạt khi một mutation có thể làm thay đổi kết quả evaluation. Các thay đổi như cập nhật lifecycle của flag, bật/tắt config theo environment, đổi serving mode, đổi flag kill switch, thay rule, đổi priority rule, assign/unassign flag group, hoặc cập nhật group kill switch đều cần loại bỏ snapshot liên quan. \texttt{EvaluationCacheInvalidator} hỗ trợ invalidation một flag hoặc nhiều flag theo project, flag key và optional environment key. Nguyên tắc triển khai là thà invalidation hơi rộng còn hơn giữ snapshot lỗi thời; ví dụ khi không chứng minh được scope environment, cache của default alias cũng được xóa để tránh stale read.

\begin{figure}[H]
    \centering
    \includegraphics[width=0.95\textwidth]{\detokenize{Images/mermaid/Cache invalidation - Tính nhất quán cấu hình.png}}
    \caption{Luồng cache invalidation và bảo đảm tính nhất quán cấu hình}
    \label{fig:cache-invalidation-consistency}
\end{figure}

TTL đóng vai trò lớp bảo vệ phụ trong trường hợp một lần invalidation bị bỏ sót hoặc provider cache gặp lỗi. TTL không thay thế invalidation sau mutation cấu hình, vì yêu cầu của hệ thống là thay đổi cấu hình phải được phản ánh nhanh trên data plane. Trong phạm vi một backend instance, mutation ghi database xong rồi invalidation cache giúp hệ thống gần với strong consistency sau thao tác quản trị. Khi mở rộng nhiều instance, Redis hoặc cơ chế pub/sub sẽ cần thiết hơn để giảm độ trễ đồng bộ giữa các instance.

Tính nhất quán của cấu hình cũng được bảo vệ bằng transaction và audit log. Các mutation quan trọng ghi thay đổi cấu hình và audit entry trong cùng một transaction, tránh tình trạng cấu hình đã đổi nhưng không có bằng chứng truy vết. Evaluation path fail-closed: nếu không tìm thấy project, environment, flag hoặc config thì trả \texttt{enabled=false} với \texttt{reason=NOT\_FOUND}; nếu quá trình đọc snapshot hoặc đánh giá gặp lỗi bất thường thì trả \texttt{enabled=false} với \texttt{reason=ERROR}. Cache chỉ là tối ưu hiệu năng, không được phép thay đổi thứ tự ưu tiên hoặc kết quả xác định của engine.

\subsection{Trust boundary của evaluation context và SDK}

Evaluation API được đánh dấu public vì ứng dụng tích hợp cần gọi endpoint này ở runtime. Do đó evaluation context do client gửi lên phải luôn được xem là dữ liệu không tin cậy. Engine có thể dùng context để quyết định rollout, nhưng không dùng context thay cho xác thực hoặc phân quyền nghiệp vụ. Trong demo, các trường như \texttt{userId}, \texttt{roles} hoặc attribute mô phỏng người dùng giúp minh họa user allowlist, role targeting và percentage rollout; trong hệ thống production, các trường nhạy cảm nên được xác nhận bởi backend nghiệp vụ hoặc hệ thống định danh đáng tin cậy.

Thiết kế rollout ưu tiên định danh ổn định nhưng không nhạy cảm. Rule percentage rollout dùng \texttt{targetingKey} cùng \texttt{projectKey} và \texttt{flagKey} để tính bucket bằng stable hashing. Cùng một user context hợp lệ sẽ nhận cùng một kết quả cho cùng một flag, trong khi hệ thống không cần lưu trạng thái rollout riêng cho từng người dùng. Nếu thiếu \texttt{targetingKey} khi có percentage rollout rule, engine trả kết quả fail-closed với \texttt{reason=INVALID\_CONTEXT}.

Ranh giới giữa client-side và server-side evaluation cũng cần được hiểu rõ. Client-side evaluation thuận tiện cho việc ẩn/hiện UI, còn server-side evaluation an toàn hơn cho chức năng nhạy cảm. Mini project dùng API tập trung ở backend để giữ logic đánh giá trong một nơi và để demo app cũng như SDK không tự triển khai lại rule engine. Tuy nhiên, feature flag chỉ quyết định phát hành tính năng; nó không thay thế authorization. Các endpoint quan trọng vẫn phải kiểm tra quyền ở backend ngay cả khi UI bị ẩn bởi feature flag.

JavaScript SDK hỗ trợ tích hợp an toàn hơn bằng cách đóng gói các hàm \texttt{evaluate}, \texttt{isEnabled} và \texttt{getVariant}. SDK validate response từ backend, giới hạn thời gian chờ mặc định 1500 ms, và khi gặp lỗi request, timeout, network, HTTP, JSON hoặc schema response thì trả kết quả fail-closed: \texttt{enabled=false}, \texttt{variant=off}, \texttt{reason=ERROR}. Nhờ đó ứng dụng tích hợp không phải tự xử lý quá nhiều trường hợp lỗi và vẫn giữ mặc định an toàn.

\subsection{Bảo vệ Management API bằng RBAC và validation}

Management API thuộc control plane vì các endpoint này thay đổi project, feature flag, rule, group, sample user, kill switch và các cấu hình liên quan. Khác với Evaluation API public, các controller quản trị dùng bearer token demo và decorator \texttt{RequirePermissions} để yêu cầu quyền phù hợp. Backend là nơi quyết định quyền cuối cùng; dashboard chỉ ẩn hoặc vô hiệu hóa thao tác không được phép nhằm cải thiện trải nghiệm người dùng, không phải lớp bảo vệ chính.

RBAC demo phân biệt ba vai trò. \texttt{admin} có toàn quyền trong phạm vi demo, bao gồm quản lý project, flag, lifecycle, rule, group, group kill switch, sample user, audit log và statistics. \texttt{developer} có quyền đọc control plane và thực hiện các thao tác cấu hình thường gặp như quản lý flag, rule và assign group, nhưng không mặc định có quyền cho các thao tác nhạy cảm hơn như quản lý project hoặc bật group kill switch. \texttt{viewer} chỉ có quyền đọc. Actor được backend resolve từ bearer token thông qua \texttt{DemoIdentityService}; actor ghi vào audit log không lấy trực tiếp từ body request, nhờ đó giảm nguy cơ giả mạo lịch sử thao tác.

Validation DTO và error handling là lớp bảo vệ bổ sung cho control plane. Các khóa project, flag, environment và group được validate ở boundary API; các request phân trang, time range, rule parameters và mutation body được chuẩn hóa trước khi đi vào service. Khi quyền không đủ hoặc dữ liệu không hợp lệ, backend trả lỗi có cấu trúc thay vì để lỗi nội bộ rò rỉ ra client. Cách này giúp API dễ demo, dễ kiểm thử và phù hợp với yêu cầu clean APIs, validation và error handling của đề bài.

\subsection{Auditability và khả năng truy vết thay đổi}

Auditability là yêu cầu xuyên suốt của hệ thống. Mỗi mutation quan trọng cần ghi actor, action, target type, target key, request ID, timestamp, metadata và snapshot trước/sau nếu có. Bảng \texttt{AuditLogEntry} được thiết kế append-only ở mức ứng dụng: service chỉ tạo bản ghi mới cho sự kiện thay đổi, không dùng audit log như bảng trạng thái hiện tại. Dữ liệu này hỗ trợ demo rollback reasoning, kiểm thử, phân tích sự cố và mở rộng sang quy trình phê duyệt trong tương lai \cite{owasp_logging}.

Hệ thống cũng dùng audit log như nguồn cho configuration history. Endpoint history của flag truy vấn các audit entry liên quan đến feature flag, flag environment config và flag rule để hiển thị lịch sử thay đổi cấu hình. Cách làm này tránh tạo một hệ thống versioning song song trong phạm vi mini project, đồng thời vẫn đáp ứng nhu cầu xem lại ai đã thay đổi gì, khi nào, trước/sau ra sao và request nào tạo ra thay đổi đó.

Audit log cần được bảo vệ như dữ liệu quản trị vì snapshot cấu hình có thể tiết lộ rule targeting, rollout strategy hoặc thông tin vận hành. Vì vậy quyền đọc audit/history nằm trong control-plane read permission và đi qua RBAC demo. Khi kết hợp audit log với statistics dashboard, người vận hành có thể trả lời hai câu hỏi khác nhau: cấu hình đã thay đổi như thế nào, và sau thay đổi đó runtime evaluation đang trả những reason nào. Đây là nền tảng để trình diễn system thinking của dự án: hiệu năng, tính nhất quán, bảo mật và khả năng quan sát được thiết kế cùng nhau thay vì xử lý rời rạc.

\section{Phương pháp kiểm thử và đánh giá}

Phần này chỉ xác định chiến lược kiểm thử dùng cho mini project. Các kết quả kiểm thử cụ thể được trình bày ở Chương 3, còn logic chi tiết của Rule Evaluation Engine, API, cache, RBAC và audit log đã được mô tả trong các phần thiết kế tương ứng.

\subsection{Phạm vi kiểm thử}

Hệ thống được kiểm thử theo bốn lớp chính:

\begin{itemize}
    \item \textbf{Domain và Rule Evaluation Engine}: kiểm tra precedence, targeting rule, percentage rollout, kill switch, default off và fallback.
    \item \textbf{Backend API}: kiểm tra contract của Management API và Evaluation API, bao gồm validation, lỗi không tìm thấy, xung đột dữ liệu và response shape.
    \item \textbf{Control-plane side effects}: kiểm tra audit log, cache invalidation, thống kê evaluation và phân quyền demo.
    \item \textbf{Workflow trình diễn}: kiểm tra Dashboard, Demo Application, JavaScript SDK và Docker Compose trên các kịch bản có thể trình bày trực tiếp.
\end{itemize}

Cách chia này giúp tránh kiểm thử trùng lặp cùng một hành vi ở nhiều tầng. Unit test tập trung vào quyết định thuần của engine; integration test xác nhận API, persistence và side effect; kịch bản demo xác nhận trải nghiệm end-to-end.

\subsection{Tiêu chí đánh giá Rule Evaluation Engine}

Rule Evaluation Engine được đánh giá theo ba tiêu chí:

\begin{itemize}
    \item \textbf{Đúng precedence}: trạng thái archive, disabled, group kill switch, flag kill switch và global on phải có thứ tự ưu tiên cao hơn các targeting rule.
    \item \textbf{Xác định}: cùng một snapshot cấu hình và cùng một evaluation context phải luôn tạo ra cùng một kết quả, đặc biệt đối với percentage rollout.
    \item \textbf{Fail-closed}: khi flag không tồn tại, context không đủ hoặc engine gặp lỗi, hệ thống không được tự động bật tính năng ngoài ý muốn.
\end{itemize}

Các test không lặp lại mô tả thuật toán; chúng chỉ tạo dữ liệu đầu vào, gọi engine hoặc Evaluation API và so sánh với kết quả kỳ vọng gồm \texttt{enabled}, \texttt{variant}, \texttt{reason} và metadata rule nếu có.

\subsection{Tiêu chí đánh giá API và dữ liệu}

Đối với API, trọng tâm kiểm thử là contract và tính nhất quán của dữ liệu. Các mutation phải validate input, ghi dữ liệu đúng ràng buộc, tạo audit log trong cùng transaction và làm mất hiệu lực cache liên quan. Các endpoint đọc phải trả về cấu trúc response ổn định, hỗ trợ phân trang khi cần và không để lộ dữ liệu nội bộ.

Đối với Evaluation API, response luôn phải có dạng evaluation-shaped kể cả khi tài nguyên không tồn tại hoặc xảy ra lỗi. Điều này giúp SDK và Demo Application xử lý fallback bằng cùng một contract thay vì phải phân nhánh theo nhiều kiểu lỗi HTTP khác nhau.

\subsection{Tiêu chí đánh giá Dashboard và Demo Application}

Dashboard được đánh giá theo khả năng hỗ trợ thao tác quản trị an toàn: hiển thị đúng trạng thái lifecycle/runtime, vô hiệu hóa thao tác không đủ quyền, phản ánh lỗi validation và cho phép theo dõi audit log/statistics. Demo Application được đánh giá theo khả năng chứng minh việc thay đổi cấu hình Feature Flag có thể làm thay đổi hành vi ứng dụng mà không cần deploy lại.

Các kịch bản demo được chọn để bao phủ những giá trị quan trọng nhất của sản phẩm: bật/tắt toàn cục, targeting theo user hoặc role, rollout theo phần trăm, flag kill switch, group kill switch, audit log, statistics và fallback của SDK.

\subsection{Tiêu chí triển khai và trình diễn}

Quy trình triển khai được đánh giá bằng khả năng khởi tạo hệ thống trong môi trường sạch. Docker Compose phải khởi chạy được PostgreSQL, Backend API, Admin Dashboard và Demo Application; migration và seed data phải có thể chạy lại theo cách idempotent; các service phải giao tiếp đúng qua cấu hình URL và biến môi trường.

Tiêu chí cuối cùng là khả năng trình bày: sau khi khởi chạy, người xem có thể đi từ Dashboard sang Demo Application, thay đổi cấu hình flag và quan sát kết quả runtime cùng audit/statistics phát sinh.
 
% ------------------------------------------------------------------------------
% Chapter 3
% ------------------------------------------------------------------------------
\chapter{Kết quả thực hiện và đánh giá} % enter the name of the chapter here

\section{Môi trường và cấu trúc triển khai}

\subsection{Cấu trúc mã nguồn}

Nhằm tách biệt các thành phần của hệ thống nhưng vẫn thuận tiện cho việc phát triển, kiểm thử và triển khai, toàn bộ mã nguồn được tổ chức theo mô hình monorepo. Kiến trúc này cho phép Backend API, Admin Dashboard, Demo Application và JavaScript SDK cùng sử dụng chung quy trình quản lý phiên bản, phụ thuộc và môi trường phát triển. Đồng thời, việc đặt các thành phần trong cùng một kho mã nguồn giúp giảm chi phí đồng bộ API contract, chia sẻ kiểu dữ liệu và triển khai Docker Compose cho toàn bộ hệ thống \cite{docker_compose}.

\subsubsection{Tổ chức monorepo}

Dự án được tổ chức theo mô hình workspace với hai thư mục chính là \texttt{apps} và \texttt{packages}. Thư mục \texttt{apps} chứa các ứng dụng có thể triển khai độc lập, bao gồm Backend API, Admin Dashboard và Demo Application. Thư mục \texttt{packages} chứa JavaScript SDK dùng chung để các ứng dụng khách gọi Evaluation API.

Cấu trúc này giúp giảm trùng lặp mã nguồn, dễ dàng chia sẻ các kiểu dữ liệu giữa các thành phần và hỗ trợ mở rộng thêm ứng dụng mới mà không ảnh hưởng đến kiến trúc hiện có. Việc phát triển theo monorepo cũng phù hợp với quy trình tích hợp và triển khai liên tục khi toàn bộ hệ thống được quản lý trong cùng một repository \cite{microsoft_devops,docker_compose}.

\subsubsection{Cấu trúc Backend API}

Backend API được xây dựng bằng NestJS theo kiến trúc module nhằm tách biệt các miền chức năng của hệ thống. Mỗi module đảm nhiệm một nghiệp vụ riêng như quản lý project, feature flag, flag group, rule evaluation, audit log, authentication, sample user và thống kê evaluation.

Tầng controller chịu trách nhiệm tiếp nhận và xác thực request, tầng service xử lý nghiệp vụ, còn tầng repository thực hiện truy cập cơ sở dữ liệu thông qua Prisma ORM. Các DTO được sử dụng để kiểm tra dữ liệu đầu vào, trong khi middleware, guard, interceptor và exception filter đảm nhiệm các chức năng cắt ngang như xác thực, phân quyền, ghi nhận request ID, logging và xử lý lỗi thống nhất. Kiến trúc phân lớp này giúp tăng khả năng bảo trì, kiểm thử và mở rộng hệ thống khi bổ sung các loại rule hoặc endpoint mới.

\subsubsection{Cấu trúc Admin Dashboard}

Admin Dashboard được xây dựng bằng React theo hướng component-based. Giao diện được chia thành các nhóm thành phần độc lập như trang quản lý project, feature flag, flag group, audit log, statistics và màn hình cấu hình rule.

Các thành phần dùng chung như layout, biểu mẫu, bảng dữ liệu, hộp thoại xác nhận, badge trạng thái và hook xử lý API được tách riêng để tái sử dụng trong toàn bộ ứng dụng. Dữ liệu được lấy từ Backend API thông qua các service riêng biệt, giúp tách logic giao diện khỏi logic giao tiếp mạng và đơn giản hóa việc bảo trì khi API thay đổi.

\subsubsection{Cấu trúc Demo Application và SDK}

Demo Application là một ứng dụng React độc lập dùng để minh họa việc tích hợp Feature Flag vào ứng dụng khách. Thay vì gọi trực tiếp Evaluation API, ứng dụng sử dụng JavaScript SDK được phát triển trong thư mục \texttt{packages/js-sdk}. SDK chịu trách nhiệm xây dựng request, gửi Evaluation API, kiểm tra tính hợp lệ của response và trả về kết quả dưới dạng giao diện lập trình thống nhất cho ứng dụng.

Việc tách SDK thành một package riêng giúp giảm sự phụ thuộc giữa ứng dụng khách và Backend API, đồng thời tạo điều kiện tái sử dụng cho nhiều ứng dụng khác nhau. Cách tổ chức này cũng phù hợp với mô hình Feature Flag Platform trong thực tế, nơi client chỉ cần tích hợp SDK mà không cần quan tâm đến chi tiết triển khai của Evaluation Engine \cite{launchdarkly_feature_flags,openfeature_flag_evaluation}.

\subsection{Môi trường phát triển}

Để bảo đảm tính ổn định, khả năng tái lập môi trường và thuận tiện trong quá trình phát triển, hệ thống được xây dựng trên nền tảng Node.js với kiến trúc monorepo. Backend API, Admin Dashboard, Demo Application và JavaScript SDK cùng sử dụng chung hệ sinh thái công cụ, đồng thời cơ sở dữ liệu được quản lý thông qua Prisma ORM và PostgreSQL. Ngoài ra, dự án tích hợp các công cụ kiểm thử, kiểm tra chất lượng mã nguồn và tài liệu hóa API nhằm nâng cao khả năng bảo trì và mở rộng hệ thống.

\subsubsection{Runtime và package manager}

Toàn bộ hệ thống được phát triển bằng TypeScript trên nền tảng Node.js. Backend sử dụng NestJS, trong khi Admin Dashboard và Demo Application được xây dựng bằng React. Các thành phần trong monorepo được quản lý bằng npm workspaces, cho phép chia sẻ package nội bộ, đồng bộ phiên bản phụ thuộc và đơn giản hóa quá trình cài đặt.

Việc sử dụng cùng một runtime giúp giảm sự khác biệt giữa các thành phần của hệ thống, đồng thời tạo điều kiện thuận lợi cho việc xây dựng quy trình build, kiểm thử và triển khai thống nhất. Đây cũng là mô hình thường được áp dụng trong các hệ thống web hiện đại theo quy trình DevOps \cite{microsoft_devops}.

\subsubsection{Database và công cụ ORM}

Hệ thống sử dụng PostgreSQL làm hệ quản trị cơ sở dữ liệu chính. PostgreSQL đáp ứng tốt các yêu cầu về tính toàn vẹn dữ liệu, hỗ trợ transaction, khóa ngoại, chỉ mục và kiểu dữ liệu JSONB phục vụ lưu trữ linh hoạt các tham số rule, thuộc tính người dùng và snapshot của audit log \cite{postgresql_json_types}.

Lớp truy cập dữ liệu được xây dựng bằng Prisma ORM. Prisma cung cấp cơ chế mô hình hóa dữ liệu thông qua schema, hỗ trợ migration, sinh mã truy cập cơ sở dữ liệu theo kiểu dữ liệu tĩnh và giảm đáng kể lượng mã SQL thủ công. Việc sử dụng Prisma cũng giúp tăng tính an toàn kiểu dữ liệu, đơn giản hóa thao tác CRUD và thuận tiện khi mở rộng mô hình dữ liệu trong tương lai.

\subsubsection{Công cụ kiểm thử và kiểm tra mã nguồn}

Dự án sử dụng Jest để xây dựng các bài kiểm thử đơn vị cho các thành phần nghiệp vụ quan trọng, đặc biệt là Rule Evaluation Engine, cache, SDK và các hàm tiện ích. Các trường hợp kiểm thử bao gồm đánh giá theo user allowlist, role, percentage rollout, group/flag kill switch, default off, \texttt{NOT\_FOUND}, \texttt{INVALID\_CONTEXT} và các tình huống ngoại lệ nhằm bảo đảm tính đúng đắn của thuật toán evaluation.

Bên cạnh đó, ESLint được sử dụng để kiểm tra chất lượng mã nguồn và phát hiện các lỗi phổ biến trong quá trình phát triển. Việc kết hợp TypeScript với ESLint giúp chuẩn hóa phong cách lập trình, hạn chế lỗi tiềm ẩn và nâng cao khả năng bảo trì của toàn bộ hệ thống.

\subsubsection{Công cụ tài liệu hóa API}

Backend API tích hợp OpenAPI thông qua Swagger để tự động sinh tài liệu REST API. Các endpoint được mô tả đầy đủ về request, response, tham số đầu vào, mã lỗi và kiểu dữ liệu, giúp lập trình viên dễ dàng kiểm thử và tích hợp với các ứng dụng khách.

Việc sử dụng Swagger giúp tài liệu API luôn đồng bộ với mã nguồn, giảm công sức cập nhật tài liệu thủ công và hỗ trợ kiểm thử trực tiếp ngay trên giao diện trình duyệt. Cách tiếp cận này phù hợp với các hệ thống Feature Flag hiện đại, nơi Evaluation API đóng vai trò là điểm tích hợp giữa nền tảng quản trị và các ứng dụng khách \cite{openfeature_flag_evaluation}.

\subsection{Database migration và seed data}

Để bảo đảm cơ sở dữ liệu luôn đồng bộ với mã nguồn và có thể tái tạo trên nhiều môi trường khác nhau, hệ thống sử dụng Prisma Migrate để quản lý schema và lịch sử thay đổi cơ sở dữ liệu. Bên cạnh đó, các tập dữ liệu mẫu được chuẩn bị sẵn nhằm hỗ trợ phát triển, kiểm thử và trình diễn các kịch bản Feature Flag mà không cần nhập dữ liệu thủ công.

\subsubsection{Khởi tạo database schema}

Toàn bộ mô hình dữ liệu của hệ thống được định nghĩa trong Prisma Schema và sinh tự động thành các bảng trong PostgreSQL thông qua Prisma Migrate. Schema bao gồm các thực thể chính như Project, Environment, Feature Flag, Flag Group, FlagEnvironmentConfig, FlagRule, SampleUserContext, AuditLogEntry và các bảng phục vụ thống kê evaluation.

Các ràng buộc toàn vẹn như khóa chính, khóa ngoại, chỉ mục, giá trị duy nhất và các quan hệ một-nhiều được khai báo trực tiếp trong schema. Ngoài các kiểu dữ liệu cơ bản, PostgreSQL còn sử dụng kiểu dữ liệu JSONB để lưu trữ linh hoạt các tham số rule, thuộc tính người dùng và snapshot của audit log, giúp hệ thống dễ dàng mở rộng khi bổ sung các loại rule mới \cite{postgresql_json_types}.

\subsubsection{Quản lý migration}

Mỗi thay đổi đối với mô hình dữ liệu đều được quản lý dưới dạng migration độc lập. Prisma tự động sinh các tập lệnh SQL tương ứng và lưu lại lịch sử migration trong cơ sở dữ liệu, bảo đảm mọi môi trường phát triển và triển khai đều sử dụng cùng một phiên bản schema.

Cách tiếp cận này giúp kiểm soát sự thay đổi của cơ sở dữ liệu theo thời gian, hỗ trợ cập nhật hoặc khôi phục schema khi cần thiết, đồng thời giảm rủi ro sai khác giữa môi trường phát triển, kiểm thử và production.

\subsubsection{Seed project, environment và Feature Flag}

Sau khi hoàn tất migration, hệ thống cung cấp các script seed để tạo dữ liệu mẫu phục vụ quá trình phát triển và trình diễn. Dữ liệu khởi tạo bao gồm một hoặc nhiều project, các environment mặc định như \texttt{development}, \texttt{staging} và \texttt{production}, các flag group cùng nhiều feature flag đại diện cho các kịch bản phổ biến.

Mỗi feature flag được cấu hình với trạng thái, serving mode, kill switch và các rule phù hợp để minh họa các chiến lược rollout như bật hoặc tắt toàn cục, targeting theo user, targeting theo role, percentage rollout và group kill switch. Nhờ đó, các chức năng của Evaluation Engine và Dashboard có thể được kiểm thử ngay sau khi khởi tạo dữ liệu.

\subsubsection{Seed sample user và thuộc tính organization mẫu}

Hệ thống đồng thời tạo tập sample user và thuộc tính organization mẫu nhằm phục vụ các kịch bản đánh giá Feature Flag. Mỗi người dùng được gán một \texttt{targetingKey} duy nhất cùng với các thuộc tính như vai trò, organization và các metadata cần thiết cho quá trình evaluation.

Việc chuẩn bị sẵn dữ liệu mẫu giúp dễ dàng kiểm thử các trường hợp targeting theo user ID, role và rollout theo organization-level targetingKey trong demo, đồng thời bảo đảm các kết quả percentage rollout luôn ổn định nhờ sử dụng cùng một định danh người dùng trong quá trình đánh giá. Các tập dữ liệu này cũng được sử dụng để minh họa các kịch bản demo và hỗ trợ kiểm thử tự động của hệ thống.

\subsection{Triển khai bằng Docker Compose}

Để đơn giản hóa quá trình cài đặt và bảo đảm tính nhất quán giữa các môi trường phát triển, hệ thống được triển khai bằng Docker Compose. Cách tiếp cận này cho phép khởi tạo đồng thời các thành phần của nền tảng Feature Flag thông qua một cấu hình thống nhất, giảm sự phụ thuộc vào môi trường cài đặt của từng máy và hỗ trợ tái tạo hệ thống một cách nhanh chóng \cite{docker_compose}.

\begin{figure}[H]
    \centering
    \includegraphics[width=0.95\textwidth]{\detokenize{Images/mermaid/Triển khai bằng Docker Compose.png}}
    \caption{Các service chính trong quy trình triển khai bằng Docker Compose}
    \label{fig:docker-compose-deployment}
\end{figure}


\subsubsection{Các service trong Docker Compose}

Docker Compose được sử dụng để quản lý các service chính của hệ thống. Trong phiên bản hiện tại, các service bao gồm cơ sở dữ liệu PostgreSQL và các ứng dụng chạy trong môi trường Node.js như Backend API, Admin Dashboard và Demo Application. Khi cần mở rộng, các service này có thể được bổ sung thêm Redis để cache snapshot cấu hình evaluation hoặc các thành phần giám sát mà không làm thay đổi kiến trúc tổng thể.

Mỗi service được khai báo độc lập với image, cổng truy cập, biến môi trường, volume và các phụ thuộc cần thiết. Việc quản lý tập trung giúp quá trình khởi động, dừng và cập nhật toàn bộ hệ thống chỉ cần thực hiện bằng một số lệnh Docker Compose thống nhất.

\subsubsection{Network giữa các service}

Các service được kết nối thông qua mạng nội bộ do Docker Compose tự động tạo. Nhờ đó, Backend API có thể truy cập PostgreSQL thông qua tên service thay vì địa chỉ IP cố định, trong khi Admin Dashboard và Demo Application giao tiếp với Backend API thông qua các endpoint được cấu hình trong biến môi trường.

Việc sử dụng mạng nội bộ giúp cô lập các thành phần của hệ thống với môi trường bên ngoài, giảm sự phụ thuộc vào cấu hình mạng của máy chủ và đơn giản hóa quá trình triển khai trên nhiều thiết bị khác nhau \cite{docker_compose}.

\subsubsection{Volume lưu trữ dữ liệu}

Để bảo đảm dữ liệu không bị mất khi container được khởi động lại hoặc cập nhật phiên bản, PostgreSQL sử dụng Docker Volume để lưu trữ dữ liệu lâu dài. Toàn bộ schema, dữ liệu nghiệp vụ, audit log và các thông tin cấu hình Feature Flag được lưu trên volume thay vì bên trong container.

Việc tách dữ liệu khỏi vòng đời của container giúp hệ thống dễ dàng sao lưu, phục hồi và nâng cấp mà không ảnh hưởng đến dữ liệu đã tồn tại. Đây cũng là phương pháp triển khai được khuyến nghị đối với các dịch vụ sử dụng cơ sở dữ liệu \cite{docker_compose}.

\subsubsection{Biến môi trường và cấu hình runtime}

Các thông số cấu hình của hệ thống được quản lý thông qua biến môi trường nhằm tách biệt mã nguồn với thông tin triển khai. Các biến quan trọng bao gồm chuỗi kết nối PostgreSQL, cổng dịch vụ, khóa bí mật phục vụ xác thực, cấu hình CORS, địa chỉ Backend API và các tham số runtime khác.

Việc sử dụng biến môi trường giúp cùng một mã nguồn có thể triển khai trên nhiều môi trường như development, staging và production mà không cần chỉnh sửa chương trình. Đồng thời, các thông tin nhạy cảm không được ghi trực tiếp trong mã nguồn, góp phần nâng cao tính bảo mật và khả năng quản lý cấu hình của hệ thống \cite{owasp_rest_security,docker_compose}.

\subsection{Quy trình khởi chạy hệ thống}

Hệ thống được thiết kế để có thể khởi tạo nhanh trên một môi trường mới thông qua quy trình gồm chuẩn bị môi trường, khởi tạo cơ sở dữ liệu, sinh dữ liệu mẫu và khởi chạy các dịch vụ. Các bước này bảo đảm toàn bộ thành phần của nền tảng Feature Flag được đồng bộ trước khi đưa vào sử dụng hoặc trình diễn.

\subsubsection{Khởi tạo môi trường}

Bước đầu tiên là cài đặt các công cụ cần thiết như Node.js, npm và Docker Desktop. Sau khi tải mã nguồn, toàn bộ các package trong monorepo được cài đặt thông qua npm workspaces để đồng bộ các phụ thuộc giữa Backend API, Admin Dashboard, Demo Application và JavaScript SDK.

Tiếp theo, các biến môi trường được cấu hình thông qua tệp cấu hình tương ứng. Các tham số như chuỗi kết nối PostgreSQL, cổng dịch vụ, khóa bí mật và địa chỉ Backend API được thiết lập trước khi khởi động hệ thống nhằm bảo đảm các thành phần có thể giao tiếp với nhau một cách chính xác.

\subsubsection{Chạy migration và seed}

Sau khi cơ sở dữ liệu PostgreSQL được khởi động, Prisma Migrate được sử dụng để tạo hoặc cập nhật schema theo phiên bản mới nhất của hệ thống. Quá trình migration đồng thời thiết lập các bảng dữ liệu, khóa ngoại, chỉ mục và các ràng buộc toàn vẹn cần thiết.

Khi migration hoàn tất, các script seed được thực thi để tạo dữ liệu mẫu phục vụ phát triển và kiểm thử. Dữ liệu khởi tạo bao gồm project, environment, feature flag, flag group, rule và sample user, cho phép hệ thống có thể sử dụng ngay sau khi cài đặt mà không cần cấu hình thủ công.

\subsubsection{Khởi chạy các service}

Các service của hệ thống được khởi động thông qua Docker Compose hoặc chạy trực tiếp trong môi trường phát triển. PostgreSQL được khởi động trước, tiếp theo là Backend API, Admin Dashboard và Demo Application. Trong quá trình khởi động, Backend API tự động kết nối đến cơ sở dữ liệu và sẵn sàng phục vụ các yêu cầu từ Dashboard cũng như Demo Application.

Việc khởi chạy theo thứ tự giúp bảo đảm các dịch vụ phụ thuộc luôn sẵn sàng trước khi ứng dụng bắt đầu xử lý yêu cầu, đồng thời giảm khả năng phát sinh lỗi kết nối trong giai đoạn khởi động \cite{docker_compose}.

\subsubsection{Kiểm tra trạng thái hệ thống}

Sau khi toàn bộ service được khởi động, trạng thái hệ thống được kiểm tra bằng cách truy cập giao diện Admin Dashboard, Demo Application và tài liệu Swagger của Backend API. Các endpoint quản lý project, feature flag và evaluation được sử dụng để xác nhận Backend API đã kết nối thành công với cơ sở dữ liệu.

Cuối cùng, một số kịch bản đánh giá Feature Flag như bật hoặc tắt toàn cục, targeting theo role và percentage rollout được thực hiện trên Demo Application nhằm xác nhận Rule Evaluation Engine, Dashboard và cơ sở dữ liệu hoạt động đúng theo thiết kế trước khi tiến hành kiểm thử hoặc trình diễn hệ thống.

\section{Kết quả xây dựng}

Kết quả xây dựng của hệ thống được lưu tại GitHub repository: \url{https://github.com/fabyanbui/feature-flag-platform}. Các minh chứng chi tiết về Backend API, Frontend Dashboard, Demo Application và SDK sẽ được trình bày chi tiết trong phần demo.

\subsection{Backend API}

\begin{figure}[H]
    \centering
    \includegraphics[width=0.95\textwidth]{\detokenize{Images/screenshot/Kết quả xây dựng Backend API.png}}
    \caption{Minh chứng kết quả xây dựng Backend API}
    \label{fig:backend-api-result}
\end{figure}

\subsection{Frontend Dashboard}

\begin{figure}[H]
    \centering
    \includegraphics[width=0.95\textwidth]{\detokenize{Images/screenshot/Kết quả xây dựng Frontend Dashboard.png}}
    \caption{Minh chứng kết quả xây dựng Frontend Dashboard}
    \label{fig:frontend-dashboard-result}
\end{figure}

\subsection{Demo Application}

Demo Application đã được mở rộng từ ứng dụng frontend thuần thành storefront có kèm demo backend API nhẹ. Mỗi feature flag trong demo có một endpoint GET tương ứng dưới \path|/api/demo/features/*|; endpoint này dùng \texttt{@ffp/js-sdk} ở phía server để gọi Evaluation API trước khi trả dữ liệu tính năng. Nếu flag tắt, bị archive, bị group/flag kill switch hoặc SDK không xác minh được trạng thái bật, demo backend trả \texttt{403 FEATURE\_DISABLED}.

Một số thao tác có sẵn trong giao diện storefront được nối với endpoint backend tương ứng với feature flag, không cần thêm panel minh họa riêng. Nhờ đó phần demo thể hiện đồng thời hai lớp tích hợp: frontend ẩn/hiện UI theo flag và backend API vẫn tự guard khi bị gọi trực tiếp.

\begin{figure}[H]
    \centering
    \includegraphics[width=0.95\textwidth]{\detokenize{Images/screenshot/Kết quả xây dựng Demo Application.png}}
    \caption{Minh chứng kết quả xây dựng Demo Application}
    \label{fig:demo-application-result}
\end{figure}

\section{Kiểm thử và đánh giá}

Kế hoạch kiểm thử của dự án tập trung vào ba mức: unit test để kiểm tra từng đơn vị nghiệp vụ quan trọng, E2E test để xác nhận các luồng API chạy qua NestJS application và PostgreSQL, và CI để tự động kiểm tra chất lượng khi có thay đổi mã nguồn. Các nhóm kiểm thử ưu tiên những rủi ro chính của hệ thống Feature Flag: evaluation phải xác định, mặc định an toàn, mutation phải có audit log, cache không làm sai kết quả, phân quyền phải được xử lý ở backend và SDK phải fail-closed khi lỗi.

\subsection{Unit test}

Unit test được triển khai bằng Jest trong hai workspace chính: \texttt{@ffp/backend} và \texttt{@ffp/js-sdk}. Ở backend, unit test kiểm tra Rule Evaluation Engine, stable rollout hash, evaluation service/controller/repository, cache, invalidation, audit log, history, RBAC, DTO validation, exception filter, request context, repository và statistics service. Các test dùng mock hoặc dữ liệu tối thiểu để cô lập đơn vị cần kiểm tra, sau đó đối chiếu output với contract mong đợi như \texttt{enabled}, \texttt{reason}, error code, audit snapshot hoặc permission result.

Với JavaScript SDK, unit test kiểm tra client gọi Evaluation API, validate response, trả helper \texttt{isEnabled}/\texttt{getVariant}, xử lý timeout, HTTP error, response sai định dạng và fallback fail-closed. Nhóm test này bảo đảm ứng dụng khách không tự bật tính năng khi Evaluation API không khả dụng.

\begin{table}[H]
    \centering
    \small
    \caption{Thống kê unit test đã triển khai}
    \label{tab:unit-test-summary}
    \begin{tabular}{p{0.24\textwidth} p{0.36\textwidth} p{0.16\textwidth} p{0.14\textwidth}}
        \toprule
        \textbf{Workspace} & \textbf{Phạm vi kiểm thử chính} & \textbf{Test suite} & \textbf{Test case} \\
        \midrule
        \texttt{@ffp/backend} & Evaluation engine, rollout hash, cache, audit, RBAC, DTO, controller, service, repository, statistics & 52 & 416 \\
        \texttt{@ffp/js-sdk} & SDK client, response validation, helper API, timeout, lỗi mạng và fallback fail-closed & 1 & 21 \\
        Tổng & Unit test toàn workspace qua \texttt{npm run test -- --runInBand} & 53 & 437 \\
        \bottomrule
    \end{tabular}
\end{table}

\subsection{Integration test và E2E test}

Repository đã cấu hình Jest cho integration test qua \texttt{apps/backend/test/jest-integration.json} và script \texttt{npm run test:integration --workspace=@ffp/backend}. Ở trạng thái hiện tại, không có file \texttt{*.integration-spec.ts} riêng; các kiểm thử tích hợp được thể hiện chủ yếu qua nhóm E2E test trong \texttt{apps/backend/test}.

E2E test dùng \texttt{@nestjs/testing} để khởi tạo \texttt{AppModule}, áp dụng global validation pipe, exception filter, request context middleware và demo RBAC giống runtime. Các test gọi HTTP bằng \texttt{supertest}, thao tác dữ liệu qua Prisma và dọn dữ liệu bằng helper test database. Hiện có 10 file \texttt{*.e2e-spec.ts}, bao phủ health check, management API, vertical slice tạo project/flag/rule/evaluate, API hardening, demo flow, flag history, group kill switch, evaluation cache, evaluation statistics và server-resolved RBAC.

\subsection{Luồng CI trong CI/CD}

Luồng CI được triển khai bằng GitHub Actions tại \texttt{.github/workflows/ci.yml}. Workflow chạy khi push hoặc pull request vào nhánh \texttt{develop} và \texttt{main}, đồng thời hỗ trợ chạy thủ công bằng \texttt{workflow\_dispatch}. Job \texttt{validate} sử dụng Ubuntu 24.04, Node.js 22.22.2 và PostgreSQL 16 service để mô phỏng môi trường kiểm thử có cơ sở dữ liệu thật.

Các bước CI lần lượt gồm kiểm tra whitespace bằng \texttt{git diff --check}, cài dependencies bằng \texttt{npm ci}, validate và generate Prisma client, chạy lint, chạy unit test backend, build toàn bộ workspace, apply database migration, chạy seed hai lần để kiểm tra idempotency, chạy integration test và cuối cùng chạy backend E2E test. Luồng này giúp phát hiện sớm lỗi định dạng, lỗi schema, lỗi TypeScript/build, lỗi unit test và lỗi end-to-end trước khi thay đổi được hợp nhất.

\section{Các kịch bản demo}

Các kịch bản dưới đây sẽ được trình bày cụ thể trong phần demo. Mục tiêu là cho thấy hệ thống có thể thay đổi hành vi ứng dụng bằng cấu hình Feature Flag, đồng thời vẫn giữ được tính xác định, khả năng rollback nhanh, audit và quan sát vận hành.

\subsection{Quan sát project, flag và ứng dụng demo}

Kịch bản mở đầu giới thiệu project \texttt{demo-project}, các feature flag đã seed, các nhóm flag như \texttt{checkout-experience} và \texttt{recommendations}, sau đó mở Demo Application để cho thấy ứng dụng đang nhận quyết định từ Evaluation API thông qua SDK.

\subsection{Bật và tắt tính năng toàn cục}

Kịch bản này thay đổi cấu hình serving mode hoặc trạng thái runtime của một flag đang bật toàn cục. Demo Application được evaluate lại với cùng user context để quan sát tính năng xuất hiện hoặc bị tắt mà không cần build/deploy lại mã nguồn.

\subsection{Targeting theo role}

Kịch bản này chuyển giữa các user mẫu như beta customer, regular customer và shop admin. Demo sẽ cho thấy cùng một flag có thể bật cho nhóm người dùng có role phù hợp và tắt cho nhóm còn lại, đồng thời hiển thị \texttt{projectKey}, \texttt{flagKey}, \texttt{enabled} và \texttt{reason}.

\subsection{Rollout theo phần trăm}

Kịch bản này dùng nhiều tài khoản mẫu có \texttt{targetingKey} ổn định để minh họa percentage rollout. Một số tài khoản nhận tính năng mới, số khác vẫn dùng trải nghiệm cũ; khi evaluate lại cùng context, kết quả không thay đổi do stable hashing.

\subsection{Kill switch cấp flag và group}

Kịch bản này bật kill switch ở cấp một flag hoặc cấp group, ví dụ nhóm \texttt{checkout-experience}. Demo sẽ nhấn mạnh rằng trạng thái runtime chuyển sang Off ngay lập tức, trong khi lifecycle/configuration của từng flag vẫn được giữ lại để có thể khôi phục nhanh.

\subsection{Audit log, history và statistics}

Kịch bản này thực hiện một số mutation cấu hình và nhiều lần evaluation từ Demo Application, sau đó mở Dashboard để xem audit log, configuration history và statistics. Phần demo sẽ cho thấy actor, action, request ID, before/after snapshot và số liệu aggregate theo flag/reason.

\subsection{Fallback của JavaScript SDK}

Kịch bản này mô phỏng Evaluation API lỗi hoặc timeout. SDK \texttt{@ffp/js-sdk} sẽ trả kết quả fail-closed với \texttt{enabled=false}, giúp Demo Application tiếp tục hoạt động ở trạng thái an toàn và phân biệt được lỗi client-local với quyết định từ backend.

\subsection{Phân quyền Admin, Developer và Viewer}

Kịch bản cuối cùng chuyển giữa ba vai trò demo trên Admin Dashboard. Viewer chỉ xem dữ liệu, Developer được thao tác trong phạm vi cấu hình được phép, còn Admin có quyền quản trị đầy đủ; các mutation bị từ chối sẽ được backend bảo vệ thay vì chỉ ẩn ở giao diện.


% ------------------------------------------------------------------------------
% Chapter 4
% ------------------------------------------------------------------------------
\chapter{Kết luận} % enter the name of the chapter here
\label{cha:ket-luan} % enter the chapter label here (for cross referencing)

\section{Tổng kết kết quả đạt được}

Đề tài đã hoàn thành mục tiêu nghiên cứu Feature Flag và xây dựng một nền tảng quản lý phát hành tính năng an toàn ở phạm vi mini project. Báo cáo đã trình bày nền tảng lý thuyết, khảo sát giải pháp hiện có, phân tích yêu cầu, thiết kế kiến trúc, thiết kế dữ liệu, Rule Evaluation Engine, API, bảo mật, kiểm thử và kết quả triển khai.

Về sản phẩm, hệ thống gồm Backend API, Admin Dashboard, Demo Application, JavaScript SDK và PostgreSQL. Các chức năng bắt buộc được triển khai đầy đủ; nhiều chức năng khuyến khích như cache snapshot, audit-backed history, RBAC demo, statistics, group kill switch và Docker Compose cũng được hoàn thành.

Về kiểm chứng, hệ thống đã được kiểm thử ở các lớp engine, API, persistence, cache, RBAC, giao diện, demo app và quy trình triển khai. Các kịch bản trình diễn cho thấy người quản trị có thể thay đổi cấu hình flag trên Dashboard và quan sát hành vi ứng dụng thay đổi thông qua Evaluation API mà không cần deploy lại.

Như vậy, kết quả đạt được không chỉ là một ứng dụng quản lý bật/tắt tính năng, mà là một mô hình nhỏ thể hiện được tư duy release management, thiết kế hệ thống, phân tách control plane/data plane, fail-closed evaluation và auditability trong vận hành phần mềm.

\section{Hạn chế và hướng phát triển}

\subsection{Hạn chế}

Mặc dù đáp ứng mục tiêu mini project, hệ thống vẫn còn một số hạn chế khi so với nền tảng Feature Flag ở mức production:

\begin{itemize}
    \item Hệ thống chưa được đánh giá dưới tải lớn; các kiểm thử hiện tập trung chủ yếu vào tính đúng đắn, luồng API và kịch bản demo.
    \item Kiến trúc hiện phù hợp với một backend service trong môi trường demo, chưa có cơ chế đồng bộ cấu hình/cache giữa nhiều instance.
    \item SDK mới có phiên bản JavaScript và vẫn gọi Evaluation API theo request/response, chưa hỗ trợ streaming hoặc push update thời gian thực.
    \item Demo backend hiện chỉ là minh họa server-side feature gating đơn giản, chưa phải Backend SDK hoàn chỉnh và chưa có middleware route guard tái sử dụng cho nhiều framework.
    \item Rule engine mới hỗ trợ các rule cơ bản như user allowlist, role targeting và percentage rollout, chưa có targeting theo thuộc tính phức tạp, thời gian, vị trí hoặc organization.
    \item Hệ thống chưa hỗ trợ multi-variant experiment và A/B testing hoàn chỉnh, dù response đã có khái niệm variant đơn giản.
    \item RBAC hiện phục vụ demo, chưa tích hợp OAuth2/OIDC, SSO, quản lý người dùng thật hoặc chính sách bảo mật ở mức doanh nghiệp.
    \item Chưa có approval workflow, stale flag detection, rate limiting đầy đủ, monitoring hạ tầng, backup/recovery và disaster recovery.
\end{itemize}

\subsection{Hướng phát triển}

Các hướng phát triển tương ứng với các hạn chế trên gồm:

\begin{itemize}
    \item Bổ sung benchmark và load test cho Evaluation API, đo latency, throughput, cache hit/miss và hành vi khi database hoặc cache gặp lỗi.
    \item Mở rộng kiến trúc sang nhiều backend instance, dùng distributed cache hoặc event-based invalidation để đồng bộ snapshot cấu hình giữa các instance.
    \item Phát triển thêm SDK cho các ngôn ngữ phổ biến như Java, Go, Python hoặc C\#, đồng thời nghiên cứu polling tối ưu, SSE/WebSocket hoặc streaming update.
    \item Phát triển Backend SDK hoặc middleware chính thức cho Node.js, NestJS hoặc Express để backend nghiệp vụ có thể guard route/service bằng Feature Flag một cách chuẩn hóa.
    \item Mở rộng Rule Evaluation Engine để hỗ trợ targeting theo thuộc tính người dùng, organization/tenant, thời gian, khu vực địa lý và điều kiện kết hợp.
    \item Thiết kế mô hình multi-variant flag và A/B testing, kết hợp statistics để đo lường hiệu quả từng biến thể.
    \item Thay thế mô hình RBAC demo bằng cơ chế xác thực chuẩn như OAuth2/OIDC hoặc SSO, kèm quản lý người dùng, nhóm quyền và chính sách truy cập chi tiết.
    \item Bổ sung approval workflow cho thay đổi production, stale flag detection để giảm feature flag debt, rate limiting, monitoring/alerting, backup/recovery và tài liệu vận hành production.
\end{itemize}

\section{Bài học kinh nghiệm}

Quá trình thực hiện đề tài cho thấy giá trị của việc tách rõ nghiên cứu, thiết kế, triển khai và kiểm thử. Các bài học dưới đây được trình bày ở mức tổng kết, không lặp lại chi tiết kỹ thuật đã phân tích ở các chương trước.

\subsection{Bài học về phân tích bài toán}

\subsubsection{Xác định phạm vi MVP}

Feature Flag là miền bài toán rộng, dễ mở sang A/B testing, realtime sync hoặc workflow phê duyệt. Việc xác định MVP giúp ưu tiên các năng lực cốt lõi: quản lý flag, rule engine xác định, audit log, dashboard và demo tích hợp.

\subsubsection{Phân tách yêu cầu bắt buộc và mở rộng}

Tách yêu cầu bắt buộc khỏi yêu cầu khuyến khích giúp sản phẩm có baseline ổn định trước khi bổ sung SDK, RBAC, statistics, group kill switch và Docker Compose.

\subsubsection{Chuyển yêu cầu thành use case và thiết kế}

Các yêu cầu ban đầu cần được chuyển thành use case, contract API, mô hình dữ liệu và tiêu chí kiểm thử cụ thể; nếu không, hệ thống dễ chỉ dừng ở giao diện bật/tắt đơn giản.

\subsection{Bài học về thiết kế hệ thống}

\subsubsection{Tầm quan trọng của evaluation precedence}

Một rule engine chỉ đáng tin cậy khi thứ tự ưu tiên được định nghĩa trước và được kiểm thử. Đây là điểm quyết định để kết quả evaluation có thể giải thích được.

\subsubsection{Phân tách định nghĩa flag và cấu hình runtime}

Tách metadata của flag khỏi cấu hình theo environment giúp cùng một flag phục vụ nhiều môi trường mà không nhân bản định nghĩa.

\subsubsection{Thiết kế auditability từ đầu}

Audit log nên được đưa vào luồng mutation ngay từ đầu, vì bổ sung sau khi hệ thống đã có nhiều endpoint sẽ dễ bỏ sót hành động quan trọng.

\subsubsection{Cân bằng giữa tính đơn giản và khả năng mở rộng}

Kiến trúc một backend service đủ đơn giản cho mini project nhưng vẫn giữ ranh giới module để có thể tách service hoặc thêm hạ tầng phân tán sau này.

\subsection{Bài học về triển khai và kiểm thử}

\subsubsection{Validation tại ranh giới hệ thống}

DTO validation và constraint database giúp phát hiện lỗi sớm, giảm trạng thái cấu hình không hợp lệ lọt vào engine.

\subsubsection{Kiểm thử các trường hợp precedence}

Các test quan trọng nhất là những test chứng minh precedence, fail-closed và stable rollout không thay đổi khi refactor.

\subsubsection{Quản lý cache invalidation}

Cache chỉ an toàn khi có chiến lược invalidation rõ. Với Feature Flag, snapshot cũ có thể trực tiếp làm sai quyết định runtime.

\subsubsection{Xây dựng dữ liệu và kịch bản demo}

Seed data và kịch bản demo cần được chuẩn bị cùng lúc với tính năng để người xem có thể quan sát giá trị hệ thống mà không phải cấu hình thủ công quá nhiều.

\section{Hạn chế}

\subsection{Hạn chế về kiến trúc}

Mặc dù hệ thống đã đáp ứng đầy đủ các chức năng của một nền tảng Feature Flag trong phạm vi mini project, kiến trúc hiện tại vẫn còn một số hạn chế khi xét tới khả năng triển khai ở quy mô lớn hoặc trong môi trường phân tán. Các hạn chế này chủ yếu liên quan đến khả năng mở rộng, đồng bộ cấu hình và đánh giá hiệu năng dưới tải cao, đồng thời cũng là những hướng phát triển phổ biến của các nền tảng Feature Flag thương mại như LaunchDarkly hay Unleash \cite{launchdarkly_feature_flags,unleash_feature_flag_best_practices}.

\subsubsection{Chưa hỗ trợ triển khai phân tán}

Kiến trúc của hệ thống được thiết kế để triển khai trên một cụm dịch vụ đơn giản phục vụ mục tiêu học tập và trình diễn. Backend API, Evaluation Engine và Dashboard sử dụng chung cơ sở dữ liệu, trong khi cơ chế cache chủ yếu hoạt động trên từng instance của ứng dụng. Mô hình này giúp giảm độ phức tạp trong quá trình phát triển nhưng chưa giải quyết các bài toán thường gặp trong môi trường nhiều node như đồng bộ cache, cân bằng tải hay đảm bảo tính nhất quán giữa các instance.

Trong các hệ thống Feature Flag thực tế, Evaluation Service thường được triển khai trên nhiều máy chủ kết hợp với cơ chế cache phân tán và đồng bộ cấu hình nhằm duy trì độ trễ thấp và khả năng chịu lỗi cao \cite{unleash_feature_flag_best_practices,launchdarkly_feature_flags}. Đây là hướng mở rộng chưa được triển khai trong phạm vi của đề tài.

\subsubsection{Chưa có cơ chế đồng bộ cấu hình realtime}

Sau khi cấu hình Feature Flag được thay đổi trên Dashboard, các instance của hệ thống hiện dựa vào cơ chế cache invalidation và tải lại dữ liệu khi cần thiết. Hệ thống chưa hỗ trợ các cơ chế truyền sự kiện thời gian thực như WebSocket, Server-Sent Events hoặc message broker để thông báo ngay lập tức tới các Evaluation Service hoặc SDK đang hoạt động.

Do đó, trong trường hợp triển khai nhiều instance hoặc sử dụng cache với thời gian sống dài, có thể tồn tại khoảng thời gian ngắn mà các instance chưa đồng bộ hoàn toàn cấu hình mới. Trong các nền tảng thương mại, việc đồng bộ cấu hình theo thời gian thực là một thành phần quan trọng nhằm giảm độ trễ giữa thao tác quản trị và thời điểm cấu hình có hiệu lực trên toàn hệ thống \cite{launchdarkly_feature_flags,unleash_feature_flag_best_practices}.

\subsubsection{Chưa đánh giá ở lưu lượng lớn}

Các thử nghiệm của đề tài chủ yếu tập trung vào việc kiểm chứng tính đúng đắn của Rule Evaluation Engine, các kịch bản rollout và khả năng hoạt động của toàn bộ hệ thống trong môi trường demo. Chưa thực hiện các bài kiểm thử hiệu năng với số lượng lớn request đồng thời, nhiều instance hoặc khối lượng dữ liệu lớn.

Vì vậy, báo cáo chưa đưa ra các chỉ số định lượng về throughput, độ trễ, khả năng mở rộng theo chiều ngang hay mức sử dụng tài nguyên của hệ thống. Việc xây dựng bộ kiểm thử tải bằng các công cụ chuyên dụng và đánh giá khả năng mở rộng trong nhiều điều kiện triển khai khác nhau sẽ là hướng nghiên cứu và hoàn thiện tiếp theo.

\subsection{Hạn chế về chức năng}

Bên cạnh việc đáp ứng đầy đủ các chức năng của mức MVP và hầu hết các chức năng khuyến khích trong đề bài, hệ thống vẫn còn một số hạn chế về mặt chức năng so với các nền tảng Feature Flag thương mại. Các hạn chế này chủ yếu nhằm giữ phạm vi mini project ở mức phù hợp, giảm độ phức tạp trong thiết kế và tập trung vào việc làm rõ nguyên lý hoạt động của Feature Flag thay vì xây dựng một nền tảng hoàn chỉnh \cite{launchdarkly_feature_flags,unleash_feature_flag_best_practices}.

\subsubsection{Số loại rule còn giới hạn}

Rule Evaluation Engine hiện hỗ trợ ba loại rule chính gồm targeting theo người dùng, targeting theo vai trò và percentage rollout. Đây là các chiến lược phổ biến, đủ để minh họa nguyên lý tách biệt giữa deploy và release cũng như đáp ứng các yêu cầu của đề tài.

Tuy nhiên, hệ thống chưa hỗ trợ các rule phức tạp như điều kiện theo nhiều thuộc tính kết hợp, biểu thức logic AND/OR lồng nhau, targeting theo thời gian, vị trí địa lý, phiên bản ứng dụng hoặc các phân khúc người dùng được tạo động. Trong các nền tảng Feature Flag hiện đại, những khả năng này giúp đáp ứng nhiều kịch bản phát hành sản phẩm và cá nhân hóa hơn \cite{launchdarkly_targeting_rules,openfeature_evaluation_context}.

\subsubsection{Chưa hỗ trợ nhiều variant đầy đủ}

Hệ thống hiện tập trung vào mô hình Feature Flag dạng bật hoặc tắt, đồng thời hỗ trợ trả về giá trị variant cơ bản phục vụ quá trình evaluation. Tuy nhiên, chưa xây dựng đầy đủ cơ chế quản lý nhiều variant với trọng số phân phối, cấu hình riêng cho từng variant hoặc rollout độc lập theo từng biến thể.

Việc mở rộng sang mô hình đa variant sẽ cho phép triển khai các kịch bản như A/B testing, multivariate testing hoặc phát hành nhiều phiên bản giao diện đồng thời. Đây là một trong những khả năng quan trọng của các nền tảng Feature Flag thương mại nhưng nằm ngoài phạm vi của mini project \cite{launchdarkly_feature_flags,fowler_feature_toggles}.

\subsubsection{Chưa hỗ trợ approval workflow}

Mọi thay đổi cấu hình Feature Flag trong hệ thống đều được áp dụng ngay sau khi người dùng có đủ quyền thực hiện thao tác. Hệ thống đã ghi nhận đầy đủ audit log và lịch sử thay đổi nhưng chưa hỗ trợ quy trình phê duyệt trước khi cấu hình được áp dụng.

Trong môi trường doanh nghiệp, các thay đổi quan trọng thường phải trải qua quy trình nhiều bước như tạo yêu cầu, xem xét, phê duyệt và mới được phát hành nhằm giảm nguy cơ cấu hình sai hoặc thay đổi ngoài ý muốn. Chức năng này chưa được triển khai do vượt quá phạm vi của đề tài.

\subsubsection{Chưa hỗ trợ stale flag detection}

Hệ thống chưa cung cấp cơ chế tự động phát hiện các Feature Flag không còn được sử dụng hoặc đã tồn tại quá lâu trong mã nguồn. Việc quản lý vòng đời của Feature Flag hiện vẫn dựa trên thao tác thủ công của người quản trị.

Trong thực tế, việc tồn tại quá nhiều Feature Flag không còn giá trị sẽ làm tăng độ phức tạp của mã nguồn, gây khó khăn cho bảo trì và hình thành \textit{feature flag debt}. Nhiều nền tảng hiện nay cung cấp khả năng theo dõi mức độ sử dụng, cảnh báo stale flag hoặc hỗ trợ dọn dẹp các flag không còn cần thiết nhằm duy trì chất lượng hệ thống \cite{mahdavi_hezaveh_feature_toggles,unleash_feature_flag_best_practices}.

\subsection{Hạn chế về bảo mật và vận hành}

Hệ thống đã áp dụng các cơ chế xác thực, phân quyền, validation dữ liệu, audit log và cache invalidation nhằm đáp ứng yêu cầu của đề tài. Tuy nhiên, nếu triển khai trong môi trường sản xuất với nhiều người dùng và lưu lượng truy cập lớn, hệ thống vẫn cần được bổ sung nhiều cơ chế về bảo mật và vận hành để đáp ứng các yêu cầu về tính sẵn sàng, khả năng giám sát và an toàn dữ liệu \cite{owasp_rest_security,owasp_logging}.

\subsubsection{Mô hình xác thực còn đơn giản}

Hệ thống sử dụng cơ chế xác thực và phân quyền phù hợp với phạm vi của mini project, cho phép phân biệt các vai trò quản trị như administrator, developer và viewer. Tuy nhiên, mô hình xác thực hiện vẫn còn đơn giản và chưa tích hợp với các hệ thống quản lý danh tính phổ biến như OAuth~2.0, OpenID Connect hoặc Single Sign-On (SSO).

Ngoài ra, hệ thống cũng chưa hỗ trợ các cơ chế như xác thực đa yếu tố (Multi-Factor Authentication), quản lý phiên đăng nhập tập trung hay chính sách quản lý khóa truy cập dành cho các hệ thống tích hợp. Đây đều là những yêu cầu thường gặp đối với các nền tảng quản trị trong môi trường doanh nghiệp \cite{owasp_rest_security}.

\subsubsection{Chưa có rate limiting đầy đủ}

Các API của hệ thống hiện chưa triển khai cơ chế giới hạn tốc độ truy cập đối với từng người dùng hoặc từng địa chỉ IP. Trong điều kiện bình thường, điều này không ảnh hưởng đến việc trình diễn và đánh giá chức năng của hệ thống. Tuy nhiên, khi triển khai thực tế, việc thiếu rate limiting có thể làm tăng nguy cơ bị khai thác thông qua các cuộc tấn công gửi số lượng lớn request hoặc gây quá tải cho dịch vụ.

Việc bổ sung rate limiting, quota theo API, giới hạn số lần evaluation trong một khoảng thời gian và cơ chế phát hiện bất thường sẽ góp phần nâng cao khả năng bảo vệ hệ thống trước các hành vi truy cập không hợp lệ \cite{owasp_rest_security}.

\subsubsection{Chưa có hệ thống monitoring chuyên biệt}

Hệ thống hiện ghi nhận audit log và thống kê số lần evaluation phục vụ việc quản trị Feature Flag, nhưng chưa tích hợp các nền tảng monitoring chuyên biệt để theo dõi trạng thái vận hành của toàn bộ hệ thống.

Các chỉ số như thời gian phản hồi của Evaluation API, tỷ lệ lỗi, mức sử dụng CPU, bộ nhớ, trạng thái cache hay cảnh báo khi dịch vụ gặp sự cố vẫn chưa được thu thập và trực quan hóa. Trong môi trường sản xuất, việc kết hợp các công cụ monitoring và alerting sẽ giúp phát hiện sớm sự cố, rút ngắn thời gian xử lý và nâng cao độ tin cậy của hệ thống.

\subsubsection{Chưa triển khai cơ chế backup và recovery}

Hệ thống chưa xây dựng quy trình sao lưu và khôi phục dữ liệu cho cơ sở dữ liệu cũng như các thành phần cấu hình của Feature Flag. Mặc dù toàn bộ thay đổi đều được lưu trong audit log và dữ liệu được quản lý tập trung, nhưng chưa có cơ chế tự động tạo bản sao lưu định kỳ hoặc phục hồi sau sự cố.

Trong môi trường triển khai thực tế, việc xây dựng chiến lược backup và disaster recovery là cần thiết nhằm giảm thiểu nguy cơ mất dữ liệu do lỗi phần cứng, thao tác nhầm hoặc sự cố hệ thống. Đây là một thành phần quan trọng để đảm bảo tính sẵn sàng và khả năng vận hành liên tục của nền tảng.

\section{Hướng phát triển}

\section{Mở rộng Rule Evaluation Engine}

Mặc dù Rule Evaluation Engine đã đáp ứng đầy đủ các yêu cầu của mini project với các chiến lược bật hoặc tắt toàn cục, targeting theo người dùng, targeting theo vai trò và percentage rollout, hệ thống vẫn còn nhiều hướng mở rộng nhằm đáp ứng các nhu cầu triển khai trong môi trường sản xuất quy mô lớn. Việc mở rộng engine theo hướng mô-đun sẽ giúp bổ sung các loại rule mới mà không ảnh hưởng đến các rule hiện có, đồng thời phù hợp với định hướng phát triển của các nền tảng Feature Flag hiện đại \cite{launchdarkly_feature_flags,unleash_feature_flag_best_practices}.

\subsection{Targeting theo thuộc tính}

Hiện tại hệ thống chủ yếu đánh giá dựa trên \texttt{userId}, \texttt{role} và \texttt{targetingKey}. Trong tương lai, Evaluation Engine có thể hỗ trợ nhiều thuộc tính hơn như loại thiết bị, hệ điều hành, phiên bản ứng dụng, gói dịch vụ, quốc gia hoặc các thuộc tính nghiệp vụ do ứng dụng khách cung cấp thông qua Evaluation Context.

Các điều kiện này có thể được biểu diễn dưới dạng tập luật với các toán tử so sánh như bằng, khác, thuộc tập giá trị, so khớp chuỗi hoặc biểu thức chính quy. Cách tiếp cận này giúp hệ thống linh hoạt hơn trong việc triển khai tính năng cho các nhóm người dùng cụ thể mà không cần thay đổi mã nguồn của ứng dụng khách. Hướng phát triển này cũng phù hợp với mô hình Evaluation Context được đề xuất trong OpenFeature \cite{openfeature_evaluation_context}.

\subsection{Targeting theo organization}

Trong nhiều hệ thống SaaS, Feature Flag thường được quản lý theo từng khách hàng hoặc từng tổ chức thay vì từng người dùng riêng lẻ. Do đó, một hướng phát triển quan trọng là bổ sung khả năng targeting dựa trên \texttt{attributes.organizationId} hoặc các thuộc tính của tổ chức.

Khi đó, Evaluation Engine có thể kích hoạt hoặc vô hiệu hóa tính năng cho toàn bộ doanh nghiệp, nhóm khách hàng hoặc tenant cụ thể. Việc này giúp đơn giản hóa quá trình quản trị khi số lượng người dùng lớn, đồng thời hỗ trợ triển khai các gói tính năng khác nhau cho từng khách hàng doanh nghiệp mà vẫn sử dụng chung một hệ thống.

\subsection{Targeting theo thời gian và vị trí}

Một hướng phát triển khác là bổ sung các rule phụ thuộc vào thời gian hoặc vị trí địa lý. Ví dụ, Feature Flag chỉ được kích hoạt trong một khoảng thời gian xác định để phục vụ chiến dịch marketing, sự kiện đặc biệt hoặc quá trình phát hành theo từng giai đoạn.

Ngoài ra, hệ thống có thể hỗ trợ triển khai tính năng theo khu vực hoặc quốc gia nhằm đáp ứng yêu cầu pháp lý hoặc chiến lược phát hành riêng của từng thị trường. Để thực hiện điều này, Evaluation Context cần được mở rộng với các thông tin về thời gian, múi giờ hoặc vị trí, đồng thời Evaluation Engine cần bổ sung các bộ đánh giá tương ứng mà vẫn đảm bảo thứ tự ưu tiên giữa các loại rule.

\subsection{Hỗ trợ nhiều variant và A/B testing}

Phiên bản hiện tại chủ yếu trả về kết quả bật hoặc tắt cùng một giá trị \texttt{variant} đơn giản. Trong tương lai, hệ thống có thể mở rộng để hỗ trợ nhiều biến thể của cùng một Feature Flag, cho phép phân phối nhiều phiên bản giao diện hoặc thuật toán khác nhau đến các nhóm người dùng.

Trên cơ sở đó, nền tảng có thể hỗ trợ các kịch bản A/B testing hoặc multivariate testing, trong đó người dùng được phân bổ ổn định vào từng biến thể thông qua cơ chế băm nhất quán. Kết hợp với hệ thống thống kê evaluation và các chỉ số nghiệp vụ, kết quả thử nghiệm có thể được sử dụng để đánh giá hiệu quả của từng biến thể trước khi triển khai rộng rãi. Đây là một trong những khả năng quan trọng của các nền tảng Feature Flag hiện đại và là hướng phát triển phù hợp cho hệ thống trong tương lai \cite{launchdarkly_percentage_rollouts,fowler_feature_toggles}.

\section{Mở rộng khả năng tích hợp}

Trong phạm vi mini project, hệ thống đã cung cấp REST API và JavaScript SDK để các ứng dụng web có thể tích hợp và sử dụng Feature Flag. Tuy nhiên, nhằm đáp ứng nhu cầu triển khai trong nhiều môi trường và ngôn ngữ lập trình khác nhau, nền tảng vẫn còn nhiều hướng phát triển về khả năng tích hợp với các hệ thống bên ngoài, đồng thời nâng cao khả năng đồng bộ cấu hình và tính tương thích với các tiêu chuẩn mở của ngành \cite{openfeature_flag_evaluation,launchdarkly_feature_flags}.

\subsection{SDK cho nhiều ngôn ngữ}

Hiện nay SDK mới được xây dựng cho ứng dụng JavaScript nhằm minh họa quá trình gọi Evaluation API. Trong tương lai, hệ thống có thể phát triển thêm SDK cho các ngôn ngữ phổ biến như Java, Go, Python, C\#, PHP hoặc Kotlin để hỗ trợ nhiều loại ứng dụng khác nhau.

Các SDK nên cung cấp giao diện lập trình thống nhất, hỗ trợ quản lý Evaluation Context, xử lý lỗi, caching cục bộ và trả về kết quả đánh giá theo cùng một định dạng. Điều này giúp giảm đáng kể khối lượng mã nguồn tích hợp phía ứng dụng khách, đồng thời đảm bảo hành vi đánh giá Feature Flag nhất quán trên nhiều nền tảng.

\subsection{Webhook khi cấu hình thay đổi}

Một hướng phát triển khác là bổ sung cơ chế webhook để thông báo cho các hệ thống bên ngoài khi cấu hình Feature Flag thay đổi. Mỗi khi Feature Flag, Rule hoặc Environment được cập nhật, hệ thống có thể gửi sự kiện đến các URL đã đăng ký trước.

Webhook giúp các hệ thống tích hợp chủ động thực hiện các tác vụ như làm mới cache, đồng bộ dữ liệu, ghi nhận sự kiện hoặc kích hoạt quy trình CI/CD mà không cần liên tục truy vấn Management API. Cơ chế này cũng góp phần giảm lưu lượng truy cập và tăng khả năng mở rộng của toàn hệ thống.

\subsection{Đồng bộ cấu hình theo thời gian thực}

Phiên bản hiện tại chủ yếu sử dụng cơ chế cache kết hợp với invalidation sau khi có thay đổi cấu hình. Trong tương lai, hệ thống có thể mở rộng theo hướng đồng bộ cấu hình theo thời gian thực thông qua các công nghệ như WebSocket hoặc Server-Sent Events.

Khi đó, các SDK hoặc Evaluation Service có thể nhận ngay thông tin cập nhật mà không cần polling định kỳ. Việc đồng bộ theo thời gian thực giúp giảm độ trễ giữa thời điểm thay đổi cấu hình và thời điểm áp dụng trong ứng dụng khách, đồng thời cải thiện tính nhất quán của dữ liệu trong môi trường có nhiều instance triển khai \cite{unleash_feature_flag_best_practices}.

\subsection{Tương thích OpenFeature}

OpenFeature đang dần trở thành tiêu chuẩn mở cho việc tích hợp Feature Flag giữa nhiều nền tảng khác nhau. Do đó, một hướng phát triển quan trọng là điều chỉnh SDK và Evaluation API để tương thích với các đặc tả của OpenFeature.

Việc hỗ trợ đầy đủ các khái niệm như Provider, Evaluation Context, Resolution Details và Evaluation API sẽ giúp ứng dụng khách có thể chuyển đổi giữa nhiều nhà cung cấp Feature Flag mà không cần thay đổi đáng kể mã nguồn. Điều này không chỉ nâng cao khả năng tái sử dụng của hệ thống mà còn tạo điều kiện để nền tảng dễ dàng tích hợp với hệ sinh thái Feature Flag hiện đại trong tương lai \cite{openfeature_evaluation_context,openfeature_flag_evaluation}.

\section{Mở rộng quản trị vòng đời flag}

Trong phiên bản hiện tại, hệ thống đã hỗ trợ các trạng thái cơ bản của Feature Flag cùng khả năng lưu vết thay đổi thông qua audit log và audit-backed history. Tuy nhiên, việc quản trị vòng đời của Feature Flag trong môi trường sản xuất còn bao gồm nhiều quy trình liên quan đến phê duyệt, lập lịch triển khai và quản lý các flag không còn sử dụng. Đây là những chức năng thường xuất hiện trong các nền tảng Feature Flag quy mô lớn nhằm giảm rủi ro vận hành và hạn chế phát sinh technical debt \cite{fowler_feature_toggles,unleash_feature_flag_best_practices}.

\subsection{Approval workflow}

Một hướng phát triển quan trọng là xây dựng quy trình phê duyệt trước khi các thay đổi cấu hình được áp dụng vào môi trường production. Thay vì cho phép cập nhật trực tiếp, các thao tác quan trọng như bật hoặc tắt Feature Flag, thay đổi rule hoặc điều chỉnh rollout có thể được chuyển sang trạng thái chờ phê duyệt.

Quy trình này có thể phân tách vai trò giữa người đề xuất và người phê duyệt, đồng thời lưu đầy đủ lịch sử của từng bước xử lý trong audit log. Cách tiếp cận này giúp giảm nguy cơ cấu hình sai, tăng khả năng kiểm soát thay đổi và đáp ứng tốt hơn các yêu cầu về quản trị trong doanh nghiệp.

\subsection{Lập lịch bật và tắt flag}

Hiện tại mọi thay đổi đều được thực hiện thủ công thông qua Dashboard hoặc API. Trong tương lai, hệ thống có thể hỗ trợ lập lịch tự động kích hoạt hoặc vô hiệu hóa Feature Flag tại một thời điểm xác định.

Chức năng này đặc biệt hữu ích đối với các chiến dịch marketing, sự kiện theo mùa hoặc kế hoạch phát hành đã được chuẩn bị trước. Việc tự động thực thi theo lịch giúp giảm thao tác thủ công, hạn chế sai sót của con người và đảm bảo các thay đổi được áp dụng đúng thời điểm dự kiến.

\subsection{Phát hiện stale flag}

Sau khi một tính năng đã được triển khai ổn định, nhiều Feature Flag không còn giá trị sử dụng nhưng vẫn tồn tại trong hệ thống và trong mã nguồn ứng dụng. Điều này làm tăng độ phức tạp của logic xử lý và gây khó khăn cho quá trình bảo trì.

Do đó, hệ thống có thể được mở rộng để phát hiện các stale flag dựa trên nhiều tiêu chí như thời gian không cập nhật, thời gian không được evaluate, trạng thái luôn bật hoặc luôn tắt trong thời gian dài hoặc không còn được tham chiếu trong ứng dụng. Việc nhận diện sớm các stale flag giúp giảm technical debt và duy trì hệ thống ở trạng thái dễ quản lý hơn \cite{fowler_feature_toggles,mahdavi_hezaveh_feature_toggles}.

\subsection{Nhắc nhở và tự động dọn dẹp flag}

Sau khi phát hiện stale flag, nền tảng có thể tiếp tục hỗ trợ cơ chế nhắc nhở hoặc tự động xử lý theo các chính sách được cấu hình trước. Ví dụ, hệ thống có thể gửi thông báo cho quản trị viên khi Feature Flag vượt quá thời gian tồn tại dự kiến hoặc đề xuất chuyển sang trạng thái lưu trữ.

Trong các phiên bản nâng cao, hệ thống còn có thể hỗ trợ quy trình tự động archive hoặc đánh dấu các Feature Flag cần được loại bỏ sau khi đã hoàn tất quá trình kiểm tra. Việc quản lý vòng đời theo hướng chủ động giúp hạn chế sự tích lũy của Feature Flag không còn sử dụng, giảm chi phí bảo trì và phù hợp với các thực tiễn được khuyến nghị trong các nền tảng Feature Flag hiện đại \cite{unleash_feature_flag_best_practices,mahdavi_hezaveh_feature_toggles}.

\section{Mở rộng khả năng vận hành}

Phiên bản hiện tại đáp ứng tốt mục tiêu nghiên cứu và triển khai trong phạm vi mini project với quy mô nhỏ. Tuy nhiên, khi số lượng Feature Flag, người dùng và yêu cầu truy cập đồng thời tăng lên, hệ thống cần được mở rộng về khả năng vận hành nhằm đảm bảo tính sẵn sàng, hiệu năng và độ tin cậy. Các hướng phát triển dưới đây tập trung vào việc nâng cao khả năng mở rộng của hạ tầng mà vẫn giữ nguyên kiến trúc tổng thể của nền tảng \cite{unleash_feature_flag_best_practices,docker_compose}.

\subsection{High availability}

Một hướng phát triển quan trọng là triển khai nhiều instance cho Backend API và Dashboard nhằm loại bỏ điểm lỗi đơn (Single Point of Failure). Các instance có thể được đặt sau bộ cân bằng tải để phân phối lưu lượng truy cập và tự động chuyển hướng khi một nút gặp sự cố.

Bên cạnh đó, cơ sở dữ liệu cũng có thể được mở rộng theo mô hình primary--replica nhằm tăng khả năng chịu lỗi và cải thiện hiệu năng đọc. Cách tiếp cận này giúp Evaluation API duy trì khả năng phục vụ ngay cả khi một phần hạ tầng gặp sự cố, từ đó nâng cao độ sẵn sàng của toàn hệ thống.

\subsection{Distributed cache}

Hiện tại hệ thống sử dụng cơ chế cache đơn giản phục vụ quá trình evaluation. Trong môi trường có nhiều Backend API hoạt động đồng thời, mỗi instance sẽ có bộ nhớ đệm riêng, dẫn đến nguy cơ dữ liệu không đồng bộ sau khi cấu hình thay đổi.

Để khắc phục hạn chế này, hệ thống có thể sử dụng distributed cache dựa trên Redis hoặc các giải pháp tương tự làm tầng cache dùng chung giữa các instance. Kết hợp với cơ chế cache invalidation hoặc publish/subscribe, mọi thay đổi cấu hình có thể được đồng bộ nhanh chóng đến toàn bộ hệ thống, góp phần giảm độ trễ và duy trì tính nhất quán của dữ liệu \cite{redis_client_side_caching}.

\subsection{Monitoring và alerting}

Một hướng phát triển khác là xây dựng hệ thống giám sát và cảnh báo chuyên biệt cho nền tảng Feature Flag. Các chỉ số như số lượng evaluation, thời gian phản hồi API, tỷ lệ lỗi, số lần cache hit hoặc cache miss và mức sử dụng tài nguyên có thể được thu thập và trực quan hóa thông qua các công cụ giám sát.

Kết hợp với cơ chế cảnh báo tự động, hệ thống có thể nhanh chóng phát hiện các bất thường như tỷ lệ lỗi tăng cao, Evaluation API không khả dụng hoặc số lượng request vượt ngưỡng cho phép. Điều này giúp đội ngũ vận hành phản ứng kịp thời và giảm thiểu ảnh hưởng đến các ứng dụng đang sử dụng Feature Flag.

\subsection{Triển khai trên Kubernetes}

Phiên bản hiện tại sử dụng Docker Compose để đơn giản hóa quá trình triển khai trong môi trường phát triển và trình diễn. Trong tương lai, hệ thống có thể được mở rộng để triển khai trên Kubernetes nhằm đáp ứng yêu cầu vận hành ở quy mô lớn.

Việc sử dụng Kubernetes cho phép tự động mở rộng số lượng instance theo tải hệ thống, hỗ trợ rolling update, tự động khởi động lại các dịch vụ gặp lỗi và quản lý tài nguyên hiệu quả hơn. Đồng thời, nền tảng cũng có thể tận dụng các khả năng như Service Discovery, ConfigMap, Secret và Horizontal Pod Autoscaler để nâng cao tính linh hoạt và khả năng vận hành trong môi trường production \cite{docker_compose}.


% ------------------------------------------------------------------------------
% Reference list
% ------------------------------------------------------------------------------
\addcontentsline{toc}{chapter}{Tài liệu tham khảo}
\renewcommand{\bibname}{Tài liệu tham khảo}
\printbibliography

% ------------------------------------------------------------------------------
% Appendices
% ------------------------------------------------------------------------------
% % --------------------------------------------------
% Appendices
% --------------------------------------------------
\appendix

\chapter{Cách sử dụng}


\end{document}